\documentclass[12pt,twoside,a4paper]{scrartcl}
\usepackage[utf8]{inputenc}
\usepackage[T1]{fontenc}
\usepackage[ngerman]{babel}
\usepackage[a4paper,scale=0.8]{geometry} 
\usepackage{fancyhdr}
\usepackage{setspace} 


%\pagestyle{fancy} %eigener Seitenstil
%\fancyhf{} %alle Kopf- und Fußzeilenfelder bereinigen
%\fancyhead[L]{} %Kopfzeile links
%\fancyhead[C]{} %zentrierte Kopfzeile
%\fancyhead[R]{\leftmark} %Kopfzeile rechts
%\renewcommand{\headrulewidth}{0.4pt} %obere Trennlinie
%\fancyfoot[C]{\thepage} %Seitennummer
%\renewcommand{\footrulewidth}{0.4pt} %untere Trennlinie
\pagestyle{fancy}
\fancyhead{}
\fancyhead[RE]{\textit{\nouppercase{\leftmark}}}
\fancyhead[LO]{\textit{Die Minimalität der Feichtingeralgebra}}
\fancyhead[RO,LE]{\thepage}
\fancyfoot{}
\addtolength{\headheight}{0.15cm}

\usepackage{paralist}
\usepackage{enumitem} 

%\usepackage {picins}
\usepackage{color}
\usepackage{float}

% -- Mathe
%\usepackage{gensymb}
%fancyhdr, lastpage, booktabs, xy
\usepackage{graphicx}
\usepackage{amsmath}
\usepackage{amsfonts}
\usepackage{amssymb}
\usepackage{units}
\usepackage{amsthm}

\usepackage{hyperref}


% -- Malen
\usepackage{tikz}
\usepackage{tikz-cd}
\usetikzlibrary{patterns, decorations.pathreplacing, decorations.pathmorphing, arrows}

\usepackage{units}

\renewcommand{\theenumi}{\textbf{(\roman{enumi})}}

\numberwithin{equation}{section}

\author{Philipp Beck}



%\makeatletter
%\renewcommand*{\thesection}{\Roman{\thesection}}
%\renewcommand*{\p@section}{\thechapter.}

%Befehle für bestimmte mathematische Symbole
\newcommand{\id}{\operatorname{Id}}

\newcommand{\nt}{\unlhd}
\newcommand{\td}[1]{  \text{ d} #1}
\newcommand{\supp}{\mathrm{supp}}
%\newcommand{\esssupp}{\operatorname{\mathrm{ess \ sup}}}
\DeclareMathOperator*{\esssupp}{ess\,sup}
\newcommand{\Bild}{\mathrm{Bild}}
\newcommand{\Kern}{\mathrm{Kern}}
%Mengensymbole
\newcommand{\N}{\mathbb{N}}
\newcommand{\Z}{\mathbb{Z}}

\newcommand{\Q}{\mathbb{Q}}
\newcommand{\R}{\mathbb{R}}
\newcommand{\C}{\mathbb{C}}
\renewcommand{\L}{\mathrm{L}}
\renewcommand{\i}{\textrm{i}}
\newcommand{\M}{\mathrm{M}}
\renewcommand{\S}{\mathcal{S}}
\newcommand{\F}{\mathcal{F}}
\newcommand{\T}{\mathcal{T}}
\newcommand{\W}{\mathrm{W}}
\renewcommand{\H}{\mathrm{H}}
\renewcommand\qedsymbol{$\blacksquare$}

%Big and italic
\newcommand{\bi}[1]{\textbf{\textit{#1}}}

\swapnumbers

\newtheoremstyle{satz}% name
{15pt}% Space above
{5pt}% Space below
{\itshape}% Body font
{}% Indent amount: Indent amount: empty = no indent, \parindent = normal paragraph indent
{\bfseries}% Theorem head font
{:}% Punctuation after theorem head
{\newline }% Space after theorem head: Space after theorem head: { } = normal interword space; \newline = linebreak
{}% Theorem head spec (can be left empty, meaning `normal')

\newtheoremstyle{def}% name
{15pt}% Space above
{5pt}% Space below
{}% Body font
{}% Indent amount: Indent amount: empty = no indent, \parindent = normal paragraph indent
{\bfseries}% Theorem head font
{:}% Punctuation after theorem head
{\newline }% Space after theorem head: Space after theorem head: { } = normal interword space; \newline = linebreak
{}% Theorem head spec (can be left empty, meaning `normal')

\newtheoremstyle{exercise}% name
{15pt}% Space above
{5pt}% Space below
{}% Body font
{}% Indent amount: Indent amount: empty = no indent, \parindent = normal paragraph indent
{\bfseries}% Theorem head font
{:}% Punctuation after theorem head
{0.5em}% Space after theorem head: Space after theorem head: { } = normal interword space; \newline = linebreak
{}% Theorem head spec (can be left empty, meaning `normal')





\newcommand{\thistheoremname}{}
\theoremstyle{satz}
%Satz
\newtheorem{sz}{Satz}[section]

%Lemma
\newtheorem{lem}[sz]{Lemma}
\newtheorem{kor}[sz]{Korollar}
\newtheorem{gsz}[sz]{\thistheoremname}
\newtheorem*{gthm_no_num}{\thistheoremname}


\theoremstyle{def}
\newtheorem{df}[sz]{Definition}
\newtheorem{gd}[sz]{\thistheoremname}

\newtheorem*{gdf_no_num}{\thistheoremname}


\theoremstyle{exercise}
\newtheorem{exe}{Aufgabe}[section]
\newtheorem{loes}{Lösung}[section]

%Für Beispiele , Bemerkungen
%Also alles was wie eine Definition aussehen soll, aber keine ist.
%\begin{genericdf}{(Hier den Text einfügen)}
%	Inhalt...
%\end{genericdf}
\newenvironment{genericdf}[1]{\renewcommand{\thistheoremname}{#1}\begin{gd}}{\end{gd}}
%Dasselbe ohne Nummerierung
\newenvironment{generic_no_num}[1]{\renewcommand{\thistheoremname}{#1}\begin{gdf_no_num}}{\end{gdf_no_num}}

%Hier für Sätze
\newenvironment{genericthm}[1]{\renewcommand{\thistheoremname}{#1}\begin{gsz}}{\end{gsz}}
%Analog zu Definitionen
\newenvironment{genericthm_no_num}[1]{\renewcommand{\thistheoremname}{#1}\begin{gthm_no_num}}{\end{gthm_no_num}}
\parindent0pt
%\parskip1ex
\newcommand{\blankpage}{
	\newpage
	\thispagestyle{empty}
	\mbox{}
	\newpage
}

%----Dokument----
%\allowdisplaybreaks 
\onehalfspacing
\begin{document}
%Deckblatt
\thispagestyle{empty}
\begin{center}
	
	
	\hspace{-1.0cm}
	\includegraphics[width=0.6\textwidth]{uni_logo.png}
	\vspace{-0.6 cm}
	
	\hspace{2,2cm}
	\\[1cm]
	{\Large Institut für Analysis, Dynamik und Modellierung\\(IADM)}
	
	
	\vspace*{3cm}
	
	% Art der Arbeit => (Bachelorarbeit ,Diplomarbeit, Masterarbeit, Seminararbeit)
	\Large{\textbf{Bachelorarbeit}\\
		
		\vspace{1cm}
		
		% Titel der Arbeit 
		\textbf{\Large{Die Minimalität der Feichtingeralgebra}}
		
		\vspace*{1mm}
		
		
		\vspace{1.5cm}
		
		% Name des/der Autors/Autoren
		%{\LARGE Alisa Baransegeta}\\[35mm]
		
		% Gutachter, Kontaktdaten und Abgabetermin
		\parbox{120mm}{
			\begin{large}
				\begin{tabbing}
					Betreuer: \hspace{0.7cm} \= PD Dr. Peter H. Lesky\\[4mm]
					Verfasser:\> Philipp Beck\\ % alphabetische Reihenfolge (Nachname)
					Matrikel-Nr.:\> 2800897\\
					Email:\>uni@phil-beck.de\\
					%\textbf{Abgabetermin:} \hspace{0,3cm} \> \textbf{21. April 2019}\\
				\end{tabbing}
			\end{large}
		}}
	\end{center}
%\pagestyle{fancy}
\blankpage	
\cleardoublepage
\setcounter{page}{1}

\tableofcontents


\blankpage

%\section{Einleitung}
\chapter{Einleitung}
Bereits in den Grundvorlesungen der Analysis begegnet man Aussagen zur Umordnung von Reihen. Eines der berühmtesten Beispiele hierzu ist der Riemannsche Umordnungssatz, welcher bereits im Jahr 1867 bewiesen wurde. 
Dieser besagt, dass eine beliebige reelle Zahl durch Umordnung einer bedingt konvergenten Reihe als Grenzwert dargestellt werden kann.
Als Konvergenzbereich $ \mathcal{K} $ bezeichnen wir die aus der Umordnung entstehenden Grenzwerte.
Der Konvergenzbereich einer bedingt konvergenten Reihe auf der reellen Achse ist demnach $ \R $. 
%Dies widerspricht zunächst dem intuitiven Verständnis der Addition.
Mit der Fragestellung der Umordnung von Reihen beschäftigt man sich in den frühen Semestern nur am Rande.
Der Fokus liegt auf der Betrachtung einer Reihe als Partialsummenfolge und auf Techniken für den Nachweis der Konvergenz. Damit bleibt der Riemannsche Umordnungssatz oft der einzige Vorstoss in diese Richtung. 
Eine weitere wichtige Tatsache ist, dass absolut konvergente Reihen beliebig umsortiert werden können und der Grenzwert erhalten bleibt.
Erfüllt eine Reihe diese Umordnungseigenschaft nennt man sie unbedingt konvergent.
Für reelle Reihen sind die Begriffe der absoluten und unbedingten Konvergenz äquivalent. 
Diese Arbeit beschäftigt sich mit der bedingten und unbedingten Konvergenz von Reihen.
In dem Bereich der bedingten Konvergenz untersuchen wir die Struktur des Konvergenzbereichs $ \mathcal{K} $.
Für die unbedingte Konvergenz betrachten wir mehrere kleine Themenkomplexe.
\\
\\
Wir untersuchen die Erweiterungen des Riemannschen Umordungssatzes zunächst auf endlich vielen Dimensionen und darauffolgend auf beliebigen Banachräume.
Die Struktur bauen wir auf der von M.I. und V.M. Kadets in \cite{Kadets1991} und \cite{Kadets1997} auf.
Beide Werke beschäftigen sich mit Fragestellungen zur bedingten und unbedingten Konvergenz in Banachräumen und wurden gemeinsam von Vater und Sohn geschrieben. 
%Im Grunde genommen ist \cite{Kadets1997} als eine Überarbeitung von \cite{Kadets1991} zu sehen.
Unter einer Erweiterung verstehen wir, dass der Konvergenzbereich in der Form $ \mathcal{K} = s + U $ vorliegt. Hierbei ist $ U $ ein Untervektorraum des Ausgangsraumes und $ s $ eine affine Verschiebung. 
Im Gegensatz zu der Aussage für die reelle Achse können sich die durch Umordnung erreichbaren Grenzwerte also in einem (echten) affinen Unterraum befinden.
Mithilfe des Riemannschen Umordnungssatzes lässt sich jeder beliebige affine Unterraum in einem endlichdimensionalen Raum konstruieren. Konkret werden wir dies in \ref{sc:construction_steiniz_prop_finite}
durchführen. Die Erweitung auf die komplexen Zahlen bzw. den zweidimensionalen Raum von Paul L\'evy wurde 1905 in \cite{Levy1905} geführt.
Den allgemeinen endlichdimensionalen Fall hat Ernst Steinitz 1913 in \cite{Steiniz1913} gezeigt.
Diesen Beweis hat er auf Anregung von Edmund Landau geführt, nachdem dieser ihn auf Fehler bezüglich der Reihentheorie in damals geläufigen Lehrbüchern hinwies. 
Dies veranlasste Herr Steinitz die Arbeit von L\'evy zu überprüfen um diese Fehler zu beheben.
Hierbei stellte er fest, dass der komplexe Fall von L\'evy bewiesen wurde. Für höhere Dimensionen war diese Arbeit jedoch unvollständig.
Steinitz hat diese Lücken gefüllt und den Riemannschen Satz  auf jede endliche Dimension erweitert. Außerdem hat er dies auch auf eine andere Weise gezeigt, welche wir in dieser Arbeit betrachten werden.
Aufgrund dieser Geschichte wird die Erweiterung auf beliebige Dimensionen auch Satz von L\'evy-Steinitz genannt.
In diesem Zusammenhang ein Zitat aus der Arbeit von Steinitz\cite{Steiniz1913}:
\begin{quote}
	\textit{Wenn ich in Verbindung damit nun doch das Reihenproblem behandele, so geht die Anregung hierzu auf Herrn Landau zurück, der in neueren Lehrbüchern der Analysis bezüglich der komplexen Reihen geradezu falsche Angaben fand, sich auch bei der Durchsicht der L\'evyschen Arbeit von ihrer Richtigkeit nicht überzeugen konnte.
	Dies hat mich veranlaßt, diese Arbeit nochmal aufs genauste durchzustudieren. Sie ist sehr knapp, ja sprunghaft gehalten und im Ausdruck oft unklar. Wer die Lösung nicht bereits kennt, wird stellenweise kaum erraten können, was gemeint ist. Kennt man die Lösung, so überzeugt man sich leicht, daß die Gesichtspunkte, in deren Auffindung die hauptsächlichen Schwierigkeiten der Aufgabe liegen, soweit sie sich auf die gewöhnlichen komplexen Zahlen bezieht, wirklich gefunden sind.
	Daß die Ausführungen im einzelnen richtig sind , wenn sie auch zur vollständigen Erledigung der Aufgabe nicht ausreichen, das zu konstatieren, ist aber auch dann noch recht mühsam. So sehr es zu mißbilligen ist, wenn Publikationen in so mangelhafter Form erfolgen, daß zum Verständnis noch ausführliche Kommentare notwendig sind, daß Herr L\'evy den angegebenen Satz für Reihen mit gewöhnlichen komplexen Zahlen in der zitierten Arbeit zum großen Teil bewiesen hat. Anders steht es mit dem Satze für die allgemeinen komplexen Zahlen, da hier nur Resultate angegeben werden. Damit soll nicht gesagt sein, daß Herr L\'evy nicht auch für den allgemeinen Fall einen Beweis gefunden hat, noch weniger, daß er nicht habe finden können. Aber man kann nicht sagen, daß dieser Beweis in seiner Arbeit enthalten ist, da hierzu ergänzende Untersuchungen im $ n $-dimensionalen Raume notwendig werden, die nur im Falle $ n =2  $ noch trivial sind. }  \ \  \ \  \  \ \ \ \ \ \ \ \  \ \  \  \ \ \ \ \ \ \ \ \ \ \  \ \ \ \ \ \ \ \  \ \ 
	Ernst Steinitz(1913)
\end{quote}

Beim Übergang zu unendlichdimensionalen Räumen reicht die Voraussetzung der (bedingten) Reihenkonvergenz nicht mehr aus.
Im Jahr 1954 fand M.I. Kadets für $ \L^p $-Räume in \cite{Kadets1954} die hinreichende Bedingung
\begin{align*}
	\sum \limits_{ i = 1}^\infty \| x_i \|^{\max\{2,p\}}_{\L^p} < \infty
\end{align*}
für einen Konvergenzbereich der Form $ s + U. $ Diese wurde 1973 
durch
\begin{align*}
	\left(\sum \limits_{ i = 1}^\infty | x_i(t) |^2\right)^\frac{1}{2} \in \L^p
\end{align*}
Nikishin\cite{Nikishin1973} abgeschwächt.
Weitere Verbesserungen wurden von Trojanski\cite{Trojanski1967}, Fonf\cite{Fonf1972} und Ostrovskii\cite{Ostrovskii1986} gefunden.
Jedoch haben all diese Erweiterungen gemein, dass sich der endlichdimensionale Steinitzsatz nicht hieraus ergibt.
Diese Tatsache wurde unabhängig voneinander von Pecherskii \cite{Pecherski1989} und Chobanyan\cite{Chobanyan1987},\cite{Chobanyan1989} beseitigt.
Für die Aussage von Pecherskii werden wir dies konkret in dem Abschnitt \ref{sc:hanani_pecherskii} nachweisen.
Insbesondere wird die von Chobanyan formulierte Aussage mithilfe einer von Vakhania\cite{Vakhania:1987} bewiesenen Äquivalenz aus der Formulierung von Pecherskii folgen.\\
\\
Der zweite Teil der Arbeit richtet sich an die unbedingte Konvergenz.
Wir werden uns hauptsächlich nach der Struktur von Diestel\cite{Diestel1995} richten und insbesondere in der gesamten Arbeit dessen Begriff der Summierbarkeit übernehmen.
In dem Abschnitt \ref{sc:char_uncon_sum} widmen wir uns der Charakterisierung der unbedingten Konvergenz, indem wir hierzu äquivalente Aussagen nachweisen. Hierbei spielt die von Schur\cite{Schur1920} bewiesene Äquivalenz von schwacher Konvergenz und Normkonvergenz in $ \ell^1 $ eine große Rolle.
Auf bestimmten unendlichdimensionalen Räumen lassen sich leicht Gegenbeispiele für die Äquivalenz unbedingter und absoluter Konvergenz konstruieren.
Von Dvoretzky und Rogers wurde in  \cite{DvoretzkyRogers1950} gezeigt, dass keine unendlichdimensionalen Räume mit dieser Eigenschaft existieren. 
Effektiv haben diese gezeigt, dass jeder unendlichdimensionale Raum eine unbedingt konvergente Reihe enthält, wofür die Reihenglieder $ \| x_n \| = \nicefrac{1}{n} $ erfüllen. Dies zeigt, warum man die Begriffe der absoluten und unbedingten Konvergenz trennen muss. Den Beweis hierzu werden wir in dem Abschnitt \ref{sc:dv_rg} führen.
Auf $ \L^p $-Räumen hat Orlicz\cite{Orlicz1930} im Jahr 1930 gezeigt, dass aus unbedingter Konvergenz
\begin{align*}
	\sum \limits_{i = 1}^\infty \|x_i\|^{\max\{2,p\}}_{\L^p} < \infty
\end{align*}
für $ 1 \leq p < \infty  $ folgt. Diesen Beweis werden wir in dem Abschnitt \ref{sc:lemma_of_orlicz} führen. In dem letzten Abschnitt \ref{sc:abs_summing} beschäftigen wir uns mit Operatoren, welche unbedingte Konvergenz in absolute Konvergenz überführt.
Der Fokus wird auf einer von Grothendieck\cite{Grothendieck1956} formulierten Aussage sein. Diese besagt, dass jeder Operator von $ \ell^1  $ nach $ \ell^2 $ unbedingt konvergente Reihen auf absolut konvergente Reihen abbildet.











\blankpage
\section{Grundlagen}
\vspace{-0.5cm}
In diesem Abschnitt werden wir grundlegende Bezeichnungen, Definitionen und Sätze einführen.
Unter 
\begin{align}
\int \limits_{\R^d} f(x) \td{x}
=
\int \limits_{-\infty}^\infty \cdots \int \limits_{-\infty}^\infty f(x_1,...,x_d) \td{(x_1,...,x_d)} 
\end{align}
mit $ \td{(x_1,...,x_d)} = \td{x_1}...\td{x_d} $ verstehen wir das Lebesgueintegral auf $ \R^d $.
Sei $ f  $ eine messbare Funktion auf $ \R^d $.
Dann gilt $ f \in \L^p(\R^d) $, falls die $ \L^p $-Norm
\begin{align}
\| f \|_p = \left(  \ \int \limits_{\R^d} | f(x) |^p \td{x} \right)^{\frac{1}{p}}
\end{align}
für $ 1 \leq p < \infty $ endlich ist.
$ \L^\infty(\R^d) $ besteht aus den essentiell beschränkten messbaren Funktionen, d.h. es gilt
\begin{align*}
\|f \|_\infty 
= 
\esssupp \limits_{x \in \R^d} | f(x) | < \infty.
\end{align*}
Wir werden $ \sup $ auch für das essentielle Supremum verwenden, falls dies sinnvoll ist.
Die $ \L^p $-Räume sind für $ 1 \leq p \leq \infty $ allesamt Banachräume. 
$ \L^2(\R^d) $ ist mit dem Skalarprodukt
\begin{align}
\langle f, g \rangle
=
\int \limits_{\R^d} f(x) \overline{g(x)} \td{x}
\end{align}
ein Hilbertraum. Wenn der Kontext klar ist, werden wir $ \L^p $ statt $ \L^p(\R^d) $ schreiben.
In wenigen Fällen ist $ \| \cdot \|_2 $ die euklidische Norm auf $ \R^d $.
Mit $ x \cdot \omega = \sum_{k=1}^d x_i \omega_i$ beschreiben wir das Skalarprodukt auf $ \R^d $.
\begin{df}
	Sei $ f \in \L^1(\R^d) $. Dann definieren wir durch
	\begin{align}\label{eq:fourier_transformation}
	\hat{f}(\omega) =
	\int \limits_{\R^d} f(x) e^{-2 \pi \i x \cdot \omega} \td{x}
	\end{align}
	die \textit{Fouriertransformation} von $ f $.	
\end{df}
Für \eqref{eq:fourier_transformation} erhalten wir:
\begin{align}\label{eq:estimate_fouriertransformation}
\| \hat{f} \|_{\infty}
\leq
\sup \limits_{\omega \in \R^d} \int \limits_{\R^d} |f(x) e^{-2 \pi \i x \cdot \omega}| \td{x} = \| f \|_1
\end{align}
Falls wir die Fouriertransformation als linearen Operator auf einem Funktionenraum betrachten, schreiben wir statt $ \hat{f} $ auch $ \F f $.
\begin{genericthm}{Lemma von Riemann-Lebesgue}
	Sei $ f \in \L^1(\R^d) $.
	Dann ist $ \hat{f} $ gleichmäßig stetig und es gilt
	$ \lim_{|\omega| \to \infty} | \hat{f}(\omega) | = 0 $.
\end{genericthm}
Mit diesem Lemma können wir die Fouriertransformation durch
\begin{align*}
\F : \L^1(\R^d) \to C_0(\R^d)
\end{align*}
beschreiben. Hierbei ist $ C_0(\R^d) :=  \{ f \in C(\R^d) |   \lim_{|\omega| \to \infty} | \hat{f}(\omega) | = 0\} $. 

\begin{genericthm}{Dichtheitsprinzip}\label{th:densitiy_priniciple}
	Gegeben seien die Banachräume  $ B_1 $ und $ B_2 $, ein dichter Unterraum $ X $ von $ B_1 $ und ein linearer Operator $ A : X \to B_2 $ mit
	\begin{align}\label{eq:density_principle}
	\| A f \|_{B_2} \leq C \| f \|_{B_1}
	\end{align}
	für  alle $  f \in X $.
	Dann gilt \eqref{eq:density_principle} für alle $ f \in B_1 $ und $ A $ wird zu einem beschränkten linearen Operator von $ B_1  $ nach $ B_2 $ fortgesetzt.
\end{genericthm}
\begin{sz}
	Sei $ f \in \L^1(\R^d) \cap \L^2(\R^d) $. Dann gilt die \textit{Plancherelgleichung}
	\begin{align}\label{eq:plancherel}
	\|f \|_2 = \| \hat{f} \|_2.
	\end{align}
	.
\end{sz}
Damit lässt sich $ \F $ zu einem unitären Operator auf $ \L^2(\R^d) $ erweitern, womit auch die \textit{Parsevalgleichung}
\begin{align}\label{eq:parseval}
\langle f, g \rangle = \langle \hat{f} , \hat{g} \rangle
\end{align}
für alle $ f,g \in \L^2(\R^d) $ erfüllt ist.
Für $ f \in \L^2(\R^d)  $ können wir $ \hat{f} $ jedoch nicht mehr wie in \eqref{eq:fourier_transformation} punktweise definieren. 
Hier ist uns das Dichtheitsprinzip \ref{th:densitiy_priniciple} behilflich.
Sei $ X = \L^1 \cap \L^2 $ und $ f \in \L^2  $ beliebig. Dann ist $ X $ dicht in $ \L^2 $ und es existiert ein $f_n  \in X$ mit $ \| f - f_n \|_2 \to 0 $.
Wegen $ f_n \in \L^1 $ ist $ \hat{f}_n $ wohldefiniert und es gilt mit \eqref{eq:plancherel} $ \| f_n - f_m \|_2 = \| \hat{f}_n - \hat{f}_m \|_2 $.Damit ist $ \hat{f}_n $ eine Cauchyfolge in $ \L^2 $ und wegen der Vollständigkeit existiert ein eindeutiger Grenzwert in $ \L^2 $.
Dementsprechend definieren wir $ \hat{f} := \lim_{n \to \infty} \hat{f}_n $.

\newpage
\begin{df}
	Seien $ f, g \in \L^1(\R^d) $. Dann definieren wir durch
	\begin{equation}\label{eq:convolution}
	(f \ast g)(x) = \int \limits_{\mathbb{R}^d}
	f(y) g(x-y) \td{y}
	\end{equation}
	die \textit{Faltung} von $ f $ und $ g $.
\end{df}
Wegen
\begin{align}\label{eq:convolution_prop_1}
\| f \ast g \|_1 &\leq \| f\|_1 \|g \|_1
\end{align}
ist  $ \L^1 $ eine Banachalgebra bezüglich der Faltung. Außerdem ist
\begin{align}\label{eq:convolution_prop_2}
\widehat{(f \ast  g)} &= \hat{f} \cdot \hat{g}
\end{align}
erfüllt.

\begin{genericdf}{Translation und Modulation}
	Seien $ x, \omega \in \R^d $ und $ f : \R^d \to \C$ eine beliebige Funktion.
	Dann ist
	\begin{align}\label{eq:translation}
	T_x f(t)=  f(t-x) 
	\end{align}
	die \textit{Translation} um $ x $ und
	\begin{align}\label{eq:modulation}
	M_\omega f(t) = e^{2 \pi \mathrm{i} t \cdot \omega} f(t)
	\end{align}
	die \textit{Modulation} um $ \omega $. 
	Wir nennen Operatoren der Form $M_\omega T_x$ oder $T_x M_\omega$ \textit{Zeit-Frequenz-Verschiebungen}. 
 \end{genericdf}

Die elementaren Eigenschaften dieser Operatoren werden uns im Laufe dieser Arbeit häufig begegnen, weswegen wir diese direkt beweisen werden.

\begin{lem}\label{th:properties_TF}
	\begin{enumerate}[label =\textbf{(\roman*)}]
		\item Es gilt
		\begin{align}\label{eq:trans_mod_commutation_relation_1}
		T_x M_\omega = e^{-2 \pi \mathrm{i} x \cdot \omega} M_\omega T_x
		\end{align}
		Damit kommutieren $T_x$ und $M_\omega$ genau dann, wenn $x \cdot \omega \in \mathbb{Z}$.
		\item
		Zeit-Frequenz-Verschiebungen sind Isometrien auf $L^p$ für $1 \leq p \leq \infty$.
		\item
		Es gelten
		\begin{equation}\label{eq:trans_mod_fourier_1}
		\widehat{(T_x f)}= M_{-x} \hat{f} \ \text{und} \ \widehat{(M_\omega f)} = T_\omega \hat{f},
		\end{equation}
		woraus
		\begin{equation}\label{eq:trans_mod_fourier_2}
		\widehat{(T_x M_\omega f)} = M_{-x}T_\omega \hat{f} = e^{-2\pi \mathrm{i} x \cdot \omega} T_\omega M_{-x} \hat{f}
		\end{equation}
		folgt.
	\end{enumerate}
\end{lem}

\begin{proof}
	\begin{enumerate}[label =\textbf{(\roman*)}] 
		\item
		Durch
		\begin{align*}
		T_x M_\omega f(t) 
		&= M_\omega f(t-x)
		= e^{2\pi \mathrm{i} (t-x) \cdot \omega} f(t-x)\\
		&= e^{-2\pi \mathrm{i} x \cdot \omega} e^{2\pi \mathrm{i} t \cdot \omega} T_x f(t)
		= e^{-2\pi \mathrm{i} x \cdot \omega} M_\omega T_x f(t)
		\end{align*}
		erhalten wir die Aussage.
		
		\item
		Sei $1\leq p < \infty$.
		Es gilt
		\begin{align*}
		\| M_\omega T_x f\|_p^p
		=
		\int \limits_{\mathbb{R}^d} | M_\omega T_x f(t) |^p \td{t}
		=
		\int \limits_{\mathbb{R}^d} | T_x f(t) |^p \td{t}
		=
		\int \limits_{\mathbb{R}^d} | f(t) |^p \td{t}
		=
		\|  f\|_p^p
		\end{align*}
		für $1\leq p < \infty$.
		Der Beweis für $p = \infty$ funktioniert analog.
		
		\item
		Durch
		\begin{align*}
		\widehat{(T_x f)} (\omega) 
		=
		\int \limits_{\R^d} f(t-x) e^{-2 \pi \mathrm{i } t \cdot \omega} \td{t}
		=
		e^{-2 \i x\cdot \omega}
		\int \limits_{\R^d} f(u) e^{-2 \pi \mathrm{i} u \cdot \omega} \td{u}
		= M_{-x} \hat{f}(\omega)
		\end{align*}
		erhalten wir die erste Aussage.
		Die Zweite folgt analog.
		Durch Anwenden der beiden Aussagen
		erhalten wir die Folgerung.
	\end{enumerate}
\end{proof}


\begin{df}
Wir bezeichnen mit
\begin{align}\label{eq:involution}
f^\ast(x) = \overline{f(-x)}
\end{align}
die \textit{Involution}.	
\end{df}
Für die Involution ist 
\begin{align}\label{eq:involution_fouriertrans}
\widehat{f^\ast} = \hat{\overline{f}}
\end{align}
erfüllt. Die Faltung lässt sich nun auch durch
\begin{align}\label{eq:convolution_with_translation}
(f \ast g)(x) = \langle f, T_x g^\ast \rangle
\end{align}
beschreiben, falls beide Seiten definiert sind.
Der Differentiations -und Multiplikationsoperator ist durch
\begin{align*}
D^\alpha f
&=
\prod \limits_{j=1}^d \left(\frac{\partial^{\alpha_j} }{\partial x_j^{\alpha_j}}\right) f\\
(X^\beta f)(t)
&= t^\beta f(t) 
=
\left(\prod \limits_{j = 1}^d t_j^{\beta_j}\right) f(t)
\end{align*}
für $ \alpha,\beta \in \N_0^d $ definiert.
Für $ \alpha ,\beta\in \N_0^d $ verwenden wir außerdem:
\begin{align*}
|\alpha | &:= \sum \limits_{j =1}^d \alpha_j\\
x^\alpha &:= \sum \limits_{j=1}^d x_j^{\alpha_j}\\
\alpha \leq \beta &:\Leftrightarrow
\forall j \in \{1,...,d\} : \alpha_j \leq \beta_j.
\end{align*}
Für die Fouriertransformation gelten
\begin{align}\label{eq:fourier_derivate_1}
\widehat{(D^\alpha f)}(\omega)
=(2 \pi  \i \omega)^\alpha \hat{f}(\omega)
\end{align}
und 
\begin{align}\label{eq:fourier_derivate_2}
\widehat{((- 2 \pi \i x)^\alpha f)}(\omega)
=
D^\alpha \hat{f}(\omega).
\end{align}

Der \textit{Schwartzraum} $ \S(\R^d) $ ist über die Endlichkeit der Seminormen $ \| D^\alpha X^\beta f \|_\infty $ für $ f \in C^\infty(\R^d) $ definiert, d.h. es gilt
\begin{align*}
\| D^\alpha X^\beta f \|_\infty
=
\sup \limits_{x \in \R^d} | D^\alpha X^\beta f(x) |
 < \infty
\end{align*}
für alle $ \alpha,\beta \in \N_0^d $.
Eine Folge $ f_n \in \S(\R^d) $ konvergiert gegen $ f \in \S(\R^d) $, falls
\begin{align}
\| D^\alpha X^\beta (f_n - f ) \|_\infty \to 0
\end{align}
für alle $ \alpha,\beta \in \N_0^d $ gilt.

\newpage 
\begin{sz}
	Die Fouriertransformation $ \F : \S(\R^d) \to \S(\R^d)  $ ist ein Isomorphismus, d.h. stetig und bijektiv mit stetiger Inversen.
	Die inverse Abbildung ist durch
	\begin{align}\label{eq:fouriertransformation_inverse}
	(\F^{-1} \hat{f} )(x) = f(x) 
	=
	\int \limits_{\R^d} \hat{f}(\omega)e^{2\pi \i x \cdot \omega} \td{\omega}
	\end{align}
	gegeben.
\end{sz}
Der Dualraum $ \S^\prime(\R^d) $ ist der Raum der temperierten Distributionen und für $ f \in \S^\prime(\R^d) $ gilt folgende Stetigkeitsabschätzung:\\
Es existieren $ C >0  $ und $ N,M > 0  $, sodass
\begin{align}\label{eq:cont_distribution}
\langle f, \varphi \rangle 
\leq
C 
\sum \limits_{|\alpha| \leq M}
\sum \limits_{|\beta| \leq N}
\| D^\alpha X^\beta \varphi \|_\infty
\end{align}
für alle $ \varphi \in \S(\R^d) $ gilt. Sei $ f \in \S^\prime(\R^d) $.\\ Die Fouriertransformation von $ f $  ist durch
\begin{align}
\hat{f} (\varphi ) := \langle \hat{f}, \varphi \rangle := \langle f , \hat{\varphi} \rangle
\end{align}
für alle $ \varphi \in \S(\R^d) $ definiert. 

\begin{genericthm}{Open-Mapping-Theorem}\label{th:open_mapping}
	Gegeben seien die Banach -oder Frecheträume  $ B_1 $ und $ B_2 $ und der surjektive beschränkte Operator $ A : B_1  \to B_2 $.
	Dann ist $ A $ eine offene Abbildung.
	Insbesondere können wir diesen auf Operatoren von $ \S(\R^d)  $ nach $ \S(\R^d) $ anwenden.
\end{genericthm}

\begin{genericthm}{Closed-Range-Theorem}\label{th:closed_range}
	Gegeben seien die Banachräume $ B_1 $ und $ B_2 $
	und ein beschränkter linearer Operator $ A : B_1 \to B_2 $.
	Der Operator $ A $ hat ein abgeschlossenes Bild genau dann, wenn der 
	adjungierte Operator $ A^\prime : B_2^\prime \to B_1^\prime $
	ein abgeschlossenes Bild besitzt.
	Insbesondere ist $ A $ genau dann surjektiv, wenn $ A^\prime $ injektiv ist und ein abgeschlossenes Bild hat.
\end{genericthm}

\newpage

\section{Kurzzeit-Fouriertransformation(STFT)}

%Wir werden unsere Theorie vollständig auf der STFT aufbauen.
In diesem Abschnitt definieren wir die Kurzzeit-Fouriertransformation und beweisen 
ihre elementaren Eigenschaften.
Zunächst werden wir einige fundamentale Identitäten zeigen, welche uns helfen, die STFT so allgemein wie möglich zu definieren. 
Im Anschluss beweisen wir die Parseval-und Plancherelgleichung für die STFT, welche wir für die Inversionsformel verwenden werden.
Wir benötigen die STFT, um später die Modulationsräume zu definieren.


\begin{df}
	Sei $g \neq 0$ eine glatte Abschneidefunktion.
	Diese bezeichnen wir als \textit{Fensterfunktion}.\index{Fensterfunktion}
	Dann ist die \textit{Kurzzeit-Fouriertransformation} einer Funktion $f$ abhängig von $g$ durch
	\begin{align*}
	V_g f(x, \omega) 
	:=
	\int \limits_{\mathbb{R}^d}
	f(t) \overline{g (t-x) }  e^{-2\pi \mathrm{i} t\cdot \omega }\td{t}
	\end{align*}
	für $x, \omega \in \mathbb{R}^d$ definiert.
\end{df}\noindent
Die Kurzzeit-Fouriertransformation ist mit 
\begin{align*}
(f,g) \mapsto V_g f
\end{align*}
eine Sesquilinearform. Meist lassen wir die Fensterfunktion $g$ jedoch fixiert.
Wir betrachten nun die asymmetrische Koordinatentransformation
\begin{align*}
\mathcal{T}_a F(x,t) = F(t,t-x)
\end{align*}
und die partielle Fouriertransformation
\begin{align*}
\mathcal{F}_2 F(x, \omega)
= 
\int \limits_{\R^d} F(x,t) e^{- 2 \pi \i t \cdot \omega} \td{t}
\end{align*}
von $F$ auf $\R^{2d}$.
Mit dem Tensorprodukt $(f \otimes g) (x,t) = f(x) g(t)$ erhalten wir somit
\begin{align*}
V_g f = \mathcal{F}_2 \mathcal{T}_a(f \otimes \overline{g})
\end{align*}
für $f, g \in \mathrm{L}^2(\R^d)$.
\newpage
\begin{sz}\label{fundamental_identities_STFT}
	Seien $f,g \in \mathrm{L}^2(\mathbb{R}^{d})$.
	Dann gelten die Identitäten
	\begin{align}
	V_g f(x,\omega)
	&=
	\widehat{(f \cdot T_x \overline{g})} (\omega)\label{eq:identity_STFT_1}\\
	&=
	\langle f, M_\omega T_x g \rangle\label{eq:identity_STFT_2}\\
	&=
	\langle \hat{f} , T_w M_{-x} \hat g \rangle\label{eq:identity_STFT_3}\\
	&=
	e^{-2 \pi \mathrm{i} x \cdot \omega} \widehat{(\hat{f} \cdot T_\omega \overline{ \hat{g} })} (-x)\label{eq:identity_STFT_4}\\
	&=
	e^{-2\pi \mathrm{i} x \cdot \omega} V_{\hat{g}} \hat{f}(\omega,-x)\label{eq:identity_STFT_5}\\
	&=
	e^{-2\pi \mathrm{i} x \cdot \omega}  (f \ast M_\omega g^\ast)(x) \label{eq:identity_STFT_7}\\
	&=
	(\hat{f} \ast M_{-x} \hat{g}^\ast)(\omega)\label{eq:identity_STFT_8}\\
	&=
	e^{-\pi \mathrm{i} x \cdot w } \int \limits_{\mathbb{R}^d} f\left( t + \frac{x}{2} \right) \overline{g\left( t - \frac{x}{2} \right)} e^{-2 \pi \mathrm{i} t \cdot \omega} \td{t}\label{eq:identity_STFT_9}
	\end{align}
	und $V_g f$ ist gleichmäßig stetig auf $\mathbb{R}^{2d}$.
\end{sz}

\begin{proof}
	Zuerst werden wir die Identitäten zeigen.
	Wegen
	\begin{align*}
	V_g f(x,\omega)
	=
	\int \limits_{\R^d} f(t) \overline{g(t-x)} 
	e^{-2\pi \i t  \cdot \omega} \td{t}
	=
	\widehat{(f \cdot T_x \overline{g})} (\omega)
	=
	\langle f, M_\omega T_x g \rangle
	\end{align*}
	sind \eqref{eq:identity_STFT_1} und \eqref{eq:identity_STFT_2} klar.
	Mit der Parsevalgleichung \eqref{eq:parseval} und \eqref{eq:trans_mod_fourier_2} folgen durch
	\begin{align*}
	\langle f, M_\omega T_x g \rangle
	&=
	\left\langle \hat{f}, \widehat{(M_\omega T_x g)} \right\rangle
	=
	\underbrace{ \langle \hat{f},e^{2\pi \i x \cdot \omega} M_{-x} T_\omega \hat{g} \rangle}_{
		= e^{-2\pi \i x \cdot \omega} V_{\hat{g}} \hat{f}(\omega,-x)
		}
	= 
	\langle \hat{f}, T_\omega M_{-x} \hat{g} \rangle
	\end{align*}
	die Identitäten \eqref{eq:identity_STFT_3} und \eqref{eq:identity_STFT_5}.
	Durch Anwenden von \eqref{eq:identity_STFT_1} erhalten wir \eqref{eq:identity_STFT_4}.
	Die Identität \eqref{eq:identity_STFT_8} folgt durch
	\begin{align*}
	\widehat{(f \cdot T_x \overline{g})} (\omega)
	=
	(\hat{f} \ast M_{-x}  \hat{\overline{g}} )(\omega)
	=
	(\hat{f} \ast M_{-x}  \hat{g^\ast} )(\omega)
	\end{align*}
	unter Verwendung von \eqref{eq:convolution_prop_2} und \eqref{eq:involution_fouriertrans}.
	Analog liefert dies mit
	\begin{align*}
	e^{-2 \pi \mathrm{i} x \cdot \omega} \widehat{(\hat{f} \cdot T_\omega \overline{ \hat{g} })} (-x)
	=
	e^{-2\pi \mathrm{i} x \cdot \omega}  (f \ast M_\omega g^\ast)(x)
	\end{align*}
	die Identität \eqref{eq:identity_STFT_7}.
	Die letzte Identität 
	\begin{align*} 
	V_g f(x,\omega)
	=
	e^{-\pi \mathrm{i} x \cdot w } \int \limits_{\mathbb{R}^d} f\left( t + \frac{x}{2} \right) \overline{g\left( t - \frac{x}{2} \right)} e^{-2 \pi \mathrm{i} t \cdot \omega} \td{t}
	\end{align*}
	folgt durch $u = t  - \frac{x}{2}$.\\
	Im zweiten Teil des Beweises werden wir die gleichmäßige Stetigkeit auf $\R^{2d}$ von $V_g f$ zeigen.
	Zunächst gilt 
	\begin{align*}
	\lim \limits_{x \to 0 } \| T_x f - f \|_2 = 0, 
	\end{align*}
	womit durch $(M_\omega f)^{\hat{ }} = T_\omega \hat{f}$
	auch die Aussage
	\begin{align*}
	\lim \limits_{\omega \to 0} \| M_\omega f - f \|_2
	= 
	\lim \limits_{\omega \to 0} \| T_\omega \hat{f} - \hat{f } \|_2
	= 0
	\end{align*}
	folgt. 
	Unser Ziel ist nun 
	\begin{align*}
	\forall_{\varepsilon > 0 } \exists_{\delta > 0}
	\forall_{(x,\omega) , (x^\prime,\omega^\prime)\in \R^{2d}}: \ | (x,\omega) - (x^\prime,\omega^\prime) | < \delta \ \Rightarrow \ | V_g f(x,\omega) -V_g f(x^\prime,\omega^\prime) | < \varepsilon
	\end{align*}
	zu zeigen.
	Zunächst gilt
	\begin{align*}
	|V_g f(x,\omega) - V_g f(x^\prime, \omega^\prime) |
	&= | \langle f, M_\omega T_x g \rangle - \langle f, M_{\omega^\prime} T_{x^\prime} g \rangle |\\
	&\leq
	\int \limits_{\R^d} |f(t)| \ |M_\omega T_x g(t) - M_{\omega^\prime} T_{x^\prime} g(t) | \td{t}\\
	&\leq 
	\| f \|_2 \cdot \| M_\omega T_x g - M_{\omega^\prime} T_{x^\prime} g \|_2
	\end{align*}
	mit Hilfe der Hölder-Ungleichung.
	Es folgt
	\begin{align*}
	\| M_\omega T_x g - M_{\omega^\prime} T_{x^\prime} g \|_2
	&= 
	\| M_\omega T_x g - M_{\omega^\prime} T_x g + M_{\omega^\prime} T_x g - M_{\omega^\prime} T_{x^\prime} g \|_2\\
	&\leq
	\| M_\omega T_x g - M_{\omega^\prime} T_x g \|_2 
	+ \| M_{\omega^\prime} T_x g - M_{\omega^\prime} T_{x^\prime} g\|_2\\
	&=
	\| (M_\omega - M_{\omega^\prime})T_x g\|_2
	+ \|M_{\omega^\prime} ( T_x - T_{x^\prime})g \|_2
	\end{align*}
	Wir werden die zwei Summanden seperat untersuchen.
	Für den Zweiten erhalten wir
	\begin{align*}
	\| M_{\omega^\prime}(T_x - T_{x^\prime}) g \|_2
	= \|(T_x - T_{x^\prime}) g \|_2
	= \|(T_{x^\prime}( T_{x-x^\prime} - \id)g \|_2
	= \| ( T_{x-x^\prime} - \id)g \|_2 < \frac{\varepsilon}{2 \| f \|_2}
	\end{align*}
	für $| x - x^\prime| <\delta_1$.
	Das Argument für den ersten Summanden ist ähnlich.
	Mit \eqref{eq:trans_mod_fourier_1} gilt
	\begin{align*}
	\F((M_\omega - M_{\omega^\prime}) T_x g)
	= 
	(T_\omega - T_{\omega}) M_{-x} g,
	\end{align*}
	wodurch mit der Parsevalgleichung 
	\begin{align*}
	\| (M_\omega - M_{\omega^\prime}) T_x g \|_2
	=
	\| (T_\omega - T_{\omega^\prime}) M_{-x} \hat{g} \|_2
	=
	\| (T_{\omega - \omega^\prime} - \id) g \|_2.
	\end{align*}
    folgt.Sei $ \tilde{g}  \in C_0^\infty(\R^d)$.
	Dann gilt
	\begin{align*}
	\| (T_{\omega - \omega^\prime} - \id) g \|_2 \leq 
	\| T_{\omega - \omega^\prime} g \|_2 + \| g \|_2 
	= 2 \| g \|_2,
	\end{align*}
	womit wir den Satz der majorisierten Konvergenz anwenden können.
	Da $ C_0^\infty $ dicht in $ \L^2 $ ist, finden wir ein $ \tilde{g}  \in C_0^\infty(\R^d)$ mit
	\begin{align*}
	\| \hat{g} - \tilde{g} \|_2 \leq \frac{\varepsilon}{6 \| f\|}.
	\end{align*}
	Damit existiert ein $ \delta_2 > 0 $ , sodass
	\begin{align*}
	\| T_\omega M_{-x} \tilde{g} - T_{\omega^\prime} M_{-x} \tilde{g} \|_2
	< 
	\frac{\varepsilon}{6 \| f \|_2}
	\end{align*} 
	für $ | \omega - \omega^\prime| < \delta_2 $ erfüllt ist.
	Insbesondere erhalten wir
	\begin{align*}
	\| (T_\omega - T_{\omega^\prime}) M_{-x} \hat{g} \|_2
	&\leq 
	\| T_\omega M_{-x} \hat{g} - T_\omega M_{-x} \tilde{g} \|_2
	+
	\| T_\omega M_{-x} \tilde{g} - T_{\omega^\prime} M_{-x} \tilde{g} \|_2
	+
	\| T_{\omega^\prime} M_{-x} \tilde{g} - T_{\omega^\prime} M_{-x} \hat{g} \|_2\\
	&= 2 \| \hat{g} - \tilde{g} \|_2
	+ 
	\| T_\omega M_{-x} \tilde{g} - T_{\omega^\prime} M_{-x} \tilde{g} \|_2
	< 
	\frac{\varepsilon}{2 \| f\|_2}
 	\end{align*}
	für $ | \omega - \omega^\prime| < \delta_2 $.
	Insgesamt folgt durch
	\begin{align*}
	| (x,\omega) - (x^\prime, \omega^\prime) | \leq
	| (x,0) -(x^\prime,0) | + | (0,\omega) - (0, \omega^\prime) | < \delta :=\min\lbrace \delta_1 , \delta_2 \rbrace
	\end{align*}
	die gewünschte Aussage.
\end{proof}

Der letzte Satz liefert uns nun die Möglichkeit, den Definitionsbereich der STFT zu diskutieren.
Sei $ f \in \L^p(\R^d) $ und $ g \in \L^{p^\prime}(\R^d) $ mit $ \nicefrac{1}{p} + \nicefrac{1}{p^\prime} = 1 $.
Dann liefert uns die Hölderungleichung $ f \cdot T_x \overline{g} \in \L^1(\R^d) $ und die STFT ist punktweise über
\begin{align*}
V_g f(x,\omega ) = \widehat{(f \cdot T_x \overline{g})} (\omega)
\end{align*}
definiert. Die Identität
\begin{align*}
V_g f(x,\omega) = \langle f, M_\omega T_x g \rangle
\end{align*}
hilft uns, falls das Integral nicht definiert ist.
Damit ist die STFT auch für $ f \in \mathcal{S}^\prime(\R^d)  $ und $ g \in \mathcal{S}(\R^d) $ im Dualitätssinne definiert.\\
In folgendem Lemma verstehen wir unter einem Isomorphismus einen linearen Homöomorphismus.
\newpage
\begin{lem}\label{th:iso_tensor_distri}
	\begin{enumerate}[label =\textbf{(\roman*)}]
		\item 
		Die Operatoren $ \mathcal{T}_a $ und $ \mathcal{F}_2 $ sind Isomorphismen auf $ \S^\prime(\R^{2d}) $.
		\item
		Seien $ f,g \in \S^\prime(\R^d) $.
		Dann gilt $ f \otimes g \in \S^\prime(\R^{2d}) $.
	\end{enumerate}
\end{lem}

\begin{proof}
	\begin{enumerate}[label =\textbf{(\roman*)}]
		\item 
		Wir betrachten zunächst $ (\tilde{\mathcal{T}}_a) \varphi(x,t)
		:= \varphi(x-t,x)  $.
		Dann gilt
		\begin{align*}
		(\tilde{\mathcal{T}}_a \circ \mathcal{T}_a) \varphi(x,t) = 
		\varphi (x,t)
		=
		(\mathcal{T}_a \circ \tilde{\mathcal{T}}_a) \varphi(x,t)
		\end{align*}
		und somit $ \mathcal{T}_a^{-1}  = \tilde{\mathcal{T}}_a$. Insbesondere ist 
		\begin{align*}
		\mathcal{T}_a : \S(\R^{2d}) \to \S(\R^{2d}), \
		\varphi \mapsto \mathcal{T}_a \varphi
		\end{align*}
		ein Homöomorphismus. 
		Sei $ f \in \S(\R^{2d}) $ und 
		\begin{align*}
		u_f(\varphi) 
		= \int \limits_{\R^{2d}} f(x,t) \overline{\varphi(x,t)} \td{(x,t)}
		\end{align*}
		die durch $ f  $ induzierte temperierte Distribution.
		Hierfür gilt
		\begin{align*}
		u_{\mathcal{T}_a f}(\varphi) 
		&=
		\int \limits_{\R^{2d}}
		f(t,t-x) \overline{\varphi(x,t)} \td{(x,t)}
		=
		\int \limits_{\R^{2d}}
		f(x,t) \overline{\varphi(x-t,x)} \td{(x,t)}\\
		&=
		\int \limits_{\R^{2d}}
		f(x,t) \overline{(\mathcal{T}_a^{-1})\varphi(x,t)} \td{(x,t)}.
		\end{align*}
		Damit können wir $ \mathcal{T}_a $ sinnvoll durch
		\begin{align*}
		\mathcal{T}_a F( \varphi)
		&:=
		F(\mathcal{T}_a^{-1} \varphi)\\
		\mathcal{T}_a^{-1} F( \varphi)
		&:=
		F(\mathcal{T}_a \varphi)
		\end{align*}
		auf $ \S^\prime(\R^{2d}) $ definieren und $ \mathcal{T}_a $ ist auf $ \S^\prime(\R^{2d}) $ bijektiv. 
		Sei $ F_n \to F $ in $ \S^\prime(\R^{2d}) $, d.h. $ F_n(\varphi) \to F(\varphi) $
		für alle $ \varphi \in \S(\R^{2d}) $.
		Dann gilt
		\begin{align*}
		\mathcal{T}_a F_n(\varphi )
		=
		F_n(\mathcal{T}_a^{-1} \varphi) 
		\to 
		F(\mathcal{T}_a^{-1}\varphi) = \mathcal{T}_a F(\varphi)
		\end{align*}
		und $ \mathcal{T}_a $ ist ein Homöomorphismus auf $ \S^\prime(\R^{2d}) $.
		Die partielle Fouriertransformation $ \mathcal{F}_2 : \S^\prime(\R^{2d}) \to \S^\prime(\R^{2d}) $
		ist nach \cite[Seite~129]{mit} ein Isomorphismus.
		
		\item
		Diese Aussage wird in \cite[Theorem~4.13]{mit} bewiesen .
	\end{enumerate}
\end{proof}

Durch das letzte Lemma lässt sich die STFT auch durch
\begin{align}\label{eq:STFT_tensorprod}
V_g f = \mathcal{F}_2 \mathcal{T}_a (f \otimes \overline{g})
\end{align}
für $ f,g \in \S^\prime(\R^{d}) $ definieren.

Ähnlich zur Fouriertransformation gibt es eine Art Parseval -und Plancherelgleichung für die  STFT.


\begin{lem}\label{th:STFT_modulation_TF_argument}
	Ist $ V_g $ definiert, so gilt
	\begin{align}
	V_g  (T_u M_\eta f) ( x, \omega)
	= 
	e^{-2 \pi \i u \cdot \omega}
	V_g f (x - u , \omega - \eta)
	=
	e^{-2 \pi \i u \cdot \omega}
	T_{(u,\eta)} V_g f(x,\omega)
	\end{align}
	für alle $ x, u, \omega , \eta \in \R^d $.
\end{lem}

\begin{proof}
	Durch Anwenden von \eqref{eq:trans_mod_commutation_relation_1} erhalten wir
	\begin{align*}
	M_{- \eta} T_{-u} M_\omega T_x
	=
	e^{2\pi \i u \cdot \omega} M_{\omega - \eta} T_{x - u }.
	\end{align*}
	Damit folgt durch \eqref{eq:identity_STFT_2}:
	\begin{align*}
	V_g (T_u M_\eta f)(x, \omega)
	= 
	\langle T_u M_\eta f , M_\omega T_x g \rangle
	=
	\langle f , M_{-\eta } T_{-u} M_\omega T_x g \rangle
	=
	e^{-2 \pi \i u \cdot \omega}
	V_g f (x - u , \omega - \eta).
	\end{align*}
\end{proof}

\begin{genericthm}{STFT-Parseval}
	Seien $ f_1,f_2, g_1,g_2 \in \L^2(\R^d) $.
	Dann gilt 
	\begin{align}\label{orthogonality_relation}
	\langle V_{g_1} f_1, V_{g_2} f_2 \rangle =
	\langle f_1,f_2 \rangle  \overline{ \langle g_1, g_2 \rangle}.
	\end{align}
	Insbesondere gilt dann auch $ V_{g_1} f_1, V_{g_2} f_2 \in \L^2(\R^{2d}) $.
\end{genericthm}

\begin{proof}
	Wir bemerken zunächst, dass die Operatoren $ \T_a  $ und $ \F_2 $ unitär auf $ \L^2(\R^{2d}) $ sind. 
	Dass $ \T_a  $ unitär ist sehen wir direkt durch Substitution. $ \F_2 $ ist unitär, da die Fouriertransformation auf $ \L^2 $ unitär ist. 
	%\textcolor{red}{\textbf{Warum? Nachweisen!}}
	Damit gilt durch 
	\begin{align*}
	\langle V_{g_1} f_1 , V_{g_2} f_2 \rangle_{\L^2(\R^{2d})}
	= 
	\langle \F_2 \T_a (f_1 \otimes \overline{g_1}),
	\F_2 \T_a (f_2 \otimes \overline{g_2})
	\rangle
	=
	\langle
	f_1 \otimes \overline{g_1}, f_2 \otimes \overline{g_2}
	\rangle
	=
	\langle f_1,f_2 \rangle \overline{ \langle g_1, g_2 \rangle}
	\end{align*}
	die gewünschte Aussage.
\end{proof}

\begin{genericthm}{STFT-Plancherel}\label{th:orthogonality_relation_kor_1}
	Für $ f,g \in \L^2(\R^d) $ gilt
	\begin{align}
	\| V_g f \|_2 = \| f \|_2 \| g \|_2.
	\end{align}
\end{genericthm}

\begin{proof}
	Seien $ f,g \in \L^2(\R^d) $. Dann folgt direkt
	\begin{align*}
	\| V_g f \|_2 = \langle V_g f , V_g f \rangle = 
	\langle f, f \rangle 
	\overline{\langle g, g \rangle}
	=
	\| f \|_2 \| g \|_2
	\end{align*}
	mithilfe des letzten Satzes.
\end{proof}

Also ist die STFT für $ \| g \|_2 = 1 $  wegen 
\begin{align*}
\| V_g f \|_2 = \| f \|_2
\end{align*}
eine Isometrie zwischen $ \L^2(\R^d) $ und $ \L^2(\R^{2d}) $.
Damit ist $ V_g $ insbesondere injektiv und wir können $ f $ eindeutig durch $ V_g f $ bestimmen. Aus der Injektivität erhalten wir außerdem, dass
\begin{align*}
V_g f (x,\omega) = \langle f, M_\omega T_x g \rangle = 0 \
\Rightarrow
\
f = 0
\end{align*}
gilt. Damit gilt für ein fixiertes $ g \in \L^2(\R^d) $:
\begin{align*}
\lbrace
M_\omega T_x g \ | \ x, \omega \in \R^d \rbrace^\perp = \lbrace 0 \rbrace
\end{align*}
Also ist die lineare Hülle von $ \lbrace M_\omega T_x g \ | \ x, \omega \in \R^d \rbrace $ ein dichter Teilraum des $ \L^2(\R^d) $.


\begin{genericthm}{Inversionsformel}\label{th:inversion_formula}
	Seien $ g,\gamma \in \L^2(\R^d) $ mit $ \langle g,\gamma \rangle \neq 0$.
	Dann gilt
	\begin{align}\label{eq:inversion_formula}
	f = 
	\frac{1}{\langle g, \gamma \rangle}
	\int \limits_{\R^{2d}} V_g f(x,\omega) M_\omega T_x \gamma \td{(\omega,x)} 
	\end{align}
	für alle $ f \in \L^2(\R^d) $.
\end{genericthm}

Bevor wir in der Lage sind, diesen Satz zu beweisen, müssen wir schwach interpretierte Integrale diskutieren. Sei $ B $ ein beliebiger Banachraum und eine Abbildung $ g \ : \ \R^d \to B $ gegeben.
Dann wird 
\begin{align*}
f  = \int \limits_{\R^d} g(x) \td{x}
\end{align*}
als
\begin{align*}
l(h) := \langle f,h \rangle 
= 
\int \limits_{\R^d} \langle g(x) , h \rangle \td{x}
\end{align*}
für  alle $ h \in B^\prime $ interpretiert. Ist nun $ l $ ein beschränktes (konjugiert-)lineares Funktional, so ist das Integral als eindeutiges Element in $ B^{\prime \prime} $ geben.
Im Folgenden betrachten wir nur reflexive Räume, d.h. es gilt $ B = B^{\prime \prime} $. Für den Satz betrachten wir Integrale der Form 
\begin{align}\label{eq:unique_inverse_function}
f = \int \limits_{\R^{2d}} F(x,\omega) M_\omega T_x g \td{(\omega,x)}.
\end{align}
Dann ist durch
\begin{align*}
l(h) := \int \limits_{\R^{2d}} F(x,\omega) \overline{\langle h , M_\omega T_x g \rangle} \td{(\omega,x)}
\end{align*}
ein beschränktes konjugiert-lineares Funktional auf $ \L^2(\R^d) $ für $ F \in \L^2(\R^{2d}) $ gegeben.
Die Linearität ist klar und die Beschränktheit folgt durch Anwendung der Cauchy-Schwarz-Ungleichung und \ref{th:orthogonality_relation_kor_1}:
\begin{align}\label{eq:unique_inverse_bounded}
| l(h) |
\leq 
\| F \|_2 \| V_g h \|_2 = \| F \|_2 \| h \|_2 \| g \|_2 < \infty. 
\end{align}
Insgesamt liefert \eqref{eq:unique_inverse_function} also eine eindeutige Funktion in $ \L^2(\R^d) $ mit der Norm $ \| f \|_2 \leq \| F\|_2 \| g \|_2 $.
Nun können wir dazu übergehen, den Satz zu beweisen.

\begin{proof}
	Wegen \ref{th:orthogonality_relation_kor_1} gilt $ V_g f \in \L^2(\R^{2d}) $.
	Mit unseren vorherigen Überlegungen ist 
	\begin{align*}
	\tilde{f} = \frac{1}{\langle \gamma , g \rangle}\int \limits_{\R^{2d}} V_g f(x,\omega) M_\omega T_x \gamma \td{(\omega,x)}
	\end{align*}
	eine eindeutige, wohldefinierte Funktion in $ \L^2(\R^d) $. Unter Verwendung von \eqref{orthogonality_relation} erhalten wir dann
	\begin{align*}
	\langle \tilde{f}, h \rangle 
	&= 
	\frac{1}{\langle \gamma, g \rangle}
	\int \limits_{\R^{2d}} V_g f(x,\omega) \overline{\langle h , M_\omega T_x \gamma \rangle} \td{(\omega,x)}
	=
	\frac{1}{\langle \gamma, g \rangle}
	\langle V_g f , V_\gamma h \rangle\\
	&=
	\frac{1}{\langle \gamma, g \rangle} 
	\langle f, h \rangle  \langle \gamma , g \rangle = \langle f, h \rangle.
	\end{align*}
	Also gilt $ f = \tilde{ f} $.
\end{proof}

Wir definieren uns nun den Operator
\begin{align*}
A_g \ : \ \L^2(\R^{2d}) \to \L^2(\R^d), \quad
F \mapsto  \int \limits_{\R^{2d}} F(x,\omega) M_\omega T_x g \td{(\omega,x)}.
\end{align*}
Wegen 
\begin{align*}
\langle A_g F , h \rangle 
=
\int \limits_{\R^{2d}} F(x,\omega) \langle M_\omega T_x  g , h \rangle \td{(\omega,x)}
= 
\langle F , V_g h \rangle = 
\langle V_g^\prime F , h \rangle 
\end{align*}
ist dieser adjungiert zu der STFT $ V_g  :  \L^2(\R^d) \to \L^2(\R^{2d}) $.
Mit \eqref{eq:unique_inverse_bounded} ist die Beschränktheit auch klar.
Dies liefert uns die Möglichkeit, die Inversionsformel \eqref{eq:inversion_formula} auch durch
\begin{align}\label{eq:inversion_formula_operator}
\frac{1}{\langle \gamma,g \rangle } V^\prime_\gamma V_g = I
\end{align}
zu beschreiben.
\newpage
%\section{Frametheorie}

\begin{df}
	Wir bezeichnen eine Folge $ (e_j)_{j \in J} $ in einem Hilbertraum $ \H $ als \textit{Frame}, wenn
	$ A,B > 0  $ mit
	\begin{align}\label{eq:frame_definition}
	A \| f \|^2 \leq \sum \limits_{j \in J} | \langle f, e_j \rangle |^2 \leq B \|f \|^2
	\end{align}
	für alle $ f \in \H $ gilt. Hierbei sind $ A,B $ die \textit{Frameschranken}.
	Wir nennen einen Frame \textit{dicht}, falls $ A = B  $ gilt.
\end{df}

\begin{df}
	Für jede Teilmenge $ \lbrace e_j \ : \ j \in J \rbrace \subseteq \H $ bezeichnen wir mit
	\begin{equation}\label{coefficients_operator}
	Cf = \lbrace \langle f, e_j\rangle \ : \ j \in J \rbrace
	\end{equation}
	den \textit{Koeffizientenoperator}. Der \textit{Rekonstruktionsoperator} $ D $ ist durch
	\begin{equation}\label{reconstruction_operator}
	D c = \sum \limits_{j \in J} c_j e_j \in \H
	\end{equation}
	für endliche Folgen $ (c_j)_{j \in J} $ definiert. Mit
	\begin{equation}
	S f = \sum\limits_{j \in J} \langle f, e_j \rangle e_j
	\end{equation}
	definieren wir den \textit{Frameoperator} auf $ \H $.
\end{df}

\begin{sz}\label{skript:2.3}
	Sei $ (e_j)_{j \in J} $ ein Frame von $ \H $.
	Dann gelten:
	\begin{enumerate}
		\item 
		$ C : \H \to \ell^2(J) $ ist ein beschränkter Operator mit abgeschlossenem Bild.
		\item 
		Die Operatoren $ C $ und $ D $ sind zueinander adjungiert, d.h. $ C^\ast = D $.
		Damit wird $ D $ zu einem stetigen Operator von $ \ell^2(J) $ nach $ \H $ mit
		\begin{equation}\label{eq:reconstruction_operator_cont_bound}
		\left\| 
		\sum \limits_{j \in J }c_j e_j
		\right\|
		\leq
		\sqrt{B}\|c\|_2
		\end{equation}
		erweitert.
		\item
		Der Frameoperator $ S = C^\ast C = D D^\ast : \H\to \H$ ist positiv und invertierbar mit 
		\begin{equation}\label{eq:frame_operator_inequal}
		A I_\H \leq S \leq B I_\H \ \text{und} \ B^{-1} I_\H \leq S^{-1} \leq A^{-1} I_\H.		
		\end{equation}
		Insbesondere ist $ (e_j)_{j \in J} $ genau dann ein dichter Frame, wenn $ S = A I_\H $ gilt.
		\item
		Die optimalen Frameschranken sind durch
		$ B_{\mathrm{opt}} = \| S \| $ und $ A_{\mathrm{opt}} = \| S^{-1} \|^{-1} $ gegeben.
	\end{enumerate}
\end{sz}

\begin{proof}
	\begin{enumerate}
		\item Die Beschränktheit folgt direkt mit \eqref{eq:frame_definition} durch
		\begin{equation*}
		\left(
		\sum \limits_{j \in J} | \langle f, e_j \rangle |^2
		\right)^{\frac{1}{2}}
		\leq
		\sqrt{B} \|  f \|.
		\end{equation*}
		Aus \eqref{eq:frame_definition} folgt auch sofort die Injektivität von $ C $. Sei nun $ C f_n $ eine konvergente Folge in $ \Bild \ C $. Dann gilt mit \eqref{eq:frame_definition}
		\begin{align*}
		\sqrt{A} \| f_n - f_m \|_\H \leq \| C f_n - C f_m \|_{\ell^2(J)} < \varepsilon
		\end{align*}
		für $ n,m > N_\varepsilon $ und aus der Vollständigkeit von $ \H $ erhalten wir 
		\begin{align*}
		\sqrt{A} \| f_n - f \|_\H < \varepsilon.
		\end{align*}
		Insbesondere gilt dann aber auch
		\begin{align*}
		\| C f_n - C f \|_{\ell^2(J)} \leq \sqrt{B} \| f_n - f \|_\H < \sqrt{\frac{B}{A}} \varepsilon.
		\end{align*}
		Damit ist das Bild von $ C $ abgeschlossen.
		\item 
		Sei $ (c_j)_{j \in J} $ eine endliche Folge.
		Dann gilt 
		\begin{align*}
		\langle C^\prime c , f \rangle
		=
		\langle c , C f \rangle
		=
		\sum \limits_{j \in J} c_j \overline{\langle f , e_j \rangle}
		=
		\left\langle 
		\sum \limits_{j \in J} c_j e_j , f \right\rangle
		=
		\langle D c , f \rangle.
		\end{align*}
		Nun ist $ C $ beschränkt auf $ \H $ mit $ \| C \| \leq \sqrt{B} $.
		Da endliche Folgen dicht in $ \ell^2(J) $ liegen ist auch durch $ C^\prime = D : \ell^2(J) \to H $ ein beschränkter Operator mit derselben Norm gegeben.
		
		\item
		Es gilt $ S = C^\prime C = D D^\prime $, womit $ S : \H \to \H$ selbstadjungiert ist.
		Wegen 
		\begin{align*}
		\langle Sf , f \rangle
		=
		\sum \limits_{j \in J } |\langle f, e_j \rangle|^2
		\end{align*}
		ist $ S $ ein positiver Operator. 
		Es gilt
		\begin{align*}
		\H  = \overline{\Bild \ S} \ \oplus \ (\Bild \ S )^\bot
		\end{align*}
		und wir nehmen an, dass ein $ f \in (\Bild \ S )^\bot \setminus \{ 0\} $ existiert. Dann folgt 
		\begin{align*}
		0 = 
		\langle S f, f \rangle 
		=
		\sum \limits_{j \in J} | \langle f, e_j \rangle |^2,
		\end{align*}
		wodurch wir einen Widerspruch zu $ A > 0  $ erhalten und $ \Bild \ S $ somit dicht in $ \H $ liegt.
		Für $ f \in \Kern \ S $ erhalten wir mit 
		\begin{align*}
		A \| f \|^2 \leq \langle Sf ,f \rangle = \langle 0 , f \rangle = 0
		\end{align*}
		die Injektivität von $ S $. Also können wir die Abbildung
		\begin{align*}
		T : \Bild  \ S \to  \H, Sf \mapsto f
		\end{align*}
		definieren. Diese ist wegen
		\begin{align*}
		\| T S f \|^2 = \| f \|^2 \leq A^{-1} \langle Sf, f \rangle \leq  A^{-1} \|S f \| \| T S f \| 
		\ \Leftrightarrow \
		\|T S f \| \leq  A^{-1} \| S f \|
		\end{align*}
		stetig. Da $ T $ stetig und $ \Bild \ S $ dicht ist, ist $ T $ zu einem stetigen Operator 
		$ \tilde{T} : \H \to \H $ fortsetzbar. Mit dieser Konstruktion erhalten wir
		$ S \tilde{T} = S T = I_\H $ auf $ \H $ und $ \tilde{T} S = T S = I_\H  $ auf $ \Bild \ S $.
		Da $ S T  $ stetig und das Bild dicht ist, folgt auch $ S\tilde{T} = I_\H $ auf ganz $ \H $.
		Also ist $ S $ bijektiv.
		Die Ungleichung $ A I_\H \leq S \leq B I_\H $ ist \eqref{eq:frame_definition} nur umgeschrieben. Daraus folgt direkt
		\begin{align*}
		A S^{-1} \leq I_\H \leq B S^{-1},
		\end{align*}
		da die Operatorungleichung unter Multiplikation mit positiven invertierbaren Operatoren erhalten bleibt.
		\item 
		Wegen 
		\begin{align*}
		\langle Sf , f \rangle \leq B \| f \|^2 
		\ 
		\Leftrightarrow 
		\ 
		\frac{\langle Sf , f \rangle}{\| f \|^2} \leq B
		= \sup_{\| f \| \neq 0 } \frac{\langle Sf , f \rangle}{\| f \|^2} 
		\end{align*}
		erhalten wir $ B_\mathrm{opt} = \| S \|  $.
		Aus $ S^{-1} \leq A^{-1} I_\H $ erhalten mit demselben Argument die Schranke $ A_\mathrm{opt} $.
 	\end{enumerate}
\end{proof}

\begin{lem}
	Sei $ (e_j)_{j \in J} $ ein Frame von $ \H $.
	Angenommen es existiert ein $ c \in \ell^2(J) $ mit $ f = \sum_{j \in J} c_j e_j $.
	Dann existiert zu jedem $ \varepsilon > 0  $ eine endliche Menge $ F_\varepsilon $, sodass 
	\begin{align*}
	\left\| f - \sum \limits_{j \in F} c_j e_j \right\| < \varepsilon
	\end{align*}
	für alle endlichen Teilmengen $  F_\varepsilon \subseteq F $ gilt.
	Wir sagen, dass die Reihe $ \sum_{j \in J} c_j e_j $ \textit{unbedingt} gegen $ f \in \H $ konvergiert.	
\end{lem}

\begin{proof}
	Sei $  \varepsilon > 0  $ beliebig aber fest. Wegen $ c \in \ell^2(J) $ existiert ein $ F_\varepsilon $, sodass 
	\begin{align*}
	\sum \limits_{j \in J} |c_j|^2 - \sum \limits_{j \in F} |c_j|^2 
	=
	\sum \limits_{j \notin F} |c_j|^2 < \frac{\varepsilon}{\sqrt{B}}
	\end{align*}
	für alle endlichen $ F \supseteq F_\varepsilon $ gilt. Wir fixieren ein solches $ F $ und definieren $ c_F = c \cdot \chi_F \in \ell^2(J) $, wodurch
	\begin{align*}
	c_{F,j} =
	\begin{cases}
	1 , &\quad j \in F\\
	0 , &\quad j \notin F
	\end{cases}
	\end{align*}
	gilt. Dann folgt $ D c_F = \sum \limits_{j \in F} c_j e_j $ und
	\begin{align*}
	\left\| f - \sum \limits_{j \in F} c_j e_j \right\|
	=
	\| D ( c- c_F) \|
	\leq 
	 \sqrt{B} \| c - c_F \|_2 < \varepsilon
	\end{align*}
	die Aussage.
\end{proof}

\begin{lem}
	Sei $ (e_j)_{j \in J} $ ein Frame mit dem Schranken $ A,B $. Dann ist der duale Frame mit den Schranken $ B^{-1}, A^{-1} $ durch $ (S^{-1} e_j)_{j \in J} $ gegeben.
	Jedes $ f \in \H $ besitzt die nicht-orthogonalen Entwicklungen
	\begin{align}\label{eq:frames_non_orthogonal_expansion}
	f = 
	\sum \limits_{j \in J} \langle f , S^{-1} e_j \rangle e_j
	=
	\sum \limits_{j \in J} \langle f , e_j \rangle   S^{-1} e_j,
	\end{align}
	wobei beide Summen unbedingt in $ \H $ konvergieren.
\end{lem}

\begin{proof}
		Zunächst erhalten wir 
		\begin{align*}
		\sum \limits_{j \in J}
		| \langle f, S^{-1} e_j \rangle |^2
		=
		\sum \limits_{j \in J}
		| \langle S^{-1} f,  e_j \rangle |^2
		= \langle S ( S^{-1} f ), S^{-1} f \rangle 
		= \langle  S^{-1} f ,  f \rangle
		\end{align*}
		mithilfe der Selbstadjungiertheit und Invertierbarkeit von $ S $.
		Aus \eqref{eq:frame_operator_inequal} erhalten wir
		\begin{align*}
		B^{-1} \| f \|^2 \leq \langle S^{-1} f ,f \rangle 
		\leq A^{-1} \| f \|^2
 		\end{align*}
 		für alle $ f \in \H $. 
		Damit ist $ (S^{-1 } e_j)_{j \in J} $ ein Frame mit den Schranken $ B^{-1} $ und $ A^{-1} $.
		Unter der Verwendung von $ I_\H = S^{-1 } S =
		S S^{-1} $ erhalten wir die Entwicklungen
		\begin{align*}
		f &= S (S^{-1} f) =
		\sum \limits_{j \in J}
		\langle S^{-1} f, e_j \rangle e_j
		=
		\sum \limits_{j \in J}
		\langle  f,  S^{-1}e_j \rangle e_j \\
		f &=
		S^{-1}(S f )=
		S^{-1}\left( \sum \limits_{j \in J} \langle f, e_j \rangle e_j \right)
		=
		\sum \limits_{j \in J} \langle f, e_j \rangle S^{-1} e_j.
		\end{align*}
		Da die Folgen $ ( \langle f, e_j \rangle )_{j \in J} $ und $ ( \langle f, S^{-1} e_j \rangle )_{j \in J} $ in $ \ell^2(J) $ liegen, folgt mit dem letzten Lemma die unbedingte Konvergenz beider Entwicklungen in $  \H $.
\end{proof}

\begin{lem}
	Sei $ (e_j)_{j \in J} $ ein Frame von $ \H $ mit
	$ f = \sum_{j \in J} c_j e_j $ für ein $ c \in \ell^2(J) $.
	Dann gilt $ \| c \|_2 \geq \| a \|_2^2 $ für $ a = (\langle f, S^{-1} e_j \rangle)_{j \in J} $ und die Gleichheit ist genau dann erfüllt, wenn  $ c = a  $ gilt.
\end{lem}

\begin{proof}
	Mit \eqref{eq:frames_non_orthogonal_expansion} und $ f = \sum_{j \in J} c_j e_j $ erhalten wir 
	\begin{align*}
	\langle f, S^{-1} f \rangle
	=
	\sum \limits_{j \in J} \langle f, S^{-1} e_j \rangle \overline{
		\langle f, S^{-1} e_j \rangle
		}
	=\| a \|^2_2 = \langle c, a \rangle.
	\end{align*}
	Durch direktes Nachrechnen folgt dann
	\begin{align*}
	\| c \|^2_2
	=
	\| c -a + a \|_2^2 
	=
	\| c-a\|_2^2 + \|a \|_2^2
	\geq 
	\|a \|_2^2
	\end{align*}
	die gewünschte Aussage.
\end{proof}

\begin{sz}
	Sei $ (e_j)_{j \in J} $ ein Frame von $ \H $.
	Dann sind folgende Aussagen äquivalent:
	\begin{enumerate}[label =\textbf{(\roman*)}]
		\item Die Koeffizienten $ c \in \ell^2(J) $ in der Reihenentwicklung
		\begin{align*}
		f = \sum \limits_{j \in J} \langle f, S^{-1} e_j \rangle e_j
		\end{align*}
		für $ f \in \H $ sind eindeutig.
		
		\item
		Der Koeffizientenoperator $ C : \H \to \ell^2(J) $ ist bijektiv.
		
		\item
		Es existieren $ A^\prime , B^\prime > 0 $, sodass
		\begin{equation}\label{eq:frame_equivalence_1}
		A^\prime \| c \|_2
		\leq
		\left\| \sum \limits_{j \in J} c_j e_j \right\|
		\leq 
		B^\prime \|c\|_2
		\end{equation}
		für alle endlichen Folgen $ c $ erfüllt ist.
		
		\item
		Die Folge $ (e_j)_{j \in J} $ ist das Bild eines Orthonormalsystems $ (f_j)_{j \in J} $
		unter einem beschränkten invertierbaren Operator $ T$.
		
		\item
		Die Grammatrix $ G $, welche durch
		$ G_{jm} = \langle e_m , e_j  \rangle$ für $ j,m \in J $
		gegeben ist,
		definiert einen invertierbaren positiven Operator auf $ \ell^2(J) $.		
		\end{enumerate}
\end{sz}

\begin{proof}
	Zunächst bemerken wir, dass der Operator $ C $
	nach \ref{skript:2.3} beschränkt und injektiv mit abgeschlossenem Bild ist.
	Desweiteren ist $ D $ nach \eqref{eq:frames_non_orthogonal_expansion}	
	surjektiv ist.
	Für jeden beschränkten linearen Operator $ A $
	gilt
	\begin{align*}
	\Bild(A)^\perp = \Kern(A^\prime).
	\end{align*}
	Damit ist $ A^\prime $ genau dann injektiv, wenn $ \Bild(A) $ dicht ist.
	\begin{description}
		\item[\textbf{\textit{ \itshape\textrm{(i)}}} $ \Leftrightarrow $ \textbf{\textit{\textrm{(ii)}}}:]
		Die Eindeutigkeit der Koeffizienten ist äquivalent zur Injektivität des Operators
		\begin{align*}
		D : \ell^2(J) \to \H , c \mapsto
		\sum \limits_{j \in J} c_j e_j.
		\end{align*}
		Wegen $ C^\prime = D $ ist dies äquivalent zu 
		\begin{align*}
		\Bild(C)^\perp = \Kern(D) = \{0\},
		\end{align*}
		womit $ C  $ aufgrund der Abgeschlossenheit des Bildes surjektiv ist.
		
		\item[\textbf{\textit{ \itshape\textrm{(i)}}} $ \Rightarrow $ \textbf{\textit{\textrm{(iii)}}}:]
		Die Stetigkeit von $ D $ liefert mit
		\eqref{eq:reconstruction_operator_cont_bound}
		die Konstante $ B^\prime $.
		Wegen der Eindeutigkeit der Koeffizienten ist $ D $ bijektiv.
		Nach dem Inverse Mapping-Theorem ist dann auch $ D^{-1} : \H \to \ell^2(J) $ stetig.
		Sei nun $ f = \sum_{j \in J} c_j e_j $. Dann erhalten wir mit
		\begin{align*}
		\frac{1}{\| D^{-1} \|} \|c \|_2 \leq \| f \|
		\end{align*}
		die Abschätzung nach unten.
		
		\item[
		\textbf{\textit{ \itshape\textrm{(iii)}}} 
		$ \Rightarrow $ \textbf{\textit{\textrm{(iv)}}}:]
		Sei $ (f_j)_{j \in J} $ ein Orthonormalsystem von $ \H $. Dann gilt
		\begin{align*}
			f = \sum \limits_{j \in J} c_j f_j
		\end{align*}
		mit $ \| f \| = \| c \|_2 $ und wir definieren
		\begin{align*}
		T f := \sum \limits_{j \in J} c_j e_j.
		\end{align*}
		Da jede $ \ell^2 $- Folge durch endliche Folgen approximiert werden kann, folgt mit der Stetigkeit von $ D $ und der Normen, dass \eqref{eq:frame_equivalence_1} auch für beliebige $ c \in \ell^2 $ erfüllt ist. Damit gilt 
		\begin{align*}
		A \|c\|_2 \leq \| T f \| \leq B \| c \|_2,
		\end{align*}
		für $ A,B > 0 $. Also ist $ T $ stetig, invertierbar und es gilt $ T f_j = e_j $.
		
		\item[
		\textbf{\textit{ \itshape\textrm{(iv)}}} 
		$ \Rightarrow $ \textbf{\textit{\textrm{(i)}}}:]
		Sei $ T : \H \to \H $ ein invertierbarer stetiger Operator mit $ T f_j = e_j $ für ein Orthonormalsystem $ (f_j)_{j \in J} $.
		Dann gilt 
		\begin{align*}
		D c = 
		T \left( \sum \limits_{j \in J} c_j f_j\right) T(f) = 0
		\end{align*}
		genau dann, wenn $ c_j = 0 $ für alle $ j \in J $ gilt.
		Damit ist $ D $ injektiv, womit die Koeffizienten eindeutig sind.
		
		\item[
		\textbf{\textit{ \itshape\textrm{(iii)}}} 
		$ \Leftrightarrow $ \textbf{\textit{\textrm{(v)}}}:]
		Zunächst gilt
		\begin{align*}
		\langle Gc ,c \rangle
		=
		\sum \limits_{j,m \in J} \langle e_m, e_j \rangle c_m \overline{c_j}
		=
		\left\| \sum \limits_{m \in J} c_m e_m \right\|^2 
		\end{align*}
		für $ c \in \ell^2(J) $.
		Mit 
		\begin{align*}
		((Gc)_j)_{j \in J}
		= \left( \left\langle \sum \limits_{m \in J } c_m e_m , e_j \right\rangle\right)_{j \in J}
		=C \circ D (c) 
		\end{align*}
		ist $ G = C \circ D $ stetig. Da endliche Folgen dicht in $ \ell^2(J) $ liegen, erhalten wir 
		\begin{align*}
		A \|  c \|^2_2 \leq \langle Gc,c \rangle \leq B \| c \|_2^2,\quad   A,B > 0 
		\end{align*}
		für alle $ c \in \ell^2(J) $. Damit ist $ G $ ein positiver invertierbarer Operator.
		Die Argumentation ist analog zu \ref{skript:2.3}.
		Die Rückrichtung folgt unmittelbar aus der Stetigkeit.
	\end{description}
\end{proof}


\begin{lem}
	\begin{enumerate}
		\item
		Sei $ (e_j)_{j \in J} $ ein dichter Frame mit $ A = B = 1 $ und $ \|e_j \| = 1 $ für alle $ j \in J $.
		Dann ist $ (e_j)_{j \in J} $ ein Orthonormalsystem.
		
		\item
		Sei $ (e_j)_{j \in J} $ ein Frame.
		Dann ist $ (S^{-\frac{1}{2} e_j})_{j \in J} $ ein dichter Frame mit den Schranken $ A = B = 1 $.
		
		\item 
		Sei $ (e_j)_{j \in J} $ ein Frame.
		Dann ist der \textit{inverse Frameoperator} $ S^{-1} $ durch
		\begin{equation}
		S^{-1} f
		=
		\sum \limits_{j \in J} \langle f, S^{-1} e_j \rangle S^{-1} e_j
 		\end{equation}
 		gegeben. Damit ist $ S^{-1} $ der Frameoperator zu dem dualen Frame $ (S^{-1} e_j)_{j \in J} $. 	 
	\end{enumerate}
	
	
\end{lem}

\begin{proof}
	\begin{enumerate}
		\item Mit \eqref{eq:frame_definition} erhalten wir
		\begin{align*}
		1 = \| e_m \|^2 = \sum \limits_{j \in J} | \langle e_m,e_j \rangle|^2 
		= 1 + \sum \limits_{j \neq m} | \langle e_m,e_j \rangle|^2.
		\end{align*}
		Damit gilt $ \delta_{mj} = \langle e_m,e_j \rangle $ und $ (e_j)_{j \in J} $ ist ein Orthonormalsystem.
		
		
		\item
		\textcolor{red}{\textbf{Fehlt.}} 
		
		\item 
		Folgt direkt durch:
		\begin{align*}
		S^{-1} f = S^{-1} S (S^{-1} f) = 
		\sum \limits_{j \in J}
		\langle f, S^{-1} e_j \rangle S^{-1} e_j
		\end{align*}
		
	\end{enumerate}
\end{proof}

\blankpage

\section{$ L^{p,q}_m $-Räume}

In diesem Abschnitt werden wir gewichtete $ \L^p $-Räume mit gemischten Normen betrachten und  übertragen die bekannten Eigenschaften der $ \L^p $-Räume auf diese Räume. 
%Die Definition der Modulationsräume $ \M_m^{p,q} $ hat die Form
Die Modulationsräume werden über die Norm 
\begin{align*}
\| f \|_{\M_m^{p,q}} := \| V_g f \|_{\L_m^{p,q}}
\end{align*}
definiert. 
Also ist es naheliegend, dass wir uns zunächst mit den Eigenschaften von $ \L_m^{p,q} $ beschäftigen.



\begin{df}\label{weight_function}
	Eine Funktion $ v $ auf $ \R^{2d} $ heißt \textit{Gewichtsfunktion}, falls $ v \in \L^1_\mathrm{loc}(\R^{2d}) $ und $ v \geq 0 $ gilt.
	\begin{enumerate}[label =\textbf{(\roman*)}]
		\item 
		Eine Gewichtsfunktion $ v $ heißt \textit{submultiplikativ}, falls
		\begin{equation}\label{submultiplicative_weight}
		v(z_1 + z_2) \leq v(z_1) \cdot v(z_1)
		\end{equation}
		für $ z_1, z_2 \in \R^{2d} $ gilt.
	
		\item
		Eine Gewichtsfunktion $ m $ heißt $ v $-\textit{moderiert}, falls ein $ C > 0 $ existiert, sodass
		\begin{equation}\label{eq:v_moderate}
		m(z_1 + z_2)\leq C v(z_1) m(z_2)
		\end{equation}
		für $ z_1, z_2 \in \R^{2d} $ gilt.
		
		\item
		Die Gewichtsfunktionen $ m_1, m_2 $ heißen \textit{äquivalent}, falls
		ein $ C > 0 $ existiert, sodass
		\begin{equation}\label{equivalence_weights}
		C^{-1} m_1(z) \leq m_2(z) \leq C m_1(z)
		\end{equation}
		für $ z \in \R^{2d} $ gilt. 
		Wir schreiben hierfür $ m_1 \asymp m_2 $.
	\end{enumerate}
\end{df}


Angenommen, es existiert ein $ z_1  \in \R^{2d}$ mit $ m(z_1) = 0 $.
Dann gilt
\begin{align*}
m(z_1 + z_2) \leq C v(z_2) m(z_1) = 0
\end{align*}
für ein beliebiges $ z_2 \in \R^{2d} $ und damit auch $ m= 0 $.
Da eine submultiplikative Gewichtsfunktion $ v $ sich selbst moderiert, werden wir $ v,m > 0 $ voraussetzen.
Wegen $ v(0) \leq (v(0))^2 $ gilt $ v(0) \geq 1 $. 
Insbesondere folgt damit wegen $ v(z - z) \leq (v(z))^2 $
auch $ v(z) \geq 1 $.  

\newpage
\begin{lem}\label{th:submult_cont_symm}
	Sei $ v $ eine submultiplikative Gewichtsfunktion und $ \psi \geq 0  $ eine beliebige kompakt getragene stetige Funktion ungleich der Nullfunktion.
	Dann gelten:
	\vspace{-0.1cm}
	\begin{enumerate}[label =\textbf{(\roman*)}]
		\item $ v \ast \psi $ ist stetig und es gilt $ v \ast \psi \asymp v $.
		\item $ \tilde{v}(x,\omega) := \max\lbrace v(x,\omega),v(-x,\omega), v(x,-\omega), v(-x,-\omega) \rbrace$ ist symmetrisch. Das heißt
		\begin{align}
		\tilde{v}(x, \omega)
		= 
		\tilde{v}(-x,\omega)
		=
		\tilde{v}(x,-\omega)
		=
		\tilde{v}(-x,-\omega)
		\end{align} 
		ist für alle $ x,\omega \in \R^d $ erfüllt.
	\end{enumerate} 
\end{lem}

\begin{proof}
	\begin{enumerate}[label =\textbf{(\roman*)}]
		\item 
		Die Stetigkeit folgt direkt aus
		\begin{align*}
		(v \ast \psi)(z_n) 
		= \int \limits_{\R^{2d}} v(y) \psi(z_n-y) \td{y}
		\ \rightarrow \
		(v \ast \psi)(z_0)
		\end{align*}
		für $ z_n \to z_0 $.
		Wir müssen nun zeigen, dass
		\begin{align*}
		\frac{1}{C} v  \leq  \psi \ast v \leq C v  
		\end{align*}
		gilt.
		Der rechte Teil der Ungleichung folgt durch
		\begin{align*}
		(\psi \ast v)(z) = 
		\int \limits_{\R^{2d}} \psi(y) v(z - y) \td{y}
		\leq
		\underbrace{\max \limits_{y \in \supp  \ \psi}  \psi(y)
			\int \limits_{ \supp  \ \psi} v(-y) \td{y}
			}_{=: C_2} \cdot v(z)
		\end{align*}
		unter Verwendung der Submultiplikativität und der Stetigkeit über dem kompakten Träger.
		Wegen 
		\begin{align*}
		v( 0 ) = v(z- z) \leq v(z) v(-z)
		\ \Leftrightarrow \
		\frac{1}{v(z)} \leq \frac{v(-z)}{v(0)}
		\end{align*}
		erhalten wir $ \nicefrac{1}{v} \in \L^1_\mathrm{loc}(\R^{2d}) $.
		Des Weiteren benötigen wir 
		\begin{align*}
			v(z-y) v(y) \geq v(z-y+y) = v(z)
		\end{align*}
		für die linke Abschätzung. 
		Wir wählen ein $ y_0  $ mit $ \psi(y_0 ) > 0 $.
		Dann existiert ein $ \delta > 0  $ mit $ \psi(y) \geq \frac{1}{2} \psi(y_0) $
 		für alle $ y \in B_\delta(y_0) $.
		Nun folgt:
		\begin{align*}
		(\psi \ast v)(z) 
		\geq 
		\int \limits_{\overline{B_\delta(y_0)}} \psi(y) \frac{v(z)}{v(y)} \td{y}
		\geq 
		\underbrace{\min \limits_{y \in B_\delta(y_0)} \psi(y) 
		\int \limits_{\overline{B_\delta(y_0)}} \frac{1}{v(y)} \td{y}}_{=: C_1}
		 \cdot v(z).
		\end{align*}
		Mit der Wahl $ C := \max \lbrace C_1^{-1}, C_2 \rbrace $ folgt dann die gewünschte Aussage.
		\item Diese Aussage ist klar.
		
	\end{enumerate}
\end{proof}
Nun können wir submultiplikative Gewichtsfunktionen allgemein als stetig und symmetrisch voraussetzen. Wir werden von nun an mit $ v $ eine submultiplikative und mit $ m $ eine $ v $-moderate Gewichtsfunktion bezeichnen.
In dieser Arbeit betrachten wir polynomielle Gewichte, d.h.
\begin{align}\label{eq:poly_weight}
v_s(z) = (1+|z|)^s
=
\left( 1 + \left( \sum \limits_{j=1}^{2d} z_j^2 \right)^{\frac{1}{2}} \right)^s
= (1 +  (|x|^2 + |\omega|^2)^\frac{1}{2})^s
\end{align}
für $ z= (x,w) \in \R^{2d} $ und $ s > 0 $.

\begin{lem}\label{th:equivalence_poly_weight}
	Die Gewichte $ (1 + |x| + |\omega|)^s $ und $ (1 + |z|^2)^{\frac{s}{2}} $
	sind äquivalent zu $ v_s $.
\end{lem}

\begin{proof}
	Wir erhalten die Äquivalenz von $ (1 + |z|^2)^{\frac{s}{2}}  $ und $ (1+|x|+|\omega|)^s $ durch das asymptotische Verhalten von
	\begin{align*}
	\frac{1+ |z|^2}{(1 + |x| + |\omega|)^2}
	=
	\frac{1+ |x|^2 + |\omega|^2}{(1 + |x| + |\omega|)^2}.
	\end{align*}
	Die zweite Äquivalenz folgt analog aus dem Verhalten von 
	\begin{align*}
	\frac{1+ |z|^2}{(1+|z|)^2}. 
	\end{align*}
	Damit sind alle drei Gewichte zueinander äquivalent.
\end{proof}




\begin{lem}\label{korrolary_1_weightfunctions}
	
	\begin{enumerate}[label =\textbf{(\roman*)}]
		\item
		Sei $ m $ eine $ v-$ moderierte Gewichtsfunktion. Dann ist
		\begin{align}\label{eq:v_mod_prop_1}
		\frac{1}{C} \frac{m(z)}{v(t)} \leq m(z- t) \leq Cv(t) m(z)
		\end{align}
		erfüllt. Damit folgt
		\begin{align}\label{eq:v_mod_prop_3}
		\frac{1}{\tilde{C} v(z) } \leq m(z) \leq \tilde{C}   v(z)
		\end{align}
		für $ z \in \R^{2d} $ und es gilt
		\begin{align}\label{eq:v_mod_prop_2}
		\frac{1}{C^\prime} m(l)\leq m(z) \leq C^\prime m(l)
		\end{align}
		für alle $ z \in l + [0,1]^{2d} $ mit $ l \in \Z^{2d} $.
		
		\item 
		Die polynomiellen Gewichte $ v_s $ sind submultiplikativ.
		Desweiteren sind $ v_t $ und $ (v_t)^{-1} $ für $ 0 \leq t \leq s $ von $ v_s $ moderiert.
		
%		\item
%		Sei $ s > 2d $. Dann gilt
%		\begin{equation}\label{eq:poly_weight_1}
%		\left( \frac{1}{v_s} \ast \frac{1}{v_s}\right)(z) \leq  C_s \frac{1}{v_s}(z).
%		\end{equation}
%		 
	\end{enumerate}
	
\end{lem}

\begin{proof}
	\begin{enumerate}[label =\textbf{(\roman*)}]
		\item
		Seien $ z,t \in \R^{2d} $ beliebig.
		Dann gelten nach Definition
		$ m(z-t) \leq C m(z) v(t) $ und
		\begin{align*}
		m(z -t + t) \leq C m(z-t) v(t) 
		\ \Leftrightarrow \
		\frac{1}{C v(t)} m(z) \leq m(z-t). 
		\end{align*}
		Wir erhalten \eqref{eq:v_mod_prop_3} durch $ z = 0 $ und $ t = - w $.
		Sei nun $ z = w + l  $, wobei $ w \in [0,1]^{2d} $ und $ l \in \Z^{2d} $.
		Dann gilt $ m(z) \leq C m(l) v(w) $ und
		\begin{align*}
		m(l - w + w) \leq C m(z) v(w)
		\ \Leftrightarrow \
		\frac{1}{C v(w)} m(l) \leq m(z). 
		\end{align*}
		Mit $ C^\prime := C \max_{w \in [0,1]^{2d}} v(w) $ erhalten wir \eqref{eq:v_mod_prop_2}.
		
		\item
		Sei $ 0 \leq t \leq s $. Dann gilt wegen $ 1 + |z_1 + z_2|\leq (1+|z_1|)(1+|z_2|) $ auch
		\begin{align*}
		v_t(z_1 + z_2)
		\leq
		v_t(z_1) v_t(z_2)
		\leq 
		v_t(z_1) v_s(z_2). 
		\end{align*}
		Damit ist $ v_s $ submultiplikativ und $ v_t $ wird von $ v_s $ moderiert.
		Durch eine Äquivalenzumformung erhalten wir
		\begin{align*}
		(v_t(z_1))^{-1} \leq (v_t(z_1 + z_2))^{-1} v_s(z_2)
		\end{align*}
		mit $ z_1 = w_1 + w_2 $ und $ z_2 = -w_2 $, dass $ (v_t)^{-1} $ von $ v_s $ moderiert wird.
		
%		\item 
%		\textcolor{red}{\textbf{noch ausführen falls nötig.}}
	\end{enumerate}
\end{proof}

\begin{df}\label{df:mixed_norm}
	Sei $ m $ eine Gewichtsfunktion auf $ \R^{2d} $. 
	\begin{enumerate}[label =\textbf{(\roman*)}]
		\item 
		Wir sagen eine messbare Funktion $ F : \R^{2d} \to \C $ liegt
		in dem gewichteten Mischnormraum $ \L^{p,q}_m(\R^{2d}) $, falls
		\begin{align*}
		\| F \|_{\L^{p,q}_m} :=
		\left(		
		\ \int \limits_{\R^{2d}} 
		\left(
		\ \int \limits_{\R^{2d}}
		|F(x,\omega)|^p m(x,\omega)^p \td{x} 
		\right)^{\frac{q}{p}}
		\td{\omega}
		\right)^{\frac{1}{q}} < \infty
		\end{align*}
		erfüllt ist.
		Wenn $ p = q $ gilt, erhalten wir den gewichteten $ \L^p $-Raum auf $ \R^{2d} $ und wir schreiben
		$ \L^p_m(\R^{2d}) := \L^{p,p}_m(\R^{2d}) $.
	\end{enumerate}
	
	\begin{enumerate}
		\item[\textbf{(ii)}]
		Falls $ p = \infty $ oder $ q = \infty $ ersetzen wir die entsprechende Norm durch die $ \infty $-Norm, d.h.
		\begin{align*}
		\| F \|_{\L^{ \infty,q}_m} &:=
		\left(  \
		\int 
		\limits_{\R^{d}}
		\left(
		\esssupp \limits_{x \in \R^d} |F(x,\omega)| m(x,\omega) 
		\right)^{q}
		\td{\omega}
		\right)^{\frac{1}{q}}\\
		\| F \|_{\L^{ p, \infty}_m} &:=
		\esssupp \limits_{\omega \in \R^d}
		\left( \ 
		\int \limits_{\R^{d}}
		|F(x,\omega)|^p m(x,\omega)^p \td{x} 
		\right)^{\frac{1}{p}}.
		\end{align*}
		Für $ p  = q = \infty  $ definieren wir $ L^\infty_m(\R^{2d})  $ über die Norm
		\begin{align*}
		\| F \|_{\L^{ \infty}_m} &:=
		\esssupp \limits_{z \in \R^{2d}}
		|F(z)| m(z).
		\end{align*}
		Auch hier liegt $ F $ in den entsprechenden Räumen, falls die Norm endlich ist.
	\end{enumerate}
\end{df}

Wenn der Zusammenhang klar ist, werden wir im weiteren Verlauf $ \L^{p,q}_m $ anstatt $ \L^{p,q}_m(\R^{2d}) $ schreiben. Analog werden wir mit den anderen Räumen vorgehen.

\begin{sz}\label{hoelder_mixed_norm}
	Seien $ 1 \leq p,q \leq \infty  $, $ F \in \L^{p,q}_m $ und $ H  \in \L^{p^\prime,q^\prime}_{\nicefrac{1}{m}}$
	mit $ \nicefrac{1}{p} + \nicefrac{1}{p^\prime}=1=
	\nicefrac{1}{q} + \nicefrac{1}{q^\prime} $.
	Dann gilt die \textit{Hölder-Ungleichung}
	\begin{align}\label{eq:hoelder_mixed_norm}
	\left|
	\
	\int \limits_{\R^{2d}} 
	F(z) \overline{H(z)} 
	\td{z}
	\right|
	\leq 
	\| F \cdot H \|_{\L^1}
	\leq 
	\| F \|_{\L^{p,q}_m} 
	\| H \|_{\L^{p^\prime, q^\prime}_{\nicefrac{1}{m}}}
	\end{align}
	und $ F \cdot H \in \L^1 $.
\end{sz}

\begin{proof}
	Unser Ziel ist es, die Hölder-Ungleichung direkt zu übertragen. Es gelten\\
	$ F(\cdot, \omega) m(\cdot \omega) \in \L^p(\R^d) $ und $ H(\cdot , \omega) \nicefrac{1}{m(\cdot, \omega)} \in \L^{p^\prime}(\R^d) $ für fast alle $ \omega \in \R^d $.
	Damit erhalten wir durch zweimaliges Anwenden der üblichen Hölder-Ungleichung mit
	\begin{align*}
	\| F \cdot H \|_{\L^1}
	&=
	\int 
	\limits_{\R^d} 
	\int 
	\limits_{\R^d} |F(\cdot, \omega)| m(\cdot,\omega) |H(\cdot,\omega)| \frac{1}{m(\cdot, \omega)} \td{(x,\omega)} \\
	&\leq 
	\int \limits_{\R^d}
	\| F(\cdot , \omega) m(\cdot, \omega) \|_{p}
	\left\| H(\cdot,\omega) \frac{1}{m(\cdot , \omega)} \right\|_{p^\prime} \td{\omega}
	\leq 
%	\| \| F(\cdot , \omega) m(\cdot,\omega) \|_p \|_q
%	\left\| \left\| H(\cdot,\omega) \frac{1}{m(\cdot , \omega)}\right\|_{p^\prime} \right\|_{q^\prime}
	 \| F \|_{\L^{p,q}_m} \| H \|_{\L^{p^\prime,q^\prime}_{\nicefrac{1}{m}}}
	< \infty
	\end{align*}
	die gewünschte Aussage.
\end{proof}

\newpage

\begin{sz}\label{minkowski_mixed_norm}
	Sei $ 1 \leq p,q \leq \infty  $ und $ F,G \in \L^{p,q}_m $.
	Dann gilt mit
	\begin{align*}
	\| F + G \|_{\L^{p,q}_m} \leq \| F \|_{\L^{p,q}_m} + \| G \|_{\L^{p,q}_m} 
	\end{align*}
	die \textit{Minkowski-Ungleichung}.
\end{sz}

\begin{proof}
	Diese Ungleichung folgt wieder durch Anwenden der üblichen Minkowski-Ungleichung. Wir erhalten mit
	\begin{align*}
	\| F + G \|_{\L^{p,q}_m}^q
	&=
	\int 
	\limits_{\R^d}
	\| (F+G)(\cdot, \omega) m(\cdot, \omega) \|_p^q \td{\omega} 
	\leq 
	\int 
	\limits_{\R^d}
	(\| F(\cdot, \omega) m(\cdot, \omega) \|_{p} + \| G(\cdot, \omega) m(\cdot, \omega) \|_{p})^q \td{\omega}\\
	&\leq
	(\| F \|_{\L^{p,q}_m} + \| G\|_{\L^{p,q}_m})^q
	\end{align*}
	die gewünschte Aussage.
\end{proof}


\begin{sz}
	Sei $ 1 \leq p,q \leq \infty  $. Dann ist $ \L^{p,q}_m $ ein Banachraum.
\end{sz}

\begin{proof}
	Sei $ (F_k)_{k \in \N} $ ein Folge in $ \L_m^{p,q} $ und $ \sum_{k=1}^\infty F_k $ absolut konvergent. Wir setzen $ \| \cdot \| := \| \cdot \|_{\L^{p,q}_m} $. 
	Wir betrachten nun 
	\begin{align*}
	(x,\omega) \mapsto \sum \limits_{k=1}^\infty | F_k(x,\omega) | = G(x,\omega)
	\end{align*}
	und setzen $ G_n(x, \omega) = \sum_{k=1}^n | F_k(x,\omega) | $.
	Dann ist $ G_n \in \L^{p,q}_m $ und die Minkowski-Ungleichung liefert
	\begin{align*}
	\| G_n \| \leq \sum \limits_{k=1}^n \| F_k \| 
	\leq \sum \limits_{k=1}^\infty \| F_k \| = a 
	\end{align*}
	mithilfe der absoluten Konvergenz. Unser nächstes Ziel, ist es den Satz der monotonen Konvergenz anzuwenden. Durch zweimaliges Anwenden erhalten wir 
	\begin{align*}
	\int \limits_{\R^d} \left( \ \int \limits_{\R^d}
	|G(x,\omega)|^p m(x,\omega)^p \td{x}
	\right)^{\frac{q}{p}} \td{\omega}
	&=
	\int \limits_{\R^d} 
	\lim \limits_{n \to \infty}
	\left( \ 
	\int \limits_{\R^d}
	|G_n(x,\omega)|^p m(x,\omega)^p \td{x} \right)^{\frac{q}{p}} \td{\omega}\\
	&= \lim \limits_{n \to \infty} \| G_n \|^q \leq a^q,
	\end{align*}
	womit $ G $ bis auf eine Nullmenge $ N $ endlich sein muss. Damit ist auch
	\begin{align*}
	F(x,\omega) :=
	\begin{cases}
	\sum_{k=1}^\infty F_k(x,\omega), &\quad (x,\omega) \in \R^{2d} \setminus N\\
	\quad \quad  0, \quad &\quad (x, \omega) \in N
	\end{cases}
	\end{align*}
	aufgrund der Vollständigkeit von $ \R $ bzw. $ \C $ eine messbare Funktion.
	Wegen 
	\begin{align*}
	\| F \| \leq \| G  \| < \infty
	\end{align*}
	folgt $ F \in \L^{p,q}_m $. Gilt nun noch $ \sum_{k=1}^\infty F_k = F $ bezüglich $ \| \cdot \| $, sind wir fertig.
	Für
	\begin{align*}
	H_n(\cdot, \omega) 
	:= 
	\left| \sum \limits_{k=n}^\infty F_k(\cdot, \omega) \right|^p
	\end{align*}
	gilt $ | H_n(\cdot,\omega) | \leq G(\cdot, \omega)  $, $ G(\cdot, \omega) \in \L^p $ und $ H_n(x, \omega) \to 0 $ für alle $x, \omega \in \R^d $.
	Mit majorisierter Konvergenz erhalten wir dann
	\begin{align*}
	\| H_n(\cdot, \omega) m(\cdot, \omega) \|_p \to 0
	\end{align*}
	für alle $ \omega \in \R^d $. 
	Außerdem gilt $ \| H_n(\cdot, \omega) m(\cdot, \omega)\|_p 
	\leq \| G(\cdot, \omega) \|_p$ und $ \| G(\cdot , \omega ) \|_p\in \L^q $.
	Mit majorisierter Konvergenz folgt dann
	\begin{align*}
	\lim \limits_{n \to \infty}
	\| H_n \|
	=
	\lim \limits_{n \to \infty} 
	\| \| H_n(\cdot , \omega) m(\cdot, \omega) \|_p \|_q
	=0
	\end{align*}
	und somit die gewünschte Aussage.
\end{proof}

\begin{sz}\label{th:dualtiy_mixed_norm}
	Seien $1 \leq  p,q < \infty $ und $ \nicefrac{1}{p} + \nicefrac{1}{p^\prime}=1=
	\nicefrac{1}{q} + \nicefrac{1}{q^\prime} $.
	Dann gilt $ (\L^{p,q}_m)^\prime = \L^{p^\prime,q^\prime}_m $ und die Dualität ist durch
	\begin{align}\label{eq:dualtiy_mixed_norm}
	\langle H,F \rangle
	=
	\int 
	\limits_{\R^{2d}} F(z) \overline{H(z)} \td{z} 
	\end{align} 
	für $ F \in \L^{p,q}_m $ und $ H \in \L^{p^\prime,q^\prime}_{\nicefrac{1}{m}}  $ gegeben.
	
\end{sz}

\begin{proof}
	Wird bewiesen in \cite{benedek1961}.
\end{proof}

\newpage

\begin{sz}
	Seien $ 1 \leq p,q \leq \infty  $ und $ m $ von $ v $ moderiert. Dann ist $  \L^{p,q}_m  $ invariant unter den Translationen $ T_z $ für $ z \in \R^{2d} $ und es gilt
	\begin{align}\label{eq:invariance_translation_mixed_norm}
	\| T_z F \|_{ \L^{p,q}_m }
	\leq C v(z)
	\| F \|_{\L^{p,q}_m}.
	\end{align}
\end{sz}

\begin{proof}
	Sei $ F \in \L^{p,q}_m $.
	Mit $ z = (u, \eta) $ erhalten wir direkt
	\begin{align*}
	\| T_z F \|_{\L^{p,q}_m}
	&= 
	\left(		
	\ \int \limits_{\R^{2d}} 
	\left(
	\ \int \limits_{\R^{2d}}
	|F(x - u ,\omega - \eta )|^p m(x,\omega)^p \td{x} 
	\right)^{\frac{q}{p}}
	\td{\omega}
	\right)^{\frac{1}{q}}\\
	&=
	\left(		
	\ \int \limits_{\R^{2d}} 
	\left(
	\ \int \limits_{\R^{2d}}
	|F(x,\omega)|^p m(x + u,\omega + \eta )^p \td{x} 
	\right)^{\frac{q}{p}}
	\td{\omega}
	\right)^{\frac{1}{q}}\\
	&\leq
	C \left(		
	\ \int \limits_{\R^{2d}} 
	\left(
	\ \int \limits_{\R^{2d}}
	|F(x,\omega)|^p m(x,\omega)^p v(u, \eta)^p \td{x} 
	\right)^{\frac{q}{p}}
	\td{\omega}
	\right)^{\frac{1}{q}}\\
	&=
	C v(z) \|  F \|_{\L^{p,q}_m} < \infty
	\end{align*}
	unter Ausnutzung, dass $ m $ von $ v $ moderiert wird.
\end{proof}

%\begin{sz}\label{attributes_weighted_spaces}
%	Sei $ m $ eine $ v $-moderierte Gewichtsfunktion und $ 1 \leq p,q \leq \infty  $.
%	Dann gelten:
%	\begin{enumerate}
%		\item $ \L^{p,q}_m $ ist ein Banachraum.
%		\item $  \L^{p,q}_m  $ ist invariant unter den Translationen $ T_z $ für $ z \in \R^{2d} $ und es gilt
%		\begin{align*}
%		\| T_z F \|_{ \L^{p,q}_m }
%		\leq C v(z)
%		\| F \|_{\L^{p,q}_m}.
%		\end{align*}
%		\item
%		Sei $ F \in \L^{p,q}_m $ und $ H  \in \L^{p^\prime,q^\prime}_{\nicefrac{1}{m}}$
%		mit $ \nicefrac{1}{p} + \nicefrac{1}{p^\prime}=1=
%		\nicefrac{1}{q} + \nicefrac{1}{q^\prime} $.
%		Dann gilt die Hölderungleichung
%		\begin{align*}
%		\left|
%		\
%		\int \limits_{\R^{2d}} 
%		F(z) \overline{H(z)} 
%		\td{z}
%		\right|
%		\leq 
%		\| F \|_{\L^{p,q}_m} 
%		\| H \|_{\L^{p^\prime, q^\prime}_{\nicefrac{1}{m}}}
%		\end{align*}
%		und $ F \cdot H \in \L^1 $.
%		
%		\item
%		Sei $ p,q < \infty $.
%		Dann gilt $ (\L^{p,q}_m)^\star = \L^{p^\prime,q^\prime}_m $ und die Dualität ist durch
%		\begin{align*}
%		H(F) 
%		= 
%		\langle F,H \rangle
%		=
%		\int 
%		\limits_{\R^{2d}} F(z) \overline{H(z)} \td{z} 
%		\end{align*} 
%		für $ F \in \L^{p,q}_m $ und $ H \in \L^{p^\prime,q^\prime}_{\nicefrac{1}{m}}  $ gegeben.
%	\end{enumerate}
%\end{sz}
%
%\begin{proof}
%	\begin{enumerate} 
%	\item 
%	\textcolor{red}{\textbf{TODO.}}
%	\item
%	Sei $ F \in \L^{p,q}_m $.
%	Mit $ z = (u, \eta) $ erhalten wir direkt
%	\begin{align*}
%	\| T_z F \|_{\L^{p,q}_m}
%	&= 
%	\left(		
%	\ \int \limits_{\R^{2d}} 
%	\left(
%	\ \int \limits_{\R^{2d}}
%	|F(x - u ,\omega - \eta )|^p m(x,\omega)^p \td{x} 
%	\right)^{\frac{q}{p}}
%	\td{\omega}
%	\right)^{\frac{1}{q}}\\
%	&=
%	\left(		
%	\ \int \limits_{\R^{2d}} 
%	\left(
%	\ \int \limits_{\R^{2d}}
%	|F(x,\omega)|^p m(x + u,\omega + \eta )^p \td{x} 
%	\right)^{\frac{q}{p}}
%	\td{\omega}
%	\right)^{\frac{1}{q}}\\
%	&<
%	C \left(		
%	\ \int \limits_{\R^{2d}} 
%	\left(
%	\ \int \limits_{\R^{2d}}
%	|F(x,\omega)|^p m(x,\omega)^p v(u, \eta) \td{x} 
%	\right)^{\frac{q}{p}}
%	\td{\omega}
%	\right)^{\frac{1}{q}}\\
%	&=
%	C v(z) \| T_z F \|_{\L^{p,q}_m} < \infty
%	\end{align*}
%	unter Ausnutzung, dass $ m $ von $ v $ moderiert wird.
%	\item 
%	\textcolor{red}{\textbf{TODO.}}
%	\item
%	\textcolor{red}{\textbf{TODO.}}
%\end{enumerate}
%\end{proof}

\begin{sz}\label{th:convolution_mixed_norm}
	Sei $ m $ von $ v $ moderiert, $ F \in \L^1_v $ und $ G \in \L^{p,q}_m $.
	Dann gilt
	\begin{align*}
	\| F \ast G \|_{\L^{p,q}_m}
	\leq 
	\| F \|_{\L^1_v} \|G\|_{\L^{p,q}_m},
	\end{align*}
	wodurch wir insbesondere
	$ \L^1_v \ast \L^{p,q}_m \subseteq \L^{p,q}_m $ erhalten.
\end{sz}

\begin{proof}
	Sei $ H \in \L^{p^\prime,q^\prime}_{\nicefrac{1}{m}} $. Dann gilt mit Fubini und der Hölder-Ungleichung
	\begin{align*}
	| \langle F \ast G , H \rangle |
	&\leq 
	\int \limits_{\R^{2d}}
	\int \limits_{\R^{2d}}
	| F(w)| |G(z-w)| | H(z)| \td{w} \td{z}
	=
	\int \limits_{\R^{2d}}
	|F(w)|
	\int \limits_{\R^{2d}}
	|T_w G(z)| |H(z) |\td{z} \td{w}\\
	&\leq
	\int \limits_{\R^{2d}}
	|F(w)|  
	\|T_\omega G \|_{\L^{p,q}_m} \| H \|_{\L^{p^\prime,q^\prime}_{\nicefrac{1}{m}}}
	\td{w}
	\leq 
	C \int \limits_{\R^{2d}}
	|F(w)| v(w)\td{w} \ 
	\| G \|_{\L^{p,q}_m} \| H \|_{\L^{p^\prime,q^\prime}_{\nicefrac{1}{m}}}\\
	&=
	C \| F \|_{\L^1_v}\| G \|_{\L^{p,q}_m} \| H \|_{\L^{p^\prime,q^\prime}_{\nicefrac{1}{m}}}.
	\end{align*}
	Wir beachten hier, dass die Einsfunktion von jedem submultiplikativem Gewicht moderiert wird. Insgesamt erhalten wir die Aussage mit 
	\begin{align*}
	\qquad \qquad \qquad \qquad \qquad 
	\| F \ast G \|_{\L^{p,q}_m}
	= \sup \limits_{\| H \| = 1}
	|	\langle F\ast G , H \rangle|
	\leq 
	C \| F \|_{\L^1_v} \|G\|_{\L^{p,q}_m}.
	 \qquad
	\qquad \qquad \qquad \ \ \ \qedhere
	\end{align*}
\end{proof}




\blankpage
\newpage
\section{Zeitfrequenzanalysis auf Distributionen}

In diesem Abschnitt untersuchen wir die STFT angewandt auf Distributionen.
Wir werden sehen, dass die STFT von $ f \in \S^\prime(\R^d) $ eine stetige Funktion und polynomiell beschränkt ist.
Außerdem werden wir mithilfe der STFT eine nützliche Charakterisierung des Schwartzraums entwickeln.
Diese wird uns helfen die Inversionsformel \eqref{eq:inversion_formula} auf den Raum der Distributionen zu erweitern.

\begin{lem}
	Sei $ g \in \S(\R^d) $.
	Dann gilt
	\begin{align}\label{eq:deri_prod_time_freq}
	D^\alpha X^\beta(M_\omega T_x g)
	=
	\sum \limits_{\gamma_1 \leq \alpha}
	\sum \limits_{\gamma_2 \leq \beta}
	\binom{\alpha}{\gamma_1}
	\binom{\beta}{\gamma_2}
	x^{\gamma_2}
	(2 \pi \i \omega)^{\gamma_1}
	M_\omega T_x(D^{\alpha- \gamma_1} X^{\beta-\gamma_2} g)
	\end{align}
	für alle $ (x,\omega) \in \R^{2d} $ und $ \alpha, \beta \in \N^d_0 $.
	
\end{lem}

\begin{proof}
	Es gilt
	\begin{align*}
	X^\beta T_x g(t)
	=
	t^\beta g(t-x)
	=
	T_x ((x+t)^\beta g(t))
	=
	T_x ((x+X)^\beta g)(t).
	\end{align*}
	Mit dem binomischen Lehrsatz erhalten wir
	\begin{align}\label{eq:proof_identity_deri_mult}
	X^\beta T_x g  
	=
	\sum 
	\limits_{\gamma_2 \leq \beta}
	\binom{\beta}{\gamma_2}
	x^{\gamma_2} T_x (X^{\beta - \gamma_2} g).
	\end{align}
	Mit dieser Formel folgt
	\begin{align*}
	\widehat{(D^\alpha M_w  h)}
	&=(2\pi \i)^{|\alpha|}
	X^\alpha T_w \hat{h}
	=
	(2\pi \i)^{|\alpha|}
	\sum \limits_{\gamma_1 \leq \alpha}
	\binom{\alpha}{\gamma_1} \omega^{\gamma_1}T_\omega (X^{\alpha - \gamma_1} \hat{h})\\
	&=
	\sum \limits_{\gamma_1 \leq \alpha}
	\binom{\alpha}{\gamma_1} (2\pi \i \omega)^{\gamma_1} \widehat{(M_\omega D^{\alpha - \gamma_1} h)} 
	\end{align*}
	und somit auch
	\begin{align}
	D^\alpha M_\omega h 
	=
	\sum \limits_{\gamma_1 \leq \alpha}
	\binom{\alpha}{\gamma_1}
	(2 \pi \i \omega)^{\gamma_1} M_\omega D^{\alpha- \gamma_1} h.
	\end{align}
	Dies wenden wir nun auf \eqref{eq:proof_identity_deri_mult} an.
	Dann folgt
	\begin{align*}
	D^\alpha M_\omega X^\beta T_x g
	= 
	D^\alpha X^\beta(M_\omega T_x g)
	=
	\sum \limits_{\gamma_1 \leq \alpha}
	\sum \limits_{\gamma_2 \leq \beta}
	\binom{\alpha}{\gamma_1}
	\binom{\beta}{\gamma_2}
	x^{\gamma_2}
	(2 \pi \i \omega)^{\gamma_1}
	M_\omega D^{\alpha - \gamma_1} T_x X^{\beta - \gamma_2} g.
	\end{align*}
	Da $ M_\omega $ und $ X^\beta $ bzw. $ T_x  $ und $ D^\alpha $ kommutieren, gilt die gewünschte Aussage.
\end{proof}

\begin{lem}
	Die operatorwertige Abbildung $ (x, \omega) \mapsto M_\omega T_x $
	ist stark stetig auf $ \S(\R^d) $ und schwach*-stetig auf $ \S^\prime(\R^d) $.
\end{lem}

\begin{proof}
	Wir müssen zeigen, dass für ein beliebiges $ \varphi \in \S(\R^d) $ alle Seminormen gegen Null gehen, d.h.
	\begin{align}\label{eq:strong_contiunity_to_show}
	\lim \limits_{|x|,|\omega| \to 0}
	\| D^\alpha X^\beta (M_\omega T_x \varphi - \varphi) \|_\infty = 0
	\end{align}
	für alle $ \alpha, \beta \in \N_0^d $. Mit \eqref{eq:deri_prod_time_freq} erhalten wir 
	\begin{align*}
	\| D^\alpha X^\beta (\M_w T_x \varphi - \varphi) \|_\infty
	&=
	\left\|
	\sum \limits_{\gamma_1 \leq \alpha}
	\sum \limits_{\gamma_2 \leq \beta}
	\binom{\alpha}{\gamma_1}
	\binom{\beta}{\gamma_2}
	x^{\gamma_2}
	(2 \pi \i \omega)^{\gamma_1}
	M_\omega T_x(D^{\alpha- \gamma_1} X^{\beta-\gamma_2} \varphi)
	-D^\alpha X^\beta \varphi
	\right\|_\infty\\
	&\leq
	\| M_\omega T_x (D^\alpha X^\beta \varphi) - D^\alpha X^\beta \varphi \|_\infty\\
	&+
	\sum \limits_{0 \neq (\gamma_1,\gamma_2 ) \leq (\alpha,\beta)}
	\binom{\alpha}{\gamma_1}
	\binom{\beta}{\gamma_2}
	x^{\gamma_2}
	(2 \pi \i \omega)^{\gamma_1}
	\|M_\omega T_x(D^{\alpha- \gamma_1} X^{\beta-\gamma_2} \varphi) \|_\infty.
	\end{align*}
	Der hintere Teil geht wegen $ \gamma_1 \neq 0 $ oder $ \gamma_2 \neq 0 $ für $ |x|, |\omega| \to 0 $  gegen Null. 
	Sei $ g \in C^\infty_0(\R^d) $. Dann gilt:
	\begin{align*}
	\| M_\omega T_x g - g \|_\infty
	\leq
	\| M_\omega T_x g - \M_\omega g\|_\infty
	+
	\| M_\omega g - g \|_\infty 
	=
	\|T_x g - g \|_\infty + \| M_\omega g - g\|_\infty
	\end{align*}
	Wir wählen $ r > 1  $ mit $ \supp \ g \in B_{r-1}(0) $. 
	Für den linken Term erhalten wir mithilfe des Mittelwertsatzes
	\begin{align*}
	\|T_x g - g \|_\infty
	= 
	\sup \limits_{t \in B_r(0)} 
	| g(t -x ) - g(t) |
	\leq 
	\sup \limits_{t , \xi  \in B_r(0)} 
	\| \nabla g(\xi) \|_2 \| t -x  -t \|_2
	\leq
	C \| x \|_2 \overset{|x| \to 0}{\to} 0
	\end{align*}
	für alle $ |x| \leq 1 $. Da $ g  $ kompakt getragen ist, gilt für den rechten Term 
	\begin{align*}
	\| M_\omega g - g \|_\infty
	=
	\sup \limits_{t \in \supp \ g}
	| g(t) (e^{2\pi \i t \cdot \omega} - 1) |
	\leq 
	C \sup \limits_{t \in \supp \ g} | e^{2 \pi \i t \cdot \omega} - 1 |.
	\end{align*}
	Wir wählen $ R > 0 $, sodass $ \supp \ g \subseteq [- R, R]^d $ gilt. Dann folgt
	\begin{align*}
	2\pi | w \cdot t | 
	\leq 
	2 \pi \sum \limits_{j =1 }^d | \omega_j| | t_j | \leq
	2 \pi R \|\omega \|_1
	\end{align*}
	für alle $ t \in \supp \ g $. Für $ \delta < \pi  $  und $ \| \omega \|_1 \leq \nicefrac{\delta}{2 \pi R} $ erhalten wir
	\begin{align*}
	| e^ {2 \pi \i \omega \cdot t} - 1 |
	\leq 
	|e^{\i \delta} - 1 | \stackrel{\delta \to 0}{\to} 0 
	\end{align*}
	unabhängig von $ t $.
	Sei nun $ (\varphi_n)_{n \in \N}  $ eine Folge in $ C_0^\infty(\R^d) $ mit 
	$ \varphi_n \to \varphi $.
	Insgesamt folgt dann \eqref{eq:strong_contiunity_to_show} aus
	\begin{align*}
	\| M_\omega T_x D^\alpha X^\beta \varphi - D^\alpha X^\beta \varphi \|
	\leq
	2 \|D^\alpha X^\beta (\varphi_n - \varphi) \|_\infty
	+
	\| M_\omega T_x  D^\alpha X^\beta \varphi_n - D^\alpha X^\beta \varphi_n \|_\infty
	\to 0 
	\end{align*}
	für $ |x|,|\omega| \to 0 $ und $ n \to \infty $.
	Sei nun $ f \in \S^\prime(\R^d) $ und $ \varphi \in \S(\R^d) $.
	Dann gilt
	\begin{align*}
	\lim \limits_{|x|, |\omega| \to 0}
	\langle M_\omega T_x f , \varphi \rangle
	=
	\lim \limits_{|x|, |\omega| \to 0}
	\langle f, T_{-x} M_{-\omega} \varphi \rangle 
	= 
	\langle f, \varphi \rangle
	\end{align*}
	womit $ M_\omega T_x $ schwach*-stetig ist.
	\end{proof}
	
	\begin{sz}\label{th:STFT_continous}
		Sei $ g \in \S(\R^d) \setminus \{0\} $ und $ f \in \S^\prime(\R^d) $.
		Dann ist die STFT $ V_g f $ stetig und es existieren $ C > 0  $, $ N \geq 0 $, sodass
		\begin{align}
		|V_g f(x,\omega) |
		\leq 
		C (1+|x| + |\omega|)^N
		\end{align}
		für alle $ x, \omega \in \R^d $ gilt.		
	\end{sz}
	
\begin{proof}
	Sei $ (x,\omega) \in \R^{2d} $ beliebig und $ (x_n, \omega_n)  $ eine Folge mit $ (x_n,\omega_n ) \to (x,\omega) $. Dann gilt mit \eqref{eq:cont_distribution}:
	\begin{align*}
	&|V_g f (x_n,\omega_n) - V_g f(x,\omega)|
	=
	| \langle f, M_{\omega_n} T_{x_n} g - M_\omega T_x g \rangle |
 	\leq
 	\sum 
 	\limits_{
 		\substack{| \alpha | \leq M_1 \\
	 			| \beta | \leq M_2}
 		}
 		\| D^\alpha X^\beta ( M_{\omega_n} T_{x_n} g - M_\omega T_x g) \|_\infty\\
 	&\qquad \leq
 	\sum 
 	\limits_{
 		\substack{| \alpha | \leq M_1 \\
 			| \beta | \leq M_2}
 	}
 	\| D^\alpha X^\beta (M_{\omega_n} T_{x_n} g - M_{\omega_n} T_x g) \|_\infty
 	+
 	\sum 
 	\limits_{
 		\substack{| \alpha | \leq M_1 \\
 			| \beta | \leq M_2}
 	}
 	\| D^\alpha X^\beta (M_{\omega_n} T_{x} g - M_{\omega} T_x g) \|_\infty.
 	\end{align*}
	Wir werden beide Summen seperat untersuchen. Mit \eqref{eq:deri_prod_time_freq} und \eqref{eq:strong_contiunity_to_show} folgt
	\begin{align*}
	&\| D^\alpha X^\beta (M_{\omega_n} T_{x_n} g - M_{\omega_n} T_x g) \|_\infty
	=
	\| D^\alpha X^\beta (M_{\omega_n} T_x (T_{x_n-x} g -  g)) \|_\infty\\
	&\leq
	\sum 
	\limits_{
		\substack{\gamma_1 \leq \alpha
		\\
		\gamma_2 \leq \beta}
		}
		\binom{\alpha}{\gamma_1}
		\binom{\beta}{\gamma_2}
		|x^{\gamma_2}
		(2  \pi \i \omega_n)^{\gamma_2}|
		\| M_{\omega_n} T_x(D^{\alpha-\gamma_1} X^{\beta - \gamma_2}
		(T_{x_n-x} g - g)) \|_{\infty}
		\\
	&\leq 
	C
	\sum 
	\limits_{
		\substack{\gamma_1 \leq \alpha
			\\
			\gamma_2 \leq \beta}
	}
	\binom{\alpha}{\gamma_1}
	\binom{\beta}{\gamma_2}
	\| (D^{\alpha-\gamma_1} X^{\beta - \gamma_2}
	(M_0T_{x_n-x} g - g) \|_{\infty}
	\to 0
	\end{align*}
	für $ (x_n,\omega_n ) \to (x,\omega) $.
	Ähnliches Vorgehen liefert
	\begin{align*}
	\| D^\alpha X^\beta (M_{\omega} T_0 ( M_{\omega - \omega_n} T_x g - T_x g)) \|_\infty
	\leq ... 
	\leq
	C 
	\sum \limits_{\gamma_2 \leq \beta}
	\binom{\beta}{\gamma_2}
	\| D^\alpha X^{\beta- \gamma_2}(M_{\omega_n - \omega} T_x g - T_x g) \|_\infty
	\to 0
	\end{align*}
	für den zweiten Summanden.
	Insgesamt erhalten wir mit
	\begin{align*}
	(x_n,\omega_n) \to (x,\omega) \ 
	\Rightarrow 
	\
	V_g f(x_n, \omega_n) \to V_g f (x, \omega)
	\end{align*}
	die Stetigkeit von $ V_g f $.
	Die Stetigkeitsabschätzung \eqref{eq:cont_distribution} liefert uns
	\begin{align*}
	| V_g f(x,\omega) |
	=
	|\langle f, M_\omega T_x g \rangle |
	\leq
	C
	\sum \limits_{\substack{|\alpha| \leq M_1\\
		|\beta| \leq M_2}
		}
	\| D^\alpha X^\beta (M_\omega T_x g)\|_\infty
	\end{align*}
	und mit \eqref{eq:deri_prod_time_freq} und der Dreiecksungleichung erhalten wir:
	\begin{align*}
	|V_g f(x,\omega)\| 
	&\leq 
	C
	\sum \limits_{
		\substack{|\alpha| \leq M_1 \\ |\beta| \leq M_2}
		}
	\sum \limits_{
		\substack{\gamma_1 \leq \alpha\\ \gamma_2 \leq \beta}
		}
	\binom{\alpha}{\gamma_1} \binom{\beta}{\gamma_2}
	|x^{\gamma_2}| |(2\pi \i \omega)^{\gamma_1} | 
	\| D^{\alpha-\gamma_1} X^{\beta - \gamma_2} g \|_\infty\\
	&\leq
	C
	\max \limits_{
		\substack{|\alpha| \leq M_1 \\ |\beta| \leq M_2}
		}
	\| D^\alpha X^\beta g \|_\infty
	\sum \limits_{
		\substack{|\alpha| \leq M_1 \\ |\beta| \leq M_2}
	}
	\sum \limits_{
		\substack{\gamma_1 \leq \alpha\\ \gamma_2 \leq \beta}
	}
	\binom{\alpha}{\gamma_1} \binom{\beta}{\gamma_2}
	|x^{\gamma_2}| |(2\pi \i \omega)^{\gamma_1} |.
	\end{align*}
	Die rechte Doppelsumme ist ein Polynom vom Grad $ M_1 $ in den Variablen $ |x_j| $ und ein Polynom vom Grad $ M_2 $ in den Variablen $ |\omega_j| $.
	Wir setzen $ N := \max\{M_1,M_2\} $. Dann existiert ein $ K > 0 $, sodass
	die Doppelsumme durch $ K (1+ |x| + |\omega|)^N $ beschränkt ist.\\
	\qedhere
\end{proof}
\vspace{-2.9cm}
\begin{lem}\label{th:rapid_decay_implies_schwartz}
Sei $ g \in \S(\R^d) \setminus \{ 0\} $ fixiert und $ F(x, \omega) $ schnell fallend auf $ \R^{2d} $, d.h. für alle $ n \geq 0 $ existiert ein $ C_n > 0 $, sodass 
\begin{align}\label{eq:rapid_decay}
| F(x,\omega) |\leq C_n (1 + |x| + |\omega|)^{-n}
\end{align}
erfüllt ist. Dann ist durch
\begin{align}\label{eq:schwartz_func_rapid_decay}
f(t) = \int \limits_{\R^{2d}} F(x, \omega) M_\omega T_x g(t) \td{(\omega,x)}
\end{align}
eine Schwartzfunktion definiert.
\end{lem}

\begin{proof}
	Wegen \eqref{eq:rapid_decay} und $ g \in \S(\R^d) $ gilt
	\begin{align*}
	\int \limits_{\R^{2d}} |F(x,\omega)| |T_x g(t) | \td{(\omega,x)}
	&\leq 
	\int \limits_{\R^{2d}} | F(x,\omega) | \| g \|_\infty \td{(\omega,x)}\\
	&\leq
	C_n \|g \|_\infty 
	\int \limits_{\R^{2d}} \frac{1}{(1+ |x| +|\omega|)^n } \td{(\omega,x)} < \infty
	\end{align*}
	für ein hinreichend großes $ n $. Dies gilt auch, wenn wir $ F  $ mit einem beliebigen Polynom multiplizieren. Also konvergiert \eqref{eq:schwartz_func_rapid_decay} absolut und wir können Differentation nach $ t $ in das Integral ziehen.
	Mit \eqref{eq:deri_prod_time_freq} gilt dann
	\begin{align*}
	D^\alpha X^\beta f(t)
	&=
	\int \limits_{\R^{2d}} F(x, \omega) D^\alpha X^\beta M_\omega T_x g(t) \td{(\omega,x)}\\
	&=
	\sum \limits_{
		\substack{\gamma_1 \leq \alpha \\ \gamma_2 \leq \beta}
		}
	\binom{\alpha}{\gamma_1} \binom{\beta}{\gamma_2}
	\int \limits_{\R^{2d}} F(x,\omega) x^{\gamma_2} (2 \pi \i \omega)^{\gamma_1}
	M_\omega T_x D^{\alpha - \gamma_1} X^{\beta - \gamma_2} g(t)\td{(\omega,x)}.
	\end{align*}
	Wegen $ g \in \S(\R^d) $ gilt $ C := \max_{\gamma_1 \leq  \alpha, \gamma_2 \leq \beta} \| M_\omega T_x D^\alpha X^\beta  g \|_\infty < \infty  $.
	Damit erhalten wir
	\begin{align}\label{eq:estimate_rapid_decay}
	\| D^\alpha X^\beta f \|_\infty
	\leq
	C\int \limits_{\R^{2d}} |F(x, \omega) | P(x,\omega) \td{(\omega,x)} < \infty
	\end{align}
	mit
	\begin{equation}
	\begin{split}
	P(x,\omega)
	&= \sum 
	\limits_{\gamma_1 \leq \alpha}
	\sum \limits_{\gamma_2 \leq \beta }
	\binom{\alpha}{\gamma_1} \binom{\beta}{\gamma_2}
	|x|^{\gamma_1}|2\pi \i \omega|^{\gamma_2}\\
	&=
	\left( \sum 
	\limits_{\gamma_1 \leq \alpha} \binom{\alpha}{\gamma_1} |x|^{\gamma_1} \right)
	\left(\sum \limits_{\gamma_2 \leq \beta } \binom{\beta}{\gamma_2} |2\pi \i \omega|^{\gamma_2} \right)
	=
	\prod \limits_{j=1}^d (1+ |x_j|)^{\alpha_j} (1+ 2\pi |\omega_j|)^{\beta_j}
	\end{split}
	\end{equation}
	woraus $ f \in \S(\R^d) $ folgt.
\end{proof}

\begin{sz}\label{th:equivalence_STFT_bounded}
	Sei $ g \in \S(\R^d) \setminus \{0\} $ fest und $ f \in \S^\prime(\R^d) $.
	Dann sind äquivalent:
	\begin{enumerate}[label =\textbf{(\roman*)}]
		\item $ f \in \S(\R^d) $
		\item $ V_g f \in \S(\R^{2d}) $
		\item Für alle $ n \geq 0 $ existiert ein $ C_n > 0 $, sodass
		für alle $ x, \omega \in \R^d $ gilt:
		\begin{align}
		| V_g f(x,\omega) | 
		\leq 
		C_n (1 + |x| + | \omega|)^{-n}.
		\end{align}
		%für alle $ x, \omega \in \R^d $ gilt.
	\end{enumerate}
\end{sz}

\begin{proof}
	\begin{description}
		\item[\textit{ \itshape\textrm{(i)}} $ \Rightarrow $ \textbf{\textit{\textrm{(ii)}}}:]
		Nach \ref{th:iso_tensor_distri} gilt $ V_g f = \mathcal{F}_2 \mathcal{T}_a (f \otimes \overline{g}) \in \S(\R^{2d})  $.
		\item[\textit{ \itshape\textrm{(ii)}} $ \Rightarrow $ \textbf{\textit{\textrm{(iii)}}}:]
		Da $ V_g f \in \S(\R^{2d}) $, folgt direkt: Für alle $ n \geq 0 $ existiert ein $ C_n > 0 $, sodass
		\begin{align*}
		| V_g f(x,\omega)| = |V_g f (z)| 
		\leq 
		C_n(1+ |z|)^{-n}
		\end{align*}
		für alle $ z \in \R^{2d} $ gilt. Mit \ref{th:equivalence_poly_weight} folgt dann die Aussage.
		\item[\textit{ \itshape\textrm{(iii)}} $ \Rightarrow $ \textbf{\textit{\textrm{(i)}}}:]
		Wegen $ \S(\R^d) \subset \L^2(\R^d) $ gilt $ \|g \|_2 < \infty .$
		Sei 
		\begin{align*}
		\tilde{f} = \frac{1}{\|g \|_2^2}
		\int \limits_{\R^{2d}}
		V_g f(x,\omega) M_\omega T_x g \td{(x,\omega)}.
		\end{align*}
		Dann sind die Voraussetzungen für \ref{th:rapid_decay_implies_schwartz} erfüllt, womit $ \tilde{f} \in \S(\R^d) $ gilt. Die Inversionsformel der STFT \eqref{eq:inversion_formula} liefert uns $ f = \tilde{f} \in \S(\R^d) $.
	\end{description}
\end{proof}

\begin{lem}\label{th:equivalence_seminorms}
	Sei $ g \in \S(\R^d) \setminus \{ 0\}$. Dann ist die Familie von Seminormen
	\begin{align}\label{eq:equivalence_seminorms_1}
	\| V_g f \|_{\L^\infty_{v_s}} 
	= 
	\sup \limits_{z \in \R^{2d}} (1 + |z|)^s |V_g f(z) |
	\end{align}
	für $ s \geq 0  $ eine äquivalente Familie von Seminormen für $ \S(\R^d) $.
\end{lem}

\begin{proof}
	Wir definieren
	\begin{align*}
	\tilde{S}(\R^d)
	:= \left\{
	f \in \L^2(\R^d) \ | \ \forall_{ s \geq 0} : \sup \limits_{z \in \R^{2d}} (1+ |z|)^s| V_g f(z)|< \infty
	\right\}.
	\end{align*}
	Für $ f \in \tilde{S}(\R^d) $ gilt mit \textbf{\textit{\itshape\textrm{(iii)}}} $ \Rightarrow $ \textbf{\textit{\textrm{(i)}}} aus \ref{th:equivalence_STFT_bounded} direkt $ f \in \S(\R^d) $.
	Sei nun $ f \in \S(\R^d) $ und $ s \geq 0  $ beliebig. Dann gilt
	\begin{align*}
	|V_g f(z) | \leq C_n \frac{1}{(1+ |z|)^n} 
	\end{align*}
	für $ n = s$, womit $ f \in \tilde{\S}(\R^d) $ folgt.
	Also ist $ \S(\R^d) = \tilde{\S}(\R^d) $ und wir müssen zeigen, dass beide dieselbe Topologie besitzen.
	Wir wenden die Abschätzung \eqref{eq:estimate_rapid_decay} auf die Inversionsformel
	\begin{align*}
	f = \frac{1}{\|g\|_2^2} \int \limits_{\R^{2d}}
	V_g f(x,\omega) M_\omega T_x g \td{(x,\omega)}
	\end{align*}
	an. Damit erhalten wir
	\begin{align*}
	\| D^\alpha X^\beta f \|_\infty
	\leq
	C \int \limits_{\R^{2d}} |V_g f(z)| P(z) \td{z}
	\leq 
	C^\prime \int \limits_{\R^{2d}} |V_g f(z)| (1+ |z|)^n \td{z}.
	\end{align*}
	für ein hinreichend großes $ n $.
	Dieses $ n $ hängt nur von $ \alpha $ und $ \beta $ ab.
	 Weiter gilt
	\begin{align}\label{eq:equivalence_seminorms_2}
	\| D^\alpha X^\beta f \|_\infty \leq
	C^\prime 
	\sup \limits_{z \in \R^{2d}} | V_g f(z) | (1+ |z|)^{n + 2d + 1}
	\int \limits_{\R^{2d}} (1+ |z|)^{-2d - 1} \td{z} 
	= K \| V_g f \|_{\L^\infty_{v_{n+2d+ 1}}} < \infty.
	\end{align}
	Wir betrachten nun die Identität $ I :\tilde{\S}(\R^d) \to \S(\R^d) $.
	Sei nun $ f_m \to f  $ in $ \tilde{\S}(\R^d) $, d.h.
	\begin{align*}
	\lim \limits_{m \to \infty} 
	\| V_g f_m - V_g f \|_{\L^\infty_{v_s}}
	=
	\| V_g (f_m - f) \|_{\L^\infty_{v_s}} =0 
	\end{align*}
	für alle $ s \geq 0 $. Dann ist $ I $ wegen
	\begin{align*}
	\| D^\alpha X^\beta (f_m - f) \|_\infty \leq
	 K \| V_g (f_m -f) \|_{\L^\infty_{v_{n+2d+ 1}}}
	\to 0
	\end{align*}
	 für $ m \to \infty $ stetig. Mit dem Open-Mapping-Theorem ist $ I $ eine offene Abbildung und somit ein linearer Homöomorphismus. Damit sind beide Topologien gleich.
\end{proof}

\begin{sz}
	Seien $ g, \gamma \in \S(\R^d) \setminus \{0\} $.
	\begin{enumerate}[label =\textbf{(\roman*)}]
		\item 
		Sei $ |F(x,\omega)| \leq C (1+|x| + |\omega|)^N $ für ein $ C,N \geq 0 $.
		Dann wird durch 
		\begin{align}
		f 
		=
		\int \limits_{\R^{2d}}
		F(x,\omega) M_\omega T_x g \td{(x,\omega)}
		\end{align}
		eine temperierte Distribution $ f \in \S^\prime(\R^d) $ im Sinne von 
		\begin{align}
		\langle f, \varphi \rangle
		= 
		\int \limits_{\R^{2d}}
		F(x,\omega)  \langle M_\omega T_x g , \varphi\rangle \td{(x,\omega)}, \quad \varphi \in \S(\R^d)
		\end{align}
		definiert.
		\item
		Sei $ f \in \S^\prime(\R^d) $, $ F = V_g f $ und $ \langle \gamma, g \rangle \neq 0 $. Dann gilt
		\begin{align}\label{eq:inverse_STFT_distribution}
		f = \frac{1}{\langle \gamma, g \rangle}\int \limits_{\R^{2d}}  V_g f(x,\omega) M_\omega T_x \gamma  \td{(\omega,x)}.
		\end{align}
		\end{enumerate}
	Insbesondere ist die STFT injektiv auf $ \S^\prime(\R^d) $.
\end{sz}
\begin{proof}
	\begin{enumerate}[label =\textbf{(\roman*)}]
		\item 
		Zunächst bemerken wir, dass $ |F(z) | = \mathcal{O}(|z|^N) $ und 
		$ |V_g \varphi(z)| \in \mathcal{O}(|z|^{-n}) $ für alle $ n \geq 0 $, $ \varphi \in \S(\R^d) $ gilt. Damit konvergiert das Integral
		\begin{align*}
		\int \limits_{\R^{2d}}
		F(x,\omega) \langle M_\omega T_x g , \varphi\rangle \td{(x,\omega)}
		=
		\int \limits_{\R^{2d}}
		F(x,\omega) \overline{V_g  \varphi(x,\omega)} \td{(x,\omega)}
		\end{align*}
		absolut. Mit den selben Abschätzungen wie in \eqref{eq:equivalence_seminorms_2} erhalten wir
		\begin{align*}
		|\langle f , \varphi \rangle |
		\leq 
		C \sup \limits_{  z \in \R^{2d}}
		|V_g \varphi(z) | (1 + |z|)^{N+2d+1} \int \limits_{\R^{2d}} \frac{1}{(1+ |z|)^{2d + 1}} \td{z}
		\end{align*}
		für alle $ \varphi \in \S(\R^d) $.
		Wegen \ref{th:equivalence_seminorms} gilt $ f \in \S^\prime(\R^d) $.
		
		\item
		Nach \textbf{(i)} wird durch
		\begin{align*}
		\langle \tilde{f}, \varphi \rangle 
		= 
		\frac{1}{\langle \gamma,g \rangle} 
		\int \limits_{\R^{2d}}
		V_g f(x,\omega) \langle M_\omega T_x \gamma ,\varphi \rangle
		\td{(x,\omega)}
		\end{align*}
		eine temperierte Distribution $ \tilde{f} $ definiert. 
		Weiter gilt 
		\begin{align*}
		\varphi
		=
		\frac{1}{\langle g, \gamma \rangle}
		\int \limits_{\R^{2d}}
		V_\gamma \varphi(x,\omega) M_\omega T_x g \td{(x,\omega)}
		\end{align*}
		auf $ \S(\R^d) $. Damit erhalten wir
		\begin{align*}
		\langle f , \varphi \rangle
		&=
		\left\langle f,  \frac{1}{\langle g, \gamma \rangle}
		\int \limits_{\R^{2d}}
		V_\gamma \varphi(x,\omega) M_\omega T_x g \td{(x,\omega)}
		\right\rangle
		=
		\frac{1}{\langle \gamma, g \rangle}
		\int \limits_{\R^{2d}}
		\langle f, V_\gamma \varphi(x,\omega) M_\omega T_x g \rangle \td{(x, \omega)}\\
		&=
		\frac{1}{\langle \gamma, g \rangle}
		\int \limits_{\R^{2d}}
		\langle f,  M_\omega T_x g \rangle \overline{V_\gamma \varphi(x,\omega)}\td{(x, \omega)}
		=
		\frac{1}{\langle \gamma, g \rangle}
		\int \limits_{\R^{2d}}
		V_g f(x,\omega) \langle M_\omega T_x \gamma, \varphi \rangle \td{(x, \omega)}\\
		&=
		\langle \tilde{f}, \varphi \rangle.
		\end{align*}
		Also gilt $ \tilde{f} = f $.
		
	\end{enumerate}
\end{proof}
\newpage
\section{Modulationsräume}
In diesem Abschnitt werden wir elementare Eigenschaften der Modulationsräume und
bereits einige Tatsachen zur Feichtinger-Algebra $ S_0(\R^d) = \M^1 $ entwickeln.
Hierbei werden wir auf die Theorie der ersten drei Abschnitte zurückgreifen. 


\begin{df}
	Sei $ g \in \mathcal{S}(\R^d) \setminus \{ 0\}$ eine Fensterfunktion, $ m $ eine $ v $-moderierte Gewichtsfunktion auf $ \R^{2d} $ und $ 1 \leq p,q \leq \infty $.
	Der Modulationsraum $ \M^{p,q}_m (\R^d) $ besteht aus allen temperierten Distributionen $ f \in \mathcal{S}^\prime(\R^d) $, sodass
	\begin{align*}
	\| f \|_{\M^{p,q}_m} := \| V_g f \|_{\L^{p,q}_m} < \infty
	\end{align*}
	erfüllt ist.
	
\end{df}

Diese Definition hängt von der Wahl des Fensters $ g $ ab. 
Im weiteren Verlauf werden wir zeigen, dass $ \M^{p,q}_m $ unabhängig von der Wahl der Fensterfunktion ist.
Bis zu diesem Zeitpunkt werden wir eine fixierte Standardfensterfunktion $ g_0 \in \S(\R^d) \setminus \{ 0\} $ festlegen und die $ \M^{p,q}_m $-Norm bezüglich $ g_0 $ betrachten.

\begin{kor}
	Sei $ g \in \S(\R^d) \setminus \{ 0\}  $ und $ v_s(z) = (1+ |z|)^s$. Dann gilt
	\begin{align}
	\S(\R^d) =
	\bigcap \limits_{s \geq 0} M^\infty_{v_s} \ \text{und} \
	\S^\prime(\R^d) = \bigcup \limits_{s \geq 0} M^\infty_{\nicefrac{1}{v_s}}.
	\end{align}
\end{kor}

\begin{proof}
	Die Äquivalenz zwischen \textbf{(i)} und \textbf{(iii)} in \ref{th:equivalence_STFT_bounded} liefert direkt
	\begin{align*}
	\S(\R^d) =
	\bigcap \limits_{s \geq 0} M^\infty_{v_s}.
	\end{align*}
	Nach \ref{th:STFT_continous} gilt für $ f \in \S^\prime(\R^d) $ auch $ f \in \bigcup_{s \geq 0} M^\infty_{\nicefrac{1}{v_s}} $ und somit 
	\begin{align*}
	\S^\prime(\R^d) = \bigcup \limits_{s \geq 0} M^\infty_{\nicefrac{1}{v_s}}.
	\end{align*}
\end{proof}
Um die Eigenschaften der Modulationsräume zu zeigen, benötigen wir zuerst eine Inversionsformel, um auf die Eigenschaften der $ \L^{p,q}_m $-Räume zurückzugreifen. Hierfür definieren wir den adjungierten Operator zu $ V_g $ formal.

\begin{df}
	Sei $ \gamma \in \S(\R^d) \setminus \{ 0\} $ eine Fensterfunktion und $ F : \L^{p,q}_m$ .
	Dann definieren wir
	\begin{align*}
	V_\gamma^\prime F 
	:=
	\int \limits_{\R^{2d}}F(x,\omega)  M_\omega T_x \gamma \td{(x,\omega)}.
	\end{align*}
\end{df}
Dieses Integral wird schwach interpretiert, d.h.
\begin{align*}
\langle V^\prime_\gamma F ,f \rangle
=
\int \limits_{\R^{2d}} F(x,\omega) \langle M_\omega T_x \gamma, f \rangle \td{(x,\omega)}
=
\int \limits_{\R^{2d}} F(x,\omega) \overline{V_\gamma f(x,\omega)} \td{(x,\omega)}
=
\langle F , V_\gamma f \rangle.
\end{align*}


\begin{sz}\label{th:inversions_formula_modulation}
	Sei $ m $ eine $ v $-moderierte Gewichtsfunktion und $ \gamma \in \S(\R^d) \setminus \{0\} $.
	Dann gilt:
	\begin{enumerate}[label =\textbf{(\roman*)}]
		\item 
		Es gilt $ \Bild \ V_\gamma^\prime \subseteq \M_m^{p,q} $ und 
		%$ V_\gamma^\prime  $ bildet $ \L^{p,q}_m(\R^{2d}) $ auf $ \M^{p,q}_m $ ab und erfüllt
		\begin{align}\label{eq:inverse_estimate_modulation}
		\| V^\prime_\gamma F \|_{\M^{p,q}_m}
		\leq
		C \| V_{g_0} \gamma \|_{\L^1_v} \| F \|_{\L^{p,q}_m}. 
		\end{align}
		ist erfüllt.
		\item
		Für $ F = V_g f $ und $ \langle \gamma, g \rangle \neq 0 $ gilt die Inversionsformel
		\begin{align}\label{eq:inversion_formula_modulation}
		f = \frac{1}{\langle \gamma, g \rangle}
		\int \limits_{\R^{2d}}V_g f(x,\omega) M_\omega T_x \gamma \td{(\omega,x)}
		\end{align}
		auf $ \M^{p,q}_m $. Wir können hierfür auch
		\begin{align}
		I_{\M^{p,q}_m} 
		= \frac{1}{\langle \gamma, g \rangle } V^\prime_\gamma V_g
		\end{align}
		schreiben.
		Insbesondere ist $ V^\prime_\gamma $ surjektiv.
	\end{enumerate}
\end{sz}

Schränken wir nun den Definitionsbereich auf $ \{ F \in \L^{p,q}_m \ | \ \exists f \in  \M_m^{p,q} : F = V_g f\}  $ ein ist 
$ V_\gamma^\prime  $ eine bijektive Abbildung mit der STFT als Inverse.

\begin{proof}
	\begin{enumerate}[label =\textbf{(\roman*)}]
		\item Zunächst zeigen wir, dass $ V_\gamma^\prime F  $ in $ \S^\prime(\R^d) $ liegt.
		Sei $ \varphi \in \S(\R^d) $ beliebig. Dann erhalten wir mit der Hölder-Ungleichung \eqref{eq:hoelder_mixed_norm} und $ \S(\R^{2d}) \subseteq \L^{p^\prime,q^\prime}_{\nicefrac{1}{m}} $
		\begin{align*}
		| \langle V^\prime_\gamma F, \varphi \rangle |
		&=
		| \langle F, V_\gamma \varphi \rangle |\\
		&\leq 
		\| F \|_{\L^{p,q}_m} \| V_\gamma \varphi \|_{\L^{p^\prime, q^\prime}_{\nicefrac{1}{m}}}\\
		&\leq 
		\| F \|_{\L^{p,q}_m} \| (1+ | \cdot |)^n V_\gamma \varphi \|_\infty
		\| (1+|\cdot|)^{-n} \|_{\L^{p^\prime,q^\prime}_{\nicefrac{1}{m}}} < \infty
		\end{align*}
		für hinreichend große $ n $. Mit der Äquivalenz der Seminormen \ref{th:equivalence_seminorms} ist der Ausdruck
		\begin{align*}
		V_\gamma^\prime F
		=
		\int \limits_{\R^{2d}}
		F(x,\omega) M_\omega T_x \gamma \td{(x,\omega)}
		\end{align*}
		eine wohldefinierte temperierte Distribution. Nach \ref{th:STFT_continous} ist die STFT von $ V_\gamma^\prime  F$ stetig. Zunächst betrachten wir
		\begin{align*}
		V_\gamma (M_\eta T_u g)(x,\omega)
		&=
		\langle M_\eta T_u g , M_\omega T_x \gamma \rangle 
		= 
		\langle M_{\eta-\omega} T_u g ,T_x \gamma \rangle
		=
		\langle T_{-x} M_{\eta-\omega} T_u g , \gamma \rangle\\
		&=
		\langle e^{2\pi \i x(\eta - \omega)} M_{\eta - \omega} T_{u-x} g, \gamma \rangle
		=
		e^{2 \pi \i x \cdot (\eta  - \omega)} 
		\overline{V_g \gamma (u-x , \eta - \omega)}
		\end{align*}
		und erhalten damit
		\begin{align*}
		V_g V_\gamma^\prime F(u,\eta)
		&=
		\langle V_\gamma^\prime F, M_\eta T_u g \rangle
		=
		\langle  F,V_\gamma (M_\eta T_u g) \rangle\\
		&=
		\int \limits_{\R^{2d}} F(x,\omega)
		\overline{V_\gamma(M_\eta T_u g)(x,\omega)} \td{(x,\omega)}\\
		&=
		\int \limits_{\R^{2d}} F(x,\omega) 
		V_g \gamma (u - x, \eta - \omega) e^{2\pi \i x(\eta - \omega)} \td{(x,\omega)}.
		\end{align*}
		Durch punktweises Abschätzen folgt dann
		\begin{align}\label{eq:invert_on_mod_space_1}
		|V_g V^\prime_\gamma F(u , \eta)|
		\leq 
		(| F | \ast |V_g \gamma |)(u,\eta)
		\end{align}
		und mit \ref{th:convolution_mixed_norm} gilt 
		\begin{align}\label{eq:invert_on_mod_space_2}
		\| V_g( V^\prime_\gamma F) \|_{\L^{p,q}_m}
		\leq 
		C \| F\|_{\L^{p,q}_m} \| V_g \gamma\|_{\L^1_v} < \infty.
		\end{align}
		Wegen $ g, \gamma \in \S(\R^d) $ ist $ V_g \gamma $ nach \ref{th:equivalence_STFT_bounded} schnell fallend und somit gilt $ V_g \gamma \in \L^1_v $.  
		Mit dem Standardfenster $ g_0  $ erhalten wir dann
		\begin{align*}
		\| V_\gamma^\prime F\|_{\M^{p,q}_m}
		\leq 
		\|V_{g_0}(V_\gamma^\prime F) \|_{\L_m^{p,q}}
		\leq
		C \| F \|_{\L_m^{p,q}} \| V_{g_0} \gamma \|_{\L^1_v}.
		\end{align*}
		Also gilt $ V_\gamma^\prime F \in \M_m^{p,q} $.
		
		\item
		Für $ V_g f \in \L_m^{p,q} $ gilt 
		\begin{align*}
		\tilde{f} = \frac{1}{\langle \gamma, g \rangle} V_\gamma^\prime V_g f  \in \M_m^{p,q}
		\end{align*}
		und mit der Inversionsformel \eqref{eq:inverse_STFT_distribution} folgt $ \tilde{f} = f $.
	\end{enumerate}
\end{proof}

\begin{sz}\label{th:independence_of_window_norm}
	Die Definition von $ \M_m^{p,q} $ ist  unabhängig von der Wahl des Fensters $ g \in \S(\R^d)  \setminus \{ 0\}$. Verschiedene Fensterfunktionen liefern äquivalente Normen.
\end{sz}

\begin{proof}
	Wir verwenden \eqref{eq:invert_on_mod_space_2} mit $ g = \gamma,  $ und $ \| g \|_2 = 1 $. Damit erhalten wir
	\begin{align*}
	\| f \|_{\M_m^{p,q}}
	=
	\| V_{g_0} f \|_{\L_m^{p,q}}
	=
	\| V_{g_0} (V_g^\prime V_g f) \|_{\L_m^{p,q}}
	\leq 
	C 
	\| V_{g_0} g \|_{\L^1_v}
	\| V_g f \|_{\L_m^{p,q}}
	=
	C^\prime \| V_g f \|_{\L_m^{p,q} }.
	\end{align*}
	Durch Vertauschen von $ g_0 $ und $ g $ folgt
	\begin{align*}
	\| V_g f \|_{\L_m^{p,q}}
	\leq 
	C \| V_g g_0 \|_{\L^1_v}
	\| V_{g_0 } f \|_{\L_m^{p,q}} =
	C^\prime \| f \|_{\M_m^{p,q}}. 
	\end{align*}
	Insgesamt erhalten wir 
	\begin{align*}
	c \| V_{g_0} f \|_{\L_m^{p,q}}
	\leq
	\| V_g f \|_{\L_m^{p,q}}
	\leq
	C \| V_{g_0} f \|_{\L_m^{p,q}}
	\end{align*}
	für ein $ c,C > 0 $. Es gilt $ f \in \M^{p,q}_m $ genau dann, wenn $ V_g f \in \L_m^{p,q} $ für beliebige $ g \in \S(\R^d) \setminus \{0\} $.
	Also gilt die Normäquivalenz.
\end{proof}

\begin{lem}\label{th:estimate_convolution_pointswise_modulation}
	Seien $ g_0, g , \gamma \in \S(\R^d) \setminus \{ 0\} $ mit $ \langle \gamma , g \rangle \neq 0 $ und $ f \in \S^\prime(\R^d) $. Dann gilt
	\begin{align}\label{eq:estimate_convolution_pointswise_modulation}
	| V_{g_0 } f(x,\omega)|
	\leq
	\frac{1}{|\langle \gamma , g \rangle |}
	(|V_g f| \ast |V_{g_0} \gamma |)(x,\omega)
	\end{align}
	für alle $ (x,\omega) \in \R^{2d} $.
\end{lem}

\begin{proof}
	Mit der Inversionsformel
	\begin{align*}
	f = \frac{1}{\langle \gamma , g  \rangle}
	\int \limits_{\R^{2d}} V_g f(x,\omega ) M_\omega T_x  \gamma \td{(\omega,x)}
	\end{align*}
	erhalten wir $V_{g_0 } f = \langle \gamma , g \rangle^{-1}V_{g_0} V^\prime_\gamma(V_g f) $.
	Mit \eqref{eq:invert_on_mod_space_1} folgt dann direkt die Aussage.
\end{proof}

\begin{sz}\label{th:density_schwartzspace}
	Sei $ m(z) \leq C (1+ |z|)^N $ und $ 1 \leq p,q < \infty $.
	Dann ist der Schwartzraum $ \S(\R^d)  $ ein dichter Teilraum von $ \M_m^{p,q} $.
\end{sz}

\begin{proof}
	Sei $ f \in \S(\R^d) $ beliebig.	
	% und o.B.d.A $ | \langle \gamma, g \rangle |= 1 $.
	Dann gilt wegen $ V_{g_0} f \in \S(\R^d) $ 
	\begin{align*}
	\| f \|_{\M_m^{p,q}}
	&=
	\| V_{g_0} f \|_{\L_m^{p,q}}
	\leq
	%\| V_g f \ast V_{g_0} \gamma \|_{\L^{p,q}_m}
	\int \limits_{\R^{2d}} | V_{g_0} f(z) | (1+ |z|)^N  \td{z}\\
	&\leq 
	\| V_{g_0} f v_{N+2d+1} \|_\infty 
	\int \limits_{\R^{2d}}
	\frac{1}{(1+|z|)^{2d+1}} \td{z}
	< \infty.
	\end{align*}
	Also gilt $ \S(\R^d) \subseteq \M^{p,q}_m $. Es bleibt die Dichtheit zu zeigen.
	Sei nun $ f \in \M_m^{p,q} $ beliebig, $ K_n := \{z \in \R^{2d} \ | \ |z| \leq n \} $ und $ g\in \S(\R^d) \setminus \{0\} $ mit $ \| g \|_2 = 1 $.
	Wir definieren $ F_n = V_g f \cdot \chi_{K_n} $.
	Wegen 
	\begin{align*}
	F_n(z) (1 + |z|)^k 
	\leq
	\sup \limits_{z \in K_n} V_g f(z) (1+|z|)^k < \infty 
	\end{align*}
	für alle $ k \in \N_0 $, sind die Voraussetzungen von \ref{th:rapid_decay_implies_schwartz} erfüllt und
	\begin{align*}
	f_n 
	=
	V_g^\prime F_n 
	=
	\int \limits_{K_n} V_g f (x,\omega) M_\omega T_x g  \td{(x,\omega)}
	\end{align*}
	ist eine Schwartzfunktion. Mit der Inversionsformel \eqref{eq:inversion_formula_modulation} erhalten wir
	\begin{align*}
	\| f - f_n \|_{\M_m^{p,q}} 
	=
	\| V_g^\prime (V_g f - F_n) \|_{\M_m^{p,q}}
	\leq C \| V_g g \|_{\L^1_v} \| V_g f - F_n \|_{\L^{p,q}_m}.
	\end{align*}
	Damit gilt $ \| f - f_n \|_{\M_m^{p,q}} \to 0  $ für $ p,q < \infty  $.
	Also ist $ \S(\R^d) $ dicht in $ \M_m^{p,q} $.
\end{proof}

\begin{sz}\label{th:mod_space_banach}
	$ \M^{p,q}_m(\R^d) $ ist ein Banachraum für $ 1 \leq p,q \leq \infty $.
\end{sz}

\begin{proof}
	Wir betrachten $ V = \{ F \in \L_m^{p,q} \ | \ \exists f \in \S^\prime(\R^d) : F = V_{g_0} f \} $. Dann ist $ V $ ein Unterraum von $ \L_m^{p,q} $ und ist isometrisch isomorph zu $ \M_m^{p,q} $. Damit erhält $ \M_m^{p,q} $ seine Vektorraumstruktur von $ \L_m^{p,q} $.
	Es bleibt zu zeigen, dass $ V $ ein abgeschlossener Unterraum von $ \L^{p,q}_m $ ist. Dann ist $ \M_m^{p,q} $ vollständig.\\
	Sei $ (f_n)_{n \in \N} $ ein Cauchyfolge in $ \M_m^{p,q} $. Nach Definition ist dann $ (V_{g_0}  f_n)_{n \in \N} $ eine Cauchyfolge in $ \L_m^{p,q} $. Damit existiert ein $ F \in \L_m^{p,q}$, sodass
	\begin{align*}
	\lim \limits_{n \to \infty}
	\| V_{g_0} f_n - F \|_{\L_m^{p,q}} = 0
	\end{align*}
	gilt. Nach \ref{th:inversions_formula_modulation} ist
	\begin{align*}
	f = \frac{1}{\| g_0 \|^2_2} V_{g_0}^\prime F 
	= 
	\frac{1}{\| g_0 \|^2_2}
	\int \limits_{\R^{2d}} F(x,\omega) M_\omega T_x g_0 \td{(x, \omega)}
	\end{align*}
	in $ \M_m^{p,q} $ und es gilt 
	\begin{align*}
	f_n = \frac{1}{\| g_0 \|_2^2} V_{g_0}^\prime V_{g_0} f_n.
	\end{align*}
	Mit \eqref{eq:inverse_estimate_modulation} erhalten wir
	\begin{align*}
	\| f - f_n \|_{\M_m^{p,q}}
	=
	\frac{1}{\| g_0\|_2^2}
	\| V_{g_0}^\prime (F - V_{g_0} f_n) \|_{\M_m^{p,q}}
	\leq
	C \| F - V_{g_0} f_n \|_{\L_m^{p,q}}  
	\end{align*}
	und somit auch die Vollständigkeit von $ \M_{m}^{p,q} $.
\end{proof}


\begin{lem}\label{th:mod_space_tf_invariance}
	$ \M^{p,q}_m $ ist invariant unter Zeit-Frequenz-Verschiebungen und es gilt
	\begin{equation}\label{eq:mod_space_tf_invariance}
	\| T_x M_\omega f \|_{\M^{p,q}_m} \leq C v(x,\omega )  \|  f \|_{\M^{p,q}_m} .
	\end{equation}
\end{lem}

\begin{proof}
	Zunächst gilt
	\begin{align*}
	| V_g (M_\omega T_x f)(u , \eta) |
	=
	| \langle T_x M_\omega f , M_\eta T_u g \rangle |
	=
	| \langle f , M_{\eta - \omega} T_{u - x } g \rangle  |
	=
	| T_{(x,\omega) }V_g f(u,\eta) |.
 	\end{align*}
 	Damit erhalten wir durch \eqref{eq:invariance_translation_mixed_norm}
 	\begin{align*}
 	\| T_x M_\omega f\|_{\M_m^{p,q}}
 	=
 	\| T_{(x,\omega) } V_{g_0} f \|_{\L_m^{p,q}} 
 	\leq 
 	C v(x,\omega)  \| V_{g_0} f \|_{\L_m^{p,q}}
 	=
 	C v(x,\omega)  \|  f \|_{\M_m^{p,q}}.
 	\end{align*}
\end{proof}

\begin{lem}\label{th:mod_space_FT_invariance}
	Sei $ p = q $ und $ m(\omega,-x) \leq C m(x,\omega) $.
	Dann ist $ \M_m^p $ invariant unter der Fouriertransformation.
\end{lem}

\begin{proof}
	Mit \eqref{eq:identity_STFT_5}, der Normäquivalenz aus \ref{th:independence_of_window_norm} und $ m(\omega,-x) \leq C m(x,\omega) $ erhalten wir
	\begin{align*}
	\| \hat{f} \|_{\M_m^p}^p
	&=
	\| V_{g_0} \hat{f} \|_{\L^p_m}^p
	\leq 
	C_1
	\| V_{\hat{g_0}} \hat{f} \|_{\L_m^p}^p\\
	&=
	C_1
	\int \limits_{\R^{2d}} |V_{\hat{g_0}} \hat{f}(x,\omega) |^p  m(x,\omega)^p \td{(x,\omega)}\\
	&=
	C_1
	\int \limits_{\R^{2d}} |V_{g_0} f(-\omega,x) |^p  m(x,\omega)^p 
 	\td{(x,\omega)} \\
 	&=
 	C_1
 	\int \limits_{\R^{2d}} |V_{g_0} f(x,\omega) |^p  m(\omega,-x)^p
 	\td{(x,\omega)}\\
 	&\leq
 	\tilde{C} \| f \|_{\M_m^p}^p. 	\qedhere
 	\end{align*}
 	%mit $ m(\omega,-x) \leq C m(x,\omega) $.
\end{proof}


\begin{sz}
	Sei $ 1 \leq p,q < \infty $.
	Dann gilt $ (\M_m^{p,q})^\prime = \M^{p^\prime,q^\prime}_{\nicefrac{1}{m}} $ mit dem Dualitätsprodukt
	\begin{align}
		\langle f, h \rangle
		=
		\int \limits_{\R^{2d}} V_{g_0} f(z) \overline{V_{g_0} h(z)} \td{z}
	\end{align}
	für $ f \in \M^{p,q}_m $ und $ h \in \M^{p^\prime,h^\prime}_{\nicefrac{1}{m}} $.
\end{sz}

\begin{proof}
	Sei $ h \in \M^{p^\prime,q^\prime}_{\nicefrac{1}{m}} $. Wir definieren
	\begin{align*}
	\ell_h(f) :=
	\int
	\limits_{\R^{2d}} V_{g_0} f(z) \overline{V_{g_0} h(z)} \td{z}.
	\end{align*}
	Dies ist wegen der Hölderungleichung \eqref{eq:hoelder_mixed_norm} mit
	\begin{align*}
	| \ell_h(f) |
	\leq
	\| V_{g_0} f \|_{\L_m^{p,q}}
	\| V_{g_0} h \|_{\L_{\nicefrac{1}{m}}^{p^\prime,q^\prime}}
	=
	\|f \|_{\M_m^{p,q}} \| h \|_{\M_{\nicefrac{1}{m}}^{p^\prime,q^\prime}}
	\end{align*}
	ein beschränktes lineares Funktional.
	Sei $ \ell \in (\M_m^{p,q})^\prime $.
	Der Modulationsraum $ \M_m^{p,q} $ ist isometrisch isomorph zu dem abgeschlossenen Unterraum $ V = \{ F \in \L_m^{p,q} \ | \ \exists f \in \S(\R^d) :  F = V_{g_0} f\} $ von $ \L_m^{p,q} $. 
	Damit induziert $ \ell $ ein lineares Funktional $ \tilde{\ell} $ auf $ V $
	mit $ \ell(f) = \tilde{\ell}(V_{g_0} f) $.
	Mit dem Satz von Hahn-Banach  lässt $ \tilde{\ell}  $ sich zu einem linearen Funktional auf $ \L_m^{p,q} $ fortsetzen.
	Nach \ref{th:dualtiy_mixed_norm} existiert ein $ H \in \L^{p^\prime,q^\prime}_{\nicefrac{1}{m}} $, sodass
	\begin{align*}
	\tilde{\ell}(V_{g_0} f)
	=
	\int \limits_{\R^{2d}}
	V_{g_0} f(z) \overline{H(z)} \td{z}
	\end{align*}
	gilt. 
	Wir setzen 
	\begin{align*}
	h = V_{g_0}^\prime H= 
	\int 
	\limits_{\R^{2d}} H(x,\omega) M_\omega T_x g_0 \td{(x, \omega)}
	\end{align*}
	und erhalten $ h \in \M^{p^\prime,q^\prime}_{\nicefrac{1}{m}} $ mit \ref{th:inversions_formula_modulation}.
	Insgesamt folgt dann
	\begin{align*}
	\langle f, h \rangle
	=
	\int 
	\limits_{\R^{2d}}
	V_{g_0} f(z) \overline{V_{g_0 } h(z)} \td{z}
	=
	\int 
	\limits_{\R^{2d}}
	V_{g_0} f(z) \overline{H(z)} \td{z}
	=
	\ell(f)
	\end{align*}
	und es gilt $ (\M_m^{p,q} )^\prime = \M^{p^\prime,q^\prime}_{\nicefrac{1}{m}}  $.
\end{proof}
\newpage
\begin{sz}\label{th:mod_space_window_extention}
	Sei $ m $ eine $ v $-moderierte Gewichtsfunktion und $ g, \gamma \in \M_v^1 $.
	\begin{enumerate}[label =\textbf{(\roman*)}]
		\item 
		$ V^\prime_\gamma : \L_m^{p,q} \to \M_m^{p,q} $ ist beschränkt und die Abschätzung \eqref{eq:inverse_estimate_modulation} ist erfüllt.
		
		\item 
		Die Inversionsformel \eqref{eq:inversion_formula_modulation} gilt für $ f \in \M_m^{p,q} $.
		
		\item
		$ \| V_g  f\|_{\L_m^{p,q}} $ ist eine äquivalente Norm auf $ \M_m^{p,q} $.
	\end{enumerate}

\end{sz}


\begin{proof}
	\begin{enumerate}[label =\textbf{(\roman*)}]
		\item 
		Sei $ F \in \L_m^{p,q} $.
		Wir betrachten die Abbildung 
		\begin{align*}
		T : \S(\R^d) \to \M_m^{p,q}, \ \gamma \mapsto V_\gamma^\prime F.
		\end{align*}
		Nach \eqref{eq:inverse_estimate_modulation} gilt
		\begin{align*}
		\| T \gamma \|_{\M_m^{p,q}} =
		\| V_\gamma^\prime F \|_{\M_m^{p,q}}
		\leq 
		C \| F\|_{\L_m^{p,q}} \| V_{g_0} \gamma \|_{\L_v^1}
		=
		C \| F\|_{\L_m^{p,q}} \|  \gamma \|_{\M_v^1}.
		\end{align*}
		Wegen \eqref{th:density_schwartzspace} ist $ \S(\R^d) $ dicht in $ \M^1_v $.
		Mit \eqref{th:densitiy_priniciple} wird $ T $ zu einem beschränkten Operator von $ \M^1_v $ nach $ \M_{m}^{p,q} $ fortgesetzt.
		
		\item
		Sei zunächst $ g \in \S(\R^d) $.
		Mit $ \gamma = g_0 $ erhalten wir nach \eqref{eq:estimate_convolution_pointswise_modulation} 
		\begin{align*}
		| V_g f (x,\omega)|
		\leq
		\frac{1}{\|g_0\|_2^2}
		(|V_{g_0} f  |\ast | V_g g_0 | )(x,\omega)
		\end{align*}
		für $ (x,\omega) \in \R^{2d} $. Des weiteren gilt
		\begin{align*}
		|V_{g}  g_0(x,\omega)|
		=
		| \langle g_0 , M_\omega T_x g \rangle | 
		=
		| \langle T_{-x} M_{-\omega} g_0, g \rangle |
		=
		| \langle g , M_{- \omega} T_{-x} g_0 \rangle |
		=
		| V_{g_0} g(-x, -\omega) |
		\end{align*}
		und mit der Symmetrie von $ v $ folgt 
		$ \| V_{g_0} g \|_{\L_v^1} = \| V_g g_0\|_{\L^1_v} $.
		Insgesamt erhalten wir
		\begin{align}\label{eq:proof_extension_window_to_M1v}
		\| V_g f \|_{\L_m^{p,q}}
		\leq 
		\frac{C}{\| g_0 \|_2^2}
		\| V_{g_0 } f \|_{\L_m^{p,q}}
		\| V_g g_0 \|_{\L_v^1}
		=
		\frac{C}{\| g_0 \|_2^2}
		\| f \|_{\M_m^{p,q}}
		\| g \|_{\M_v^1}.
		\end{align}
		Der Schwartzraum $ \S(R^d) $ ist dicht in $ \M_v^1 $. Sei $ f \in \M_{m}^{p,q} $ fixiert.
		Wegen \eqref{eq:proof_extension_window_to_M1v} ist 
		\begin{align*}
		T : \M^1_v \to \L_m^{p,q},\
		 g \mapsto V_g f
		\end{align*} 
		mit \ref{th:densitiy_priniciple} ein beschränkter linearer Operator.
		Sei nun $ g, \gamma \in \M_v^1 $ und $ g_n, \gamma_n \in \S(\R^d) $, sodass
		$ \| g_n - g \|_{\M_v^1} \to 0  $ und $ \| \gamma_n - \gamma \|_{\M_v^1} \to 0 $ gilt. Mit \eqref{eq:proof_extension_window_to_M1v} erhalten wir
		$\| V_{g_n } f - V_g f \|_{\L_m^{p,q}} \to 0$
		und durch \textbf{(i)} $ \| V^\prime_{\gamma_n} F - V_{\gamma}^\prime F \|_{\M_m^{p,q}} \to 0$.
		Für $ \tilde{f}\in \M_v^1 $ gilt mit der Plancherelgleichung der STFT \ref{th:orthogonality_relation_kor_1}:
		\begin{align*}
		\| \tilde{f} \|_2^2
		&=
		\frac{1}{\| g_0 \|_2^2}
		\int \limits_{\R^{2d}}
		| V_{g_0} \tilde{f}(x,\omega)|^2 \td{(x,\omega)} \\
		&\leq
		\frac{1}{\| g_0 \|_2^2}
		\| V_{g_0} \tilde{f} \|_\infty
		\int \limits_{\R^{2d}}
		| V_{g_0} \tilde{f}(x,\omega)| v(x,\omega) \td{(x,\omega)},
		\end{align*}
		Hier ist zu beachten,dass $ v $ von der Einsfunktion moderiert wird.
		%Wir wissen jedoch noch
		Mit der Hölder-Ungleichung erhalten wir:
		\begin{align*}
		\| V_{g_0}  f\|_\infty
		\leq
		\sup \limits_{(x,\omega) \in \R^{2d}}
		\int \limits_{\R^d}
		| f(t) | \ | g(t - x )| \td{t}
		\leq
		\|f \|_2 \|g_0 \|_2.
		\end{align*}
		Insgesamt folgt dann
		\begin{align*}
		\|f \|^2_2 
		\leq \frac{1}{\| g_0 \|_2} \| f \|_2 \| f \|_{\M_v^1}
		\end{align*}
		und damit ist $ \M_v^1  $ in $ \L^2 $ eingebettet. Somit gilt
		$ \langle \gamma_n,g_n \rangle \to \langle \gamma , g \rangle  $ für $ n \to \infty $.
		Also erhalten wir mit \eqref{eq:inversion_formula_modulation}
		\begin{align*}
		f=
		\lim \limits_{n \to \infty}
		\frac{1}{\langle \gamma_n,g_n \rangle} V_{\gamma_n}^\prime V_{g_n} f = 
		\frac{1}{\langle \gamma , g \rangle} V_\gamma^\prime V_g f.
		\end{align*}
		\item 
		Die Schritte sind analog zu \ref{th:independence_of_window_norm}.
		
	\end{enumerate}
\end{proof}
\newpage\blankpage
\section{Eigenschaften und Minimalität der Feichtingeralgebra}

Unser Ziel in diesem Abschnitt ist es, die Eigenschaften der Feichtingeralgebra zu diskutieren.
Zuerst beschäftigen wir uns mit den algebraischen Eigenschaften von $ \M^1_v $.
Hieraus folgern wir, dass $ S_0 := \M^1 $ eine Banachalgebra  bezüglich Faltung und punktweiser Multiplikation ist.
Danach zeigen wir die Minimalität der Feichtingeralgebra und zum Abschluss einige Folgerungen daraus.


\begin{genericthm}{Algebraische Eigenschaften}\label{th:algebraic_properties}
	\begin{enumerate}[label =\textbf{(\roman*)}]
		\item Sei $ v(x,\omega) = u_1(x)$. Dann gilt $ \L^1_u \ast \M_u^1 \subseteq \M^1_u $ mit
		\begin{align*}
		\| k \ast f \|_{\M^1_{u_1}} \leq \| k \|_{\L^1_{u_1}} \| f \|_{\M^1_{u_1}}.
		\end{align*}
		%Außerdem ist $ \M^1_{u_1} $ eine Banachalgebra bezüglich der Faltung.
		\item 
		Sei $ v(x,\omega) = u_2(\omega) $. Dann gilt die punktweise Multiplikation $ \mathcal{F} \L^1_{u_2} \cdot \M_{u_2}^1 \subseteq M_{u_2}^1 $.
	\end{enumerate}
\end{genericthm}

\begin{proof}
	\begin{enumerate}[label =\textbf{(\roman*)}]
		\item Nach \eqref{eq:identity_STFT_7} können wir die STFT durch
		$ V_g f(x,\omega) = e^{-2\pi \i x \cdot \omega} (f \ast M_\omega g^\ast)(x) $ beschreiben.
		Damit gilt
		\begin{align*}
		\| f \|_{\M^1_{u_1}} = 
		\int \limits_{\R^{2d}} |V_g f(x,\omega) | u_1(x) \td{(x,\omega)}
		=
		\int \limits_{\R^d} \| f \ast M_\omega^\ast\|_{\L^1_{u_1}} \td{\omega}
		\end{align*}
		für $ f \in \M^1_{u_1} $. Mit \ref{th:convolution_mixed_norm} erhalten wir durch
		\begin{align*}
		\| k \ast f \|_{\M^1_{u_1}} 
		=
		\int \limits_{\R^d}
		\| (k \ast f) \ast M_\omega g^\ast \|_{\L^1_{u_1}} \td{\omega}
		\leq
		\int \limits_{\R^d}\| k \|_{\L^1_{u_1}}
		\| f \ast M_\omega g^\ast \|_{\L^1_{u_1}} \td{\omega}
		\end{align*}
		für $ k \in \L^1_{u_1} $ die passende Abschätzung.
		
		\item
		%Folgt aus der ersten Aussage.
		Wegen \eqref{eq:identity_STFT_5} gilt
		\begin{align*}
		\| f \|_{\M_{u_2}^1}
		&=
		\| f \|_{\M_v^1}
		=
		\int \limits_{\R^{2d}}
		| V_g f (x,\omega) | v(x,\omega ) \td{(x,\omega)}
		=
		\int \limits_{\R^{2d}}
		| V_g \hat{ f} (- \omega,x) | v(x,\omega ) \td{(x,\omega)}\\
		&=
		\int \limits_{\R^{2d}}
		| V_g \hat{f } (x,\omega) | v(- \omega,x ) \td{(x,\omega)}
		=
		\| \hat{f } \|_{\M^1_{\tilde{ v}}}
		\end{align*}
		mit $ \tilde{ v}(x,\omega) := v(-\omega,x) = u_2(x) $.
		Also erhalten wir $ f \in \M^1_v $ genau dann, wenn $ \hat{ f} \in \M_{\tilde{ v}}^1 $ erfüllt ist.
		Mit der ersten Aussage folgt dann 
		\begin{align*}
		\| kf \|_{\M^1_{u_2}} = 
		\| \hat{k} \ast \hat{f} \|_{\M^1_{\tilde{v}}}
		\leq 
		\| \hat{k} \|_{\L^1_{\tilde{v}}}
		\| \hat{f} \|_{\M^1_{\tilde{v}}}
		\leq 
		C \| k \|_{\mathcal{F}\L^1_{u_2}} \| f \|_{\M^1_{u_2}}
		\end{align*}
		für $ k \in \mathcal{F} \L^1_{u_2} $ und $ f \in \M^1_{u_2} $.
	\end{enumerate}
\end{proof}

\begin{df}
	Sei $ f $ eine messbare Funktion und $ u $ ein submultiplikatives Gewicht auf $ \R^d $. 
	Der \textit{Amalgamraum} $ \W(\L^1(\R^d)) $ ist über die Norm
	\begin{equation}
	\| f \|_{\W(\L^1_u)}
	:=
	\sum \limits_{k \in \Z^d}
	\esssupp \limits_{x \in [0,1]^d} |f(x+k)| u(k)
	\end{equation}
	definiert.
\end{df}

\begin{sz}\label{th:embedding_amalagm_L1}
	Sei $ u $ ein submultiplikatives Gewicht.
	Dann ist die Einbettung $ \W(\L^1_u(\R^d)) \hookrightarrow \L^1_u(\R^d) $ stetig.
\end{sz}

\begin{proof}
	Sei $ f \in  \W(\L^1_u(\R^d))$.
	Dann gilt mit der Submultiplikativität
	\begin{align*}
	\| f \|_{\L^1_u}
	&= 
	\sum \limits_{k \in \Z^d}  \int \limits_{[0,1]^d} | f(x+k) | u(x+ k) \td{x}
	\leq 
	\sum \limits_{k \in \Z^d} \esssupp_{x \in [0,1]^d} |f(x+k)| u(k) 
	\int \limits_{[0,1]^d} u(x) \td{x}\\
	&= C \| f \|_{\W(\L^1_u)}.
	\end{align*}
\end{proof}

\begin{df}
	Sei $ v $ ein submultiplikatives Gewicht auf $ \R^{2d} $.
	Wir führen die submultiplikativen Gewichte
	\begin{align*}
	u_1(x) = \inf \limits_{\omega \in \R^d} v(x,\omega),
	\quad 
	u_2(\omega) = \inf \limits_{x \in \R^d} v(x,\omega)
	\end{align*}
	auf $ \R^d $ ein.
\end{df}
%\textcolor{red}{\textbf{Beweisen?}}

\begin{lem}\label{th:am_space}
	Sei $ f \in \M^1_v $.
	Dann gelten $ f \in \W(\L^1_{u_1}) $ und $ \hat{f} \in W(\L^1_{u_2}) $.
\end{lem}

\begin{proof}
	Zuerst zeigen wir $ f \in \W(\L^1_{u_1}) $.
	Sei $ g \in C^\infty_0(\R^d) $ mit $ 0 \leq g(x) \leq 1 $ auf $ \R^d $ und $ g(x) = 1  $ auf $ [-1,1]^d $. Dann gilt 
	\begin{align*}
	\chi_Q(t) \leq T_x g(t)
	\end{align*}
	für $ t \in \R^d $ und $ x \in Q = [0,1]^d $. 
	Hieraus erhalten wir
	\begin{align*}
	\| f \cdot T_k \chi_Q \|_\infty
	=
	\esssupp_{t \in \R^d} | f(t) \chi_Q(t - k) |
	\leq 
	\esssupp_{t \in \R^d} | f(t) T_{k+x} \overline{g}(t) |
	=
	\| f \cdot T_{k+x} \overline{g} \|_\infty.
	\end{align*}
	Wegen $ \| h \|_\infty \leq \| \hat{h} \|_1 $ folgt dann mit \eqref{eq:identity_STFT_1}
	\begin{align*}
	\| f \cdot T_{k+x} \overline{g} \|_\infty
	\leq
	\left\|(f \cdot T_{k+x} \overline{g})^{\hat{}} \right\|_1
	=
	\int \limits_{\R^d}
	| V_g f(k+x, \omega) | \td{\omega}.
	\end{align*}
	Mit \eqref{eq:v_mod_prop_2} gilt $ u_1(k) \leq C u_1(k+x) $ für $ x \in Q $ und mit $ u_1(x) \leq v(x,\omega) $ folgt 
	\begin{align*}
	\| f \|_{\W(\L^1_{u_1})} 
	&= 
	\sum \limits_{k \in \Z^d}
	\| f \cdot T_k\chi_Q \|_\infty u_1(k)
	\leq 
	C
	\sum \limits_{k \in \Z^d} 
	\int \limits_Q 
	\int \limits_{\R^d}
	| V_g f(x+ k, \omega) |  \td{\omega}  \ u_1(k + x) \td{x}\\
	&=
	C\int \limits_{\R^{2d}} |V_g f(x, \omega) | u_1(x) \td{(x,\omega)} 
	\leq C \| V_g f \|_{\L^1_v(\R^{2d})}
	= C \| f \|_{\M^1_v}.
	\end{align*}
	Es bleibt also noch $ \hat{f} \in \W(\L^1_{u_2}) $ zu zeigen.
	Wir wählen $ \hat{g} \in C^\infty_0(\R^d) $, sodass $ 0 \leq \hat{g}(\omega) \leq 1 $
	auf $ \R^d $ und $ g(\omega) = 1 $ auf $ [-1,1]^d $ gilt. 
	Dann erhalten wir mit denselben Argumenten
	\begin{align*}
	\| \hat{f} \|_{\W(\L^1_{u_2})}
	\leq 
	C 
	\int \limits_{\R^{2d}}
	| V_{\hat{g}} \hat{f}(x,\omega) | u_2(x) \td{(x,\omega)}.
	\end{align*}
	Durch \eqref{eq:identity_STFT_5} gilt $ | V_{\hat{g}} \hat{f}(x,\omega) |
	= | V_g f (-\omega,x)|  $ und mit $ u_2(x) \leq v(-\omega,x) $ folgt dann
	\begin{align*}
	\| \hat{f} \|_{\W(\L^1_{u_2})}
	\leq 
	C \int \limits_{\R^{2d}}
	| V_g f(-\omega,x) | v(-\omega,x) \td{(x,\omega)}
	=
	C \int \limits_{\R^{2d}}
	| V_g f(x,\omega) | v(x,\omega) \td{(x,\omega)}
	= \| f \|_{\M^1_v}.
	\end{align*}
	\end{proof}
	
Für $ f \in \M^1_v $ folgt
wegen \ref{th:embedding_amalagm_L1} dann $ f \in \L^1_{u_1} $ und $ \hat{f } \in \L^1_{u_2} $.
Damit gilt auch $ f \in \F \L_{u_2}^1 $ und $ f $ ist stetig, da das Fourierbild einer $ \L^1 $-Funktion eine stetige Funktion liefert.
Wir erhalten $ \M^1_{u_1} \subseteq \L^1_{u_1} $ und $ \M^1_{u_2} \subseteq \L^1_{u_2} $.
Insbesondere ist der ungewichtete Raum $ \M^1 $ nach \ref{th:algebraic_properties} eine Banachalgebra bezüglich Faltung und punktweiser Multiplikation.


\begin{df}
	Wir bezeichnen den Modulationsraum $ \M^1 $ als \textit{Feichtingeralgebra} und
	setzen $S_0 := S_0(\R^d) := \M^1 $, 
	$ S_0^\prime := S_0^\prime(\R^d) := \M^\infty $.
\end{df}




\begin{df}
	Sei $ g \in \M_v^1\setminus \{0\} $ und
	\begin{align}\label{eq:gabor_expansion_modspace}
	f = \sum \limits_{n=1}^\infty c_n T_{x_n} M_{\omega_n} g
	\end{align}
	für eine beliebige Menge $ \lbrace (x_n, \omega_n)\in \R^{2d} \ | \ n \in \N \rbrace $ gegeben. 
	Wir nennen $ f $ eine \textit{nicht-uniforme Gaborentwicklung}, falls
	\begin{align*}
	\sum \limits_{n=1}^\infty |c_n | v(x_n,\omega_n) < \infty 
	\end{align*}
	erfüllt ist. Wir schreiben für den Raum dieser Entwicklungen $ \mathcal{M}_v $.
	Die Norm ist durch
	\begin{align*}
	\| f \|_{\mathcal{M}_v} = \inf \sum \limits_{n = 1}^\infty |c_n| v(x_n,\omega_n)
	\end{align*}
	gegeben. Hierbei wird das Infimum über alle zulässigen Entwicklungen von $ f  $ gebildet.
\end{df}
Wir werden zeigen, dass $ \M_v^1 = \mathcal{M}_v $ gilt. 
$ \mathcal{M}_v $ ist ein Banachraum, was direkt aus der \\Normäquivalenz von $ \| \cdot \|_{\M^1_v} $ und $ \| \cdot \|_{\mathcal{M}_v} $ folgt.

\begin{lem}\label{th:gabor_rep_m_1_v}
	Es gilt $ \M_v^1 = \mathcal{M}_v $.
	Dementsprechend besitzen alle $ f \in \M^1_v $ eine Gaborentwicklung
	\begin{align*}
	f = \sum \limits_{n=1}^\infty c_n T_{x_n} M_{\omega_n} g
	\end{align*}
	und $ \| \cdot \|_{\mathcal{M}_v} $ ist äquivalent zu $ \| \cdot\|_{\M_v^1} $.
\end{lem}

\begin{proof}
	Wir betrachten $ \ell^1_v(\R^{2d}) $ mit der Norm 
	\begin{align*}
	\| c \|_{\ell^1_v} 
	=
	\sum \limits_{(x,\omega) \in \R^{2d}} |c(x,\omega)| v(x,\omega) < \infty.
	\end{align*}
	Die Elemente in $ \ell^1_v(\R^{2d}) $ bilden nach $ \C $ ab und die Norm kann nur endlich sein, falls $ c(x,\omega) \neq 0 $ nur für abzählbare $ (x,\omega) $ gilt. Wir definieren
	\begin{align*}
	\mathcal{D}_g \ : \ \ell^1_v(\R^{2d}) \to \M^1_v , \
	c \mapsto \sum \limits_{(x, \omega) \in \R^{2d} } c(x,\omega) T_x M_\omega g.
  	\end{align*}
  	Dies entspricht aufgrund des Abzählbarkeitsarguments gerade \eqref{eq:gabor_expansion_modspace}. Mit \eqref{eq:mod_space_tf_invariance} erhalten wir durch
  	\begin{align*}
  	\| \mathcal{D}_g c \|_{\M^1_v} 
  	\leq 
  	\sum \limits_{c(x,\omega) \neq 0}
  	|c(x,\omega) | \| T_x M_\omega g \|_{\M^1_v} 
  	\leq
  	\sum \limits_{c(x,\omega) \neq 0}
  	|c(x,\omega) | C v(x,\omega) \|  g \|_{\M^1_v}
  	=
  	C \|  g \|_{\M^1_v} \| c\|_{\ell^1_v} 
  	\end{align*}
  	die Stetigkeit von $ \mathcal{D}_g $.
  	Nach Konstruktion gilt also $ \Bild \ \mathcal{D}_g = \mathcal{M}_v  \subseteq \M^1_v$.
  	Für die Surjektivität bestimmen wir den adjungierten Operator
	$ \mathcal{D}_g^\prime  : \M^\infty_{\nicefrac{1}{v}} \to \ell^1_{\nicefrac{1}{v}}$.
	Es gilt
	\begin{align*}
	\langle \mathcal{D}_g^\prime f , c \rangle 
	=
	\langle f, \mathcal{D}_g c \rangle
	=
	\sum \limits_{c(x,\omega) \neq 0} 
	\langle f, T_x M_\omega g \rangle \overline{c(x,\omega)}.
	\end{align*}
	
	Mit
	\begin{align*}
	\mathcal{D}_g^\prime f(x,\omega) 
	= 
	\langle f, T_x M_\omega g \rangle
	=
	e^{2\pi \i x \cdot \omega} V_g f(x,\omega)
	\end{align*}
	erhalten wir durch
	\begin{align*}
	\| \mathcal{D}_g^\prime f \|_{\ell^\infty_{\nicefrac{1}{v}}}
	= \| V_g f \|_{\L^\infty_{\nicefrac{1}{v}}} 
	=  \| f \|_{\M_{\nicefrac{1}{v}}^\infty}, 
	\end{align*}
	dass $ \mathcal{D}_g^\prime $ eine Isometrie ist.
	Damit ist $ \mathcal{D}_g^\prime $ injektiv und besitzt ein abgeschlossenes Bild.
	Hieraus folgt mit dem Satz vom abgeschlossenen Bild die Surjektivität von $ \mathcal{D}_g $. Also gilt $ \M_v^1 = \mathcal{M}_v $.
	Nun wenden wir uns der Normäquivalenz von $ \| \cdot \|_{\M_v^1} $ und $ \| \cdot
	 \|_{\mathcal{M}_v} $ zu.
	Sei $ N = \Kern \ \mathcal{D}_g $. Dann induziert $ \mathcal{D}_g $ einen Isomorphismus $ \tilde{\mathcal{D}}_g $:
	\begin{center}
		\begin{tikzcd}
		\ell^1_v \arrow{d}[swap]{\pi}  \arrow{r}{\mathcal{D}_g} & \M_v^1  \\
		
		 \ell^1_v / N \arrow{ur}[swap]{\tilde{\mathcal{D}}_g} &  \\
		\end{tikzcd}
	\end{center}
	Für diesen gilt $ \tilde{\mathcal{D}}_g(c + N)  = \mathcal{D}_g(c) $.
	Mit dem Inverse-Mapping-Theorem erhalten wir die Normäquivalenz von 
	$ \| c + N \|_{\ell_v^1 / N} := \inf_{n \in N} \| c + n \|_{\ell^1_v}  $ 
	und $ \|\tilde{\mathcal{D}}_g(c + N)\|_{\M_v^1} = \| \mathcal{D}_g( c ) \|_{\M^1_v}   $.
	Gilt nun $ f = \mathcal{D}_g c  $, erhalten wir nach Konstruktion
	$ \| f \|_{\mathcal{M}_v } = \| c + N \|_{\ell_v^1 / N} $ und wir haben die Normäquivalenz gezeigt.
\end{proof}



\begin{genericthm}{Minimalitätseigenschaft}\label{th:minimality_feichtinger}
	Sei $ B \subset \mathcal{S}^\prime(\R^d) $ ein Banachraum mit:
	\begin{enumerate}[label =\textbf{(\roman*)}]
		\item 
		$ B $ ist invariant unter Zeit-Frequenz-Verschiebungen und es gilt
		\begin{align*}
		\| T_x M_\omega f\|_B \leq C  v(x, \omega) \| f \|_B
		\end{align*}
		für alle $ f \in B $.
		Hierbei ist $ v $ ein submultiplikatives Gewicht.
		\item
		$ \M_v^1 \cap B \neq \{ 0\} $
	\end{enumerate}
	Dann existiert eine stetige Einbettung von $ \M_v^1 $ in $ B $.
\end{genericthm}

\begin{proof}
	Sei $ g \in \M^1_v \cap B  \setminus \{ 0\}$ .
	Nach \ref{th:gabor_rep_m_1_v} besteht $ \M_v^1 $ aus allen nicht-uniformen Gaborentwicklungen 
	\begin{align*}
	f = \sum \limits_{n = 1}^\infty c_n T_{x_n} M_{\omega_n} g
	\end{align*}
	mit $ \sum_{n = 1}^\infty |c_n| v(x_n,\omega_n) < \infty$. Mit der Invarianz unter Zeit-Frequenz-Verschiebungen erhalten wir
	\begin{align*}
	\| f \|_B \leq
	\sum \limits_{n=1}^\infty |c_n| \| T_{x_n} M_{\omega_n} g \|_B
	\leq 
	C \sum \limits_{n=1}^\infty |c_n| v(x_n, \omega_n) \| g \|_B 
	\end{align*}
	für alle $ f \in \M_v^1 $. Damit gilt $ \M^1_v \subseteq B $.
	Das Infimum über alle zulässigen Entwicklungen und die Normäquivalenz liefern dann mit
	\begin{align*}
	\| f \|_B \leq \|f \|_{\mathcal{M}_v} \|g \|_B \leq C \| f \|_{\M_v^1}
	\end{align*}
	die Stetigkeit der Einbettung $ \M_v^1 \hookrightarrow B$.
\end{proof}




\begin{kor}
	Sei $ m $ von $ v $ moderiert.
	Dann gilt 
	\begin{align}
	\M^1_v 
	\subseteq
	\M_m^{p,q}
	\subseteq
	\M^\infty_{\nicefrac{1}{v}}
	\end{align}
	für $ 1 \leq p,q \leq \infty $.
\end{kor}

\begin{proof}
	Mit \eqref{eq:mod_space_tf_invariance} folgt
	\begin{align*}
	\| T_x M_\omega f \|_{\M_m^{p,q}}
	\leq 
	C v(x,\omega) \| f \|_{\M_m^{p,q}}
	\end{align*}
	und es gilt $ \S(\R^d) \subset \M^1_v \cap \M_m^{p,q} $.
	Damit sind die Voraussetzungen von \ref{th:minimality_feichtinger} erfüllt und wir erhalten $ \M_v^1 \subseteq \M_m^{p,q} $.
	Für $ p,q < \infty $ gilt wegen $ \M_v^1 \subseteq \M_{\nicefrac{1}{m}}^{p^\prime, q^\prime} $ auch
	\begin{align*}
	\M_{m}^{p,q} 
	= 
	\left(
	\M_{\nicefrac{1}{m}}^{p^\prime,q^\prime} \right)^\prime
	\subseteq 
	(\M_v^1)^\prime 
	=
	\M_{\nicefrac{1}{v}}^\infty.
	\end{align*}
	Für $ p  = \infty $ und $ q < \infty $ gilt
	\begin{align*}
	\M^1_v \subseteq \M^{1,q^\prime}_{\nicefrac{1}{m}}
	\
	\Rightarrow
	\
	\M^{\infty,q}_m \subseteq \M^\infty_{\nicefrac{1}{v}}
	,
	\end{align*}
	da $ \nicefrac{1}{m} $ nach \eqref{eq:v_mod_prop_1} von $ v $ moderiert wird. Der umgekehrte Fall folgt analog.
	Für $ p,q = \infty $ erhalten wir mit \eqref{eq:v_mod_prop_3}
	\begin{align*}
	\| f \|_{\M_{\nicefrac{1}{v}}^\infty}
	=
	\| V_g f \|_{\L^\infty_{\nicefrac{1}{v}}}
	=
	\sup \limits_{z \in \R^{2d}}
	\left| V_g f(z) \frac{1}{v(z)} \right|
	\leq 
	\tilde{C} \| f \|_{\M_m^\infty}.
	\end{align*}
\end{proof}

Insbesondere ist mit \ref{th:minimality_feichtinger} auch die Feichtingeralgebra $ S_0(\R^d)  $ minimal.
Zum Abschluss wollen wir noch einige Folgerungen aus der Minimalität von $ S_0(\R^d) $ betrachten.
In \eqref{th:mod_space_tf_invariance} wurde bereits gezeigt, dass $ S_0 $ und $ S_0^\prime  $ invariant unter Zeit-Frequenzverschiebungen sind.
Durch das Fehlen einer Gewichtsfunktion in der Norm von $ S_0 $ folgt die Invarianz sofort und die von $ S_0^\prime $ erhalten wir unmittelbar daraus.


\begin{lem}
	Die Räume $ S_0 $ und $ S_0^\prime  $ sind invariant unter Zeit-Frequenzverschiebungen und es gilt
	\begin{align}
	\| T_x M_\omega f \|_{S_0} = \| f \|_{S_0}.
	\end{align}
\end{lem}

\begin{proof}
	Die Invarianz von $ S_0  $ folgt direkt aus der Normdefinition, \ref{th:properties_TF} und \ref{th:STFT_modulation_TF_argument}.
	Desweiteren gilt damit $ \| T_x M_\omega f \|_{S_0}  = \| f \|_{S_0} $.
	Damit erhalten wir 
	\begin{align*}
	\|T_x M_\omega f \|_{S_0^\prime}
	=
	\sup \limits_{\| \varphi \|_{S_0} = 1} | (T_x M_\omega f)(\varphi) |  
	=
	\sup \limits_{\| \varphi \|_{S_0} = 1} | f(\underbrace{M_{-\omega} T_{-x} \varphi}_{\in S_0}) |
	\leq 
	\sup \limits_{\| \varphi \|_{S_0} = 1} | f( \varphi) |
	= \| f \|_{S_0^\prime}.
	\end{align*}
\end{proof}

\begin{sz}
	Sei $ \alpha : \R^d \to \R^d $ eine lineare und bijektive Abbildung.
	Dann definiert
	\begin{align}
	T_\alpha : S_0(\R^d) \to S_0(\R^d), \ f \mapsto  f \circ \alpha
	\end{align}
	eine lineare Bijektion. 
\end{sz}

\begin{proof}
	Die Linearität dieser Abbildung ist klar.
	Die Injektivität folgt durch
	\begin{align*}
	f \circ \alpha = 
	g \circ \alpha 
	\ \Leftrightarrow \ f = g.
	\end{align*}
	Sei $ f\in S_0 $ beliebig.
	Dann gilt $ T_\alpha (f \circ \alpha^{-1}) = f $.
	Insgesamt ist $ T_\alpha $ bijektiv.
\end{proof}

\begin{df}
	Sei $ B $ eine Banachraum, der dicht in einem Hilbertraum $ H $ eingebettet ist.
	Dann bezeichnen wir $ (B,H,B^\prime) $ als \textit{Gelfandtripel}.
\end{df}

Sei $ i : B \to H  $ die Einbettungsabbildung. Mit $ H = H^\prime $ erhalten wir $ i^\prime : H \to B^\prime $.
Damit können wir die Dualität auf $ B^\prime $ mit dem Skalarprodukt auf $ H $ identifizieren. Das heißt
\begin{align}
\langle h, b \rangle_{B^\prime \times B}
=
\langle b, h \rangle_H
\end{align}
für $ h \in H \subseteq B^\prime $ und $ b \in B \subset H $.

\begin{sz}
	$ S_0 $ liegt dicht in $ \L^2 $ und es gilt $ S_0 \subseteq \L^2 \subseteq S_0^\prime $.
	Insbesondere ist $ (S_0,\L^2, S_0^\prime) $ ein Gelfandtripel.
\end{sz}
 
\begin{proof}
	Wegen $ \| T_x M_\omega f\|_{\L^2} = \| f \|_{\L^2} $ und $ \S(\R^d) \subset \L^2(\R^d) $ sind die Voraussetzungen von \ref{th:minimality_feichtinger} erfüllt.
	Damit wird $ S_0(\R^d)$ stetig in $\L^2(\R^d) $ eingebettet.Weiter ist $ \S(\R^d)  $ dicht in $ S_0(\R^d) $ und in $ L^2(\R^d) $.
	Also ist $ S_0 $ dicht in $ \L^2 $ und es 
	folgt $ S_0 \subseteq \L^2 = (\L^2)^\prime \subseteq S_0^\prime $.
\end{proof}




\begin{df}
	Seien $ (B_1, H_1,B_1^\prime) $ und $ (B_2, H_2,B_2^\prime) $ Gelfandtripel.
	Ein linearer Operator $ T $ heißt unitärer Gelfandtripel-Isomorphismus, falls gilt:
	\begin{enumerate}[label =\textbf{(\roman*)}]
		\item 
		$ T : B_1 \to B_2 $ ist ein Isomorphismus.
		\item
		$ T : H_1 \to H_2 $ ist unitär.
		\item 
		$ T : B_1^\prime \to B_2^\prime $ ist ein schwach*-stetiger Isomorphismus.
		Ebenso ist dieser normstetig.
	\end{enumerate}
\end{df}


\begin{sz}\label{th:fourier_feichtinger}
	Die Fouriertransformation $ \F : S_0(\R^d) \to S_0(\R^d) $ ist eine bijektive Isometrie.
\end{sz}

\begin{proof}
		
		Nach \ref{th:independence_of_window_norm} können wir $ g(x) = e^{-\pi x^2} $ setzen, womit $ g = \hat{g} $ gilt.
		Wegen 
		\begin{align*}
		\| f \|_{S_0} 
		=
		\| V_g f \|_{\L^1} 
		=
		\|V_{\hat{g}} \hat{f} \|_{\L^1}
		=
		\| V_{g} \hat{f} \|_{\L^1}
		=
		\| \hat{f} \|_{S_0}
		\end{align*}
		ist $ \F $ eine Isometrie auf $ S_0 $. Da $ S_0 $ vollständig ist folgt auch die Vollständigkeit von $ \Bild \ \F $.
%		Sei $ \alpha_n $ eine Cauchyfolge in $ \Bild \ \F $. Dann existiert für alle $ \alpha_n $ ein eindeutiges
%		$ f_n $, sodass $ \F f_n = \alpha_n $ gilt.
%		Da $ \F $ eine Isometrie ist, ist $ f_n  $ auch eine Cauchyfolge in $ S_0 $.
%		Damit existiert ein $ f \in S_0 $ mit $ \| f_n - f\|_{S_0} \to 0 $ und es folgt
%		\begin{align*}
%		\| f_n -f \|_{S_0} 
%		=
%		\|\alpha_n - \F f \|_{S_0} \to 0.
%		\end{align*}
%		Also ist das $ \Bild \ \F $ ein Banachraum. 
%		Wegen 
		\begin{align*}
		\| T_x M_\omega \hat{f} \|_{S_0}
		=
		\|  M_x T_{- \omega} f \|_{S_0}
		= 
		\| f \|_{S_0} =  \| \hat{f } \|_{S_0}
		\end{align*}
		ist das Bild von $ \F $ invariant unter Zeit-Frequenzverschiebungen.
		Dann gilt nach \ref{th:minimality_feichtinger} $ S_0 \subseteq \Bild \ \F $.
		Also ist  $ \F $ bijektiv.
\end{proof}

Analog zu der Fouriertransformation auf $ \S^\prime(\R^d)  $ definieren wir die Fouriertransformation auf $ S_0^\prime(\R^d) $.
Sei $ f \in S_0^\prime $ und $ \varphi \in S_0 $ beliebig. Dann definieren wir $ \hat{f} $ durch $  \hat{f} (\varphi) = \langle  \hat{f}, \varphi \rangle := \langle f , \hat{\varphi} \rangle $.


\begin{sz}\label{th:fourier_feicht_dual}
	Die Fouriertransformation
	\begin{align}
	\F : S_0^\prime \to S_0^\prime, \  f \mapsto \hat{f}
	\end{align}
	ist eine bijektive Isometrie.
	Desweiteren ist $ \F $  schwach*-stetig.
\end{sz}

\begin{proof}
	Da die Fouriertransformation auf $ S_0 $ eine bijektive Isometrie ist, erhalten wir
	\begin{align*}
	\|f \|_{S_0^\prime}
	= 
	\sup \limits_{\| \varphi \|_{S_0} = 1} |\hat{f}(\varphi) |
	=
	\sup \limits_{\| \hat{\varphi} \|_{S_0} = 1} |f(\hat{\varphi}) |
	= 
	\| f \|_{S_0^\prime}
	\end{align*}
	und $ \F $ ist surjektiv.
	Sei $ f \in S_0 $ und $ f_n \in S_0^\prime $ eine Folge mit
	\begin{align*}
	\langle f_n ,\varphi \rangle \to \langle f,\varphi \rangle
	\end{align*}
	für alle $ \varphi \in S_0 $.
	Dann folgt mit
	\begin{align*}
	\F f_n(\varphi)
	= \langle f_n , \hat{\varphi} \rangle
	\to 
	\langle f , \hat{\varphi} \rangle
	= \F f (\varphi)
	\end{align*}
	die Schwach*-Stetigkeit.
\end{proof}

Mit den Sätzen \ref{th:fourier_feichtinger} und \ref{th:fourier_feicht_dual} ist die Fouriertransformation $ \F $ ein Gelfandtripel-Automorphismus auf $ (S_0, H, S_0^\prime) $.
Insbesondere ist damit die Parsevalgleichung 
\begin{align}
\langle \hat{f},\hat{g} \rangle
=
\langle f, g \rangle
\end{align}
für  $ (f,g) \in S_0^\prime(\R^d) \times S_0(\R^d)  $ und $ (f,g) \in \L^2(\R^d) \times \L^2(\R^d) $ oder andere sinnvolle Paarungen des Gelfandtripels definiert.






\begin{genericthm}{Kernsatz}
	Sei $ K : S_0(\R^d) \to S_0^\prime(\R^d) $ ein beschränkter Operator.
	Dann existiert ein eindeutiger Kern $ k \in S_0^\prime(\R^{2d}) $, sodass 
	\begin{align}
	\langle K f, g \rangle 
	=
	\langle k , f \otimes g \rangle
	\end{align}
	mit $ f,g \in S_0(\R^d) $ gilt. Das Funktional $ Kf \in S_0^\prime(\R^d) $ lässt sich also durch
	\begin{align}
	Kf(g)
	= 
	\int \limits_{\R^{2d}} k(x,\omega)  g(x) f(\omega) \td{(x,\omega)}
	\end{align}
	beschreiben.
\end{genericthm}

\begin{proof}
	\cite[Abschnitt~14.4]{noauthor2009Foundationsof}
\end{proof}
%\newpage
%\section{Zusammenfassung}
In diesem Abschnitt wollen wir noch einmal die Resultate für $ S_0(\R^d) $ zusammenfassen.

Nach \ref{th:density_schwartzspace} liegt $ \S(\R^d) $ dicht in $ S_0(\R^d) $ und \ref{th:mod_space_banach} liefert uns die Banachraumeigenschaft.
Die Lemmas \ref{th:mod_space_tf_invariance} und \ref{th:mod_space_FT_invariance} liefern die Zeit-Frequenz-Invarianz und die Invarianz unter der Fouriertransformation.
Mit den Sätzen \ref{th:independence_of_window_norm} und \ref{th:mod_space_window_extention} kann ein beliebiges Fenster aus $ S_0 \setminus \{ 0\} $ für die Norm von $ S_0 $ gewählt werden.
Mit \ref{th:algebraic_properties} und \ref{th:am_space} wird $ S_0 $ zu einer Banachalgebra.
Die Minimaltät von $ S_0 $ erhalten wir aus \ref{th:minimality_feichtinger}.
Insbesondere folgen hieraus die restlichen Eigenschaften von $ S_0 $.
 





%\newpage
%\appendix
%\section{Anhang}

%\begin{lem}
%	Sei $ (X, \| \cdot \|) $ ein normierter Raum. Dann gilt sind äquivalent:
%	\begin{enumerate}
%		\item $ X $ ist vollständig.
%		\item Jede absolut konvergente Reihe konvergiert in $ (X, \| \cdot \|) $
%	\end{enumerate}
%\end{lem}

%\begin{genericthm}{Dichtheitsprinzip}\label{th:densitiy_priniciple}
%	Gegeben seien die Banachräume  $ B_1 $ und $ B_2 $, ein dichter Unterraum $ X $ von $ B_1 $ und ein linearer Operator $ A : X \to B_2 $ mit
%	\begin{align}\label{eq:density_principle}
%	\| A f \|_{B_2} \leq C \| f \|_{B_1}
%	\end{align}
%	für  alle $  f \in X $.
%	Dann gilt \eqref{eq:density_principle} für alle $ f \in B_1 $ und $ A $ wird zu einem beschränkten linearen Operator von $ B_1  $ nach $ B_2 $ fortgesetzt.
%\end{genericthm}

%\begin{genericthm}{Open-Mapping-Theorem}\label{th:open_mapping}
%	Gegeben seien die Banach -oder Frecheträume  $ B_1 $ und $ B_2 $ und der surjektive beschränkte Operator $ A : B_1  \to B_2 $.
%	Dann ist $ A $ eine offene Abbildung.
%	Insbesondere können wir diesen auf Operatoren von $ \S(\R^d)  $ nach $ \S(\R^d) $ anwenden.
%\end{genericthm}
%
%\begin{genericthm}{Closed Range Theorem}\label{th:closed_range}
%	Gegeben seien die Banachräume $ B_1 $ und $ B_2 $
%	und ein beschränkter linearer Operator $ A : B_1 \to B_2 $.
%	Der Operator $ A $ hat ein abgeschlossenes Bild genau dann, wenn der 
%	adjungierte Operator $ A^\prime : B_2^\prime \to B_1^\prime $
%	ein abgeschlossenes Bild besitzt.
%	Insbesondere ist $ A $ genau dann surjektiv, wenn $ A^\prime $ injektiv ist und ein abgeschlossenes Bild hat.
%\end{genericthm}
\newpage
\nocite{*}
\bibliographystyle{abbrv}
\bibliography{bib_file}
%\allowdisplaybreaks
\newpage \blankpage
\thispagestyle{empty}
\section*{Eigenständigkeitserklärung}

Hiermit bestätige ich, Philipp Beck, dass ich die vorliegende Arbeit eigenständig und ausschließlich unter Verwendung der angegebenen Quellen erstellt habe. 
Ich versichere, dass ich alle wörtlich oder sinngemäß aus anderen Werken übernommenen Aussagen als solche gekennzeichnet habe, und dass die eingereichte Arbeit weder vollständig noch in wesentlichen Teilen Gegenstand eines anderen Prüfungsverfahren gewesen ist.\\
\newline\\
\newline
\begin{tabular}{lp{2em}l}
	\hspace{3cm}   && \hspace{3cm} \\\cline{1-1}\cline{2-3}
	Datum     && Unterschrift
\end{tabular}
\vspace{1cm}
\section*{Danksagung}

Ein großer Dank geht an meinen Betreuer Professor Dr. Guido Schneider. Er hatte stets ein offenes Ohr für auftretende Probleme und gute Ratschläge, um diese zu behandeln.
%%Er hatte die ganze Zeit ein offenes Ohr und gute Ratschläge für meine Probleme.
%Ein herzliches Dankeschön geht an Alexander Marroquin, ohne den mir manche Unstimmigkeiten nicht aufgefallen wären.
Meinen Kommilitionen Alisa Baransegeta, Aleksandar Arsenijevic, Arthur Günthner, Daniel Winkle, Paul Schwahn und Volkan Bodur möchte ich ebenso für das Korrekturlesen und anregenden Diskussionen danken.
%Ein weiteres Dankeschön geht an Moritz Gösling für die Korrektur der Rechtschreibung.
Außerdem danke ich meiner Familie für die Unterstützung meines gesamten Studiums.
%%Natürlich sollte ich meine Familie nicht vergessen, die mir in der ganzen Zeit den Rücken freigehalten hat.%Ein großer Dank geht auch an meine Familie für den Halt den sie mir gibt. 

\end{document}
